% Les acides carboxyliques et leurs dérivés — Chimie Tle TI — Cours signé OnBuch+
% Programme officiel MINESEC, Terminale TI. Compiler avec : tectonic L02-acides-carboxyliques-et-derives.tex
\newcommand{\DOCMATIERE}{Chimie}
\newcommand{\DOCNIVEAU}{Tle TI}
\newcommand{\DOCMODULE}{Module 1 · Chimie organique}
\newcommand{\DOCLECON}{02}
\newcommand{\DOCTITRE}{Les acides carboxyliques et leurs dérivés}
% OnBuch+ — préambule « cours signés » ENRICHI (compilé avec tectonic / XeLaTeX)
% Système visuel « Pop » d'OnBuch+ : fond crème (jamais blanc), bordures encre
% épaisses, ombres « dures » décalées, titres Archivo Black en capitales.
% Version enrichie pour les cours de chimie : chimie (mhchem, chemfig),
% graphes (pgfplots), boîtes pédagogiques supplémentaires (définition,
% propriété, méthode, attention, le savais-tu, exercice résolu, exercice,
% TP, bilan), page de garde refaite et signature OnBuch+ ancrée en bas.
%
% Macros attendues dans main.tex AVANT \input{preamble} :
%   \DOCMATIERE  \DOCNIVEAU  \DOCMODULE  \DOCLECON (numéro)  \DOCTITRE
\documentclass[11pt]{article}

\usepackage{fontspec}
\usepackage[french]{babel}
\usepackage[a4paper,margin=2cm,top=3.4cm,bottom=2.8cm,headheight=1.4cm,footskip=1.3cm]{geometry}
\usepackage{amsmath,amssymb,amsfonts}
\usepackage[table]{xcolor} % [table] : fournit aussi \rowcolors (alternance de lignes)
\usepackage{pagecolor}
\usepackage{tikz}
\usetikzlibrary{arrows.meta,positioning,calc,shapes.geometric,shapes.misc,shapes.arrows,decorations.pathmorphing,decorations.markings,patterns,shadows,mindmap,trees,fit,backgrounds,matrix,intersections}
\usepackage{pgfplots}
\pgfplotsset{compat=1.18}
\usepackage[version=4]{mhchem}
\usepackage{chemfig}
\usepackage{enumitem}
\usepackage{fancyhdr}
% [most] charge toutes les bibliothèques tcolorbox (listings, minted, raster,
% documentation...) et gonfle inutilement la mémoire TeX sur un document long
% (~300 boîtes) : on ne charge que ce qui est réellement utilisé.
\usepackage[skins,breakable]{tcolorbox}
\usepackage{graphicx}
\usepackage{colortbl}
\usepackage{booktabs}
\usepackage{tabularx}
\usepackage{array}
\usepackage{multicol}
\usepackage{siunitx}
% parse-numbers=false : accepte aussi \SI{2,5 \times 10^{-3}}{...}
\sisetup{locale=FR,output-decimal-marker={,},group-minimum-digits=4,parse-numbers=false}
\DeclareSIUnit\ohm{\ensuremath{\symup{\Omega}}} % U+2126 absent de la police
\usepackage{qrcode}
% \degree : commande fréquemment utilisée par les modèles pour un angle ou une
% température en dehors de \SI{}{} (non fournie par les paquets chargés).
\providecommand{\degree}{^{\circ}}
\usepackage{lastpage}
\usepackage{hyperref}
\hypersetup{hidelinks}

% ============================================================
% Palette « Pop » OnBuch+ — mêmes hex que la classe P (plus_ui.dart)
% ============================================================
\definecolor{poporange}{HTML}{F59321}
\definecolor{popdark}{HTML}{E07A0C}
\definecolor{popcream}{HTML}{FFF6E8}
\definecolor{popink}{HTML}{1C1714}
\definecolor{poppurple}{HTML}{7A5AE0}
\definecolor{popgreen}{HTML}{2E9E6A}
\definecolor{poppink}{HTML}{E0457B}
\definecolor{popblue}{HTML}{2F7DE1}
\definecolor{popgold}{HTML}{C98A12}
\definecolor{popmuted}{HTML}{8A7F76}
\definecolor{popline}{HTML}{EADFD2}
\definecolor{popgray}{HTML}{8A7F76}
\colorlet{popgrey}{popgray}
% Alias tolérants (le modèle nomme parfois les couleurs différemment)
\colorlet{popwhite}{white}
\colorlet{popblack}{popink}
\colorlet{popred}{poppink}
\colorlet{popyellow}{popgold}
% Teintes claires pour fonds de tableaux / zones
\colorlet{popblueL}{popblue!10!white}
\colorlet{poporangeL}{poporange!14!white}
\colorlet{popgreenL}{popgreen!12!white}
\colorlet{poppurpleL}{poppurple!12!white}
\colorlet{poppinkL}{poppink!12!white}
\colorlet{popgoldL}{popgold!14!white}

\pagecolor{popcream}

% ============================================================
% Typographie
% ============================================================
\defaultfontfeatures{Ligatures=TeX}
\setmainfont{PlusJakartaSans-Medium.ttf}[Path=../../../fonts/,BoldFont=PlusJakartaSans-Bold.ttf]
% Maths en sans-serif (Fira Math) avec chiffres et lettres droites de Plus
% Jakarta, pour que les formules se fondent dans le texte.
\usepackage[math-style=ISO,bold-style=ISO]{unicode-math}
\setmathfont{FiraMath-Regular.otf}[Path=../../../fonts/]
\setmathfont{PlusJakartaSans-Medium.ttf}[Path=../../../fonts/,range={up/num,up/{Latin,latin},\mathup}]
\usepackage{icomma} % virgule décimale sans espace en mode math
% Caractères absents de la police de texte (sinon carré vide) : rendus par
% leur équivalent mathématique, en texte comme en formule.
\usepackage{newunicodechar}
\newunicodechar{α}{\ensuremath{\alpha}}
\newunicodechar{Α}{\ensuremath{\alpha}}
\newunicodechar{β}{\ensuremath{\beta}}
\newunicodechar{δ}{\ensuremath{\delta}}
\newunicodechar{✓}{\popcheck{popgreen}}
\newfontfamily\popdisplay{ArchivoBlack-Regular.ttf}[Path=../../../fonts/]
\newfontfamily\poplogo{SpaceGrotesk-Bold.ttf}[Path=../../../fonts/]
\newfontfamily\popsemi{PlusJakartaSans-SemiBold.ttf}[Path=../../../fonts/]
\newfontfamily\popxbold{PlusJakartaSans-ExtraBold.ttf}[Path=../../../fonts/]
\newfontfamily\popxboldls{PlusJakartaSans-ExtraBold.ttf}[Path=../../../fonts/,LetterSpace=3.0]

\setlist[enumerate]{leftmargin=1.7em,itemsep=5pt,topsep=4pt}
\setlist[itemize]{leftmargin=1.7em,itemsep=4pt,topsep=4pt,label={\color{poporange}\textbullet}}
\setlength{\parindent}{0pt}
\setlength{\parskip}{5pt}
\renewcommand{\arraystretch}{1.25}

% chemfig : liaisons un peu plus épaisses, encre Pop
\setchemfig{atom sep=2em,bond style={line width=0.9pt,popink},cram width=3pt}

% pgfplots : style Pop par défaut
\pgfplotsset{
  every axis/.append style={axis line style={popink,line width=1pt},tick style={popink},
    grid style={popline,line width=0.5pt},label style={font=\small},tick label style={font=\footnotesize}},
  popcurve/.style={poporange,line width=1.8pt,no markers,smooth},
}

% ============================================================
% Pill / squiggle
% ============================================================
\newcommand{\poppill}[3]{%
  \tikz[baseline=-0.55ex]\node[draw=popink,line width=1.2pt,rounded corners=2.6pt,%
    fill=#2,inner xsep=6.5pt,inner ysep=3pt]{%
    \popxboldls\fontsize{7}{7}\selectfont\color{#3}\MakeUppercase{#1}};%
}
\newcommand{\popsquiggle}[1]{%
  \tikz[baseline=-0.4ex]{
    \draw[#1,line width=2.6pt,line cap=round]
      plot [smooth] coordinates {(0,0.09) (0.34,0.01) (0.68,0.21) (1.02,0.01) (1.36,0.10)};
  }%
}
% Mot-clé surligné (marqueur orange sous le texte)
\newcommand{\cle}[1]{{\popxbold\color{popdark}#1}}

% ============================================================
% En-tête / pied
% ============================================================
\pagestyle{fancy}
\fancyhf{}
\renewcommand{\headrulewidth}{0pt}
\renewcommand{\footrulewidth}{0pt}
\lhead{\raisebox{-0.15ex}{\poplogo\fontsize{13}{13}\selectfont\color{popink}OnBuch\color{poporange}+}}
\rhead{\poppill{\DOCMATIERE\ \textbullet\ \DOCNIVEAU\ \textbullet\ Leçon \DOCLECON}{white}{popink}}
\lfoot{\popxbold\fontsize{7.6}{7.6}\selectfont\color{popmuted}\MakeUppercase{Cours signé OnBuch+}\ \textbullet\ {\popsemi\DOCTITRE}}
\rfoot{\tikz[baseline=-0.6ex]\node[rounded corners=8.5pt,fill=popink,minimum height=17pt,minimum width=17pt,inner xsep=5pt,inner sep=0]{\popxbold\fontsize{8}{8}\selectfont\color{popcream}\thepage\,/\,\pageref*{LastPage}};}

% ============================================================
% Titres
% ============================================================
\newcommand{\chaptitre}[2]{%
  \begin{tcolorbox}[enhanced,colback=poporange,colframe=popink,boxrule=2.6pt,arc=11pt,
      drop shadow=popink,left=15pt,right=15pt,top=12pt,bottom=11pt,before skip=3pt,after skip=18pt]
    \poppill{#2}{popink}{popcream}\par\vspace{9pt}
    {\raggedright\hyphenpenalty=10000\exhyphenpenalty=10000\popdisplay\fontsize{22}{24}\selectfont\color{popcream}\MakeUppercase{#1}\par}
  \end{tcolorbox}
}
% Section numérotée : pastille orange avec le numéro + titre Archivo
\newcounter{popsec}
\newcommand{\coursec}[1]{%
  \stepcounter{popsec}\setcounter{popsub}{0}%
  \par\vspace{16pt}\noindent
  \begin{minipage}{\linewidth}
  \tikz[baseline=-0.7ex]\node[draw=popink,line width=1.6pt,fill=poporange,rounded corners=4pt,minimum size=22pt,inner sep=0]{\popdisplay\fontsize{12}{12}\selectfont\color{popcream}\thepopsec};%
  \hspace{8pt}{\popdisplay\fontsize{15}{17}\selectfont\color{popink}\MakeUppercase{#1}}\par
  \vspace{4pt}\noindent{\color{poporange}\rule{36pt}{3.6pt}}
  \end{minipage}\par\nopagebreak\vspace{8pt}
}
\newcounter{popsub}
\newcommand{\courssub}[1]{%
  \stepcounter{popsub}%
  \par\vspace{9pt}\noindent{\popxbold\fontsize{11.5}{13}\selectfont\color{popink}{\color{poporange}\thepopsec.\thepopsub}\hspace{6pt}#1}\par\nopagebreak\vspace{4pt}
}

% ============================================================
% Boîtes pédagogiques — même recette : fond blanc, bordure épaisse colorée,
% ombre dure encre, badge plein. Titre optionnel [#1].
% ============================================================
\tcbset{popbase/.style={breakable,enhanced,colback=white,boxrule=2.3pt,arc=9pt,
  shadow={3pt}{-3pt}{0pt}{popink},left=14pt,right=14pt,top=11pt,bottom=11pt,before skip=11pt,after skip=15pt,
  fontupper=\popsemi\fontsize{10.6}{14.5}\selectfont\color{popink},
  fontlower=\popsemi\fontsize{10.6}{14.5}\selectfont\color{popink}}}
% Coche dessinée (le glyphe ✓ n'existe pas dans les polices du document)
\newcommand{\popcheck}[1]{\tikz[baseline=-0.5ex]\draw[#1,line width=1.6pt,line cap=round,line join=round](-0.13,0.0)--(-0.04,-0.09)--(0.13,0.1);}
\newcommand{\popboxhead}[3]{\poppill{#1}{#2}{white}\ifx\relax#3\relax\else\hspace{6pt}{\popxbold\fontsize{10.6}{12}\selectfont\color{popink}#3}\fi\par\vspace{7pt}}

\newtcolorbox{aretenir}[1][]{popbase,colframe=popgreen,before upper={\popboxhead{À retenir}{popgreen}{#1}}}
\newtcolorbox{exemplebox}[1][]{popbase,colframe=poporange,before upper={\popboxhead{Exemple}{poporange}{#1}}}
\newtcolorbox{definition}[1][]{popbase,colframe=popblue,before upper={\popboxhead{Définition}{popblue}{#1}}}
\newtcolorbox{propriete}[1][]{popbase,colframe=poppurple,before upper={\popboxhead{Propriété}{poppurple}{#1}}}
\newtcolorbox{methode}[1][]{popbase,colframe=popgold,before upper={\popboxhead{Méthode}{popgold}{#1}}}
\newtcolorbox{attention}[1][]{popbase,colframe=poppink,before upper={\popboxhead{Attention}{poppink}{#1}}}
\newtcolorbox{experience}[1][]{popbase,colframe=popblue,colback=popblueL,before upper={\popboxhead{Expérience}{popblue}{#1}}}
% « Le savais-tu ? » : carte violette pleine (contexte, culture, Cameroun)
\newtcolorbox{savaistu}[1][]{popbase,colback=poppurple,colframe=popink,
  fontupper=\popsemi\fontsize{10.4}{14.2}\selectfont\color{white},
  before upper={\poppill{Le savais-tu\,?}{white}{poppurple}\ifx\relax#1\relax\else\hspace{6pt}{\popxbold\color{white}#1}\fi\par\vspace{7pt}}}
% Exercice résolu : énoncé en haut, \tcblower puis la solution sur fond crème
\newcounter{popexo}\newcounter{popexr}
\newtcolorbox{exoresolu}[1][]{popbase,colframe=poporange,colbacklower=popcream,
  segmentation style={popink,line width=1.2pt,dashed},
  before upper={\stepcounter{popexr}\popboxhead{Exercice résolu \thepopexr}{poporange}{#1}},
  before lower={\poppill{Solution}{popink}{popcream}\par\vspace{6pt}}}
% Exercice (non résolu en place) — la correction est dans la partie Corrigés
\newtcolorbox{exercice}[1][]{popbase,colframe=popink,
  before upper={\stepcounter{popexo}\popboxhead{Exercice \thepopexo}{popink}{#1}}}
% Corrigé
\newtcolorbox{corrige}[1][]{popbase,colframe=popgreen,colback=popgreenL,
  before upper={\popboxhead{Corrigé}{popgreen}{#1}}}
% Objectifs / prérequis / bilan
\newtcolorbox{objectifs}{popbase,colframe=popink,colback=poporangeL,
  before upper={\popboxhead{Objectifs de la leçon}{popink}{}}}
\newtcolorbox{prerequis}{popbase,colframe=popink,colback=popblueL,
  before upper={\popboxhead{Prérequis}{popblue}{}}}
\newtcolorbox{bilan}{popbase,colframe=popink,colback=white,boxrule=2.8pt,
  before upper={\popboxhead{Fiche bilan}{poporange}{L'essentiel à connaître pour l'examen}}}
% Figure encadrée avec légende
\newtcolorbox{popfigure}[1][]{enhanced,breakable=false,colback=white,colframe=popline,boxrule=1.4pt,arc=8pt,
  left=10pt,right=10pt,top=10pt,bottom=8pt,before skip=10pt,after skip=12pt,halign=center,
  lower separated=false,halign lower=center,
  fontlower=\popsemi\fontsize{9}{11.5}\selectfont\color{popmuted},
  #1}
\newcommand{\legende}[1]{\par\vspace{6pt}{\popsemi\fontsize{9}{11.5}\selectfont\color{popmuted}#1}}

% ============================================================
% Signature OnBuch+
% ============================================================
\newcommand{\poplogoinline}[1]{{\poplogo\fontsize{#1}{#1}\selectfont\color{popink}OnBuch\color{poporange}+}}
\newcommand{\signatureonbuch}{%
  \begin{tcolorbox}[enhanced,colback=popink,colframe=popink,boxrule=2.2pt,arc=10pt,
      drop shadow=poporange,left=14pt,right=14pt,top=10pt,bottom=10pt,before skip=0pt,after skip=0pt]
    \tikz[baseline=-0.7ex]\node[circle,fill=popgreen,draw=popcream,line width=1.3pt,minimum size=22pt,inner sep=0]{\popcheck{white}};%
    \hspace{9pt}\begin{minipage}[c]{0.62\linewidth}
      {\popxbold\fontsize{12}{13}\selectfont\color{popcream}Signé\ \textbullet\ {\poplogo OnBuch\color{poporange}+}}\par\vspace{2pt}
      {\raggedright\popsemi\fontsize{8}{10.5}\selectfont\color{popline}\MakeUppercase{Équipe pédagogique OnBuch+}\\\MakeUppercase{Conforme au programme officiel MINESEC}\par}
    \end{minipage}\hfill
    {\poplogo\fontsize{22}{22}\selectfont\color{popcream}OnBuch\color{poporange}+}
  \end{tcolorbox}%
}

% ============================================================
% Page de garde
% #1 titre · #2 matière · #3 niveau · #4 module · #5 n° leçon · #6 durée ·
% #7 description / objectifs (texte ou itemize)
% ============================================================
\newcommand{\pagegarde}[7]{%
  \thispagestyle{empty}
  \newgeometry{a4paper,margin=1.7cm,top=1.6cm,bottom=1.4cm}%
  \begin{tikzpicture}[remember picture,overlay]
    \fill[poporange!18] ([xshift=-3.2cm,yshift=-3.6cm]current page.north east) circle (4.6cm);
    \fill[poppurple!14] ([xshift=2.4cm,yshift=7.6cm]current page.south west) circle (3.8cm);
  \end{tikzpicture}%
  {\poplogo\fontsize{30}{30}\selectfont\color{popink}OnBuch\color{poporange}+}\hfill
  \raisebox{0.6ex}{\poppill{Cours signé \textbullet\ Premium}{popink}{popcream}}\par
  \vspace{1pt}\popsquiggle{poppurple}\par
  \vspace{8pt}
  \begin{tcolorbox}[enhanced,colback=poporange,colframe=popink,boxrule=2.8pt,arc=13pt,
      drop shadow=popink,left=18pt,right=18pt,top=12pt,bottom=13pt,after skip=10pt]
    \poppill{#2 \textbullet\ #3}{popink}{popcream}\hspace{5pt}\poppill{#4}{white}{popink}\par\vspace{8pt}
    {\popdisplay\fontsize{14}{14}\selectfont\color{popink}LEÇON #5}\par\vspace{4pt}
    {\raggedright\hyphenpenalty=10000\exhyphenpenalty=10000\popdisplay\fontsize{25}{27}\selectfont\color{popcream}\MakeUppercase{#1}\par}
    \vspace{8pt}
    {\popxbold\fontsize{9.5}{9.5}\selectfont\color{popink}\MakeUppercase{Durée conseillée : #6}}
  \end{tcolorbox}
  \begin{tcolorbox}[enhanced,colback=white,colframe=popink,boxrule=2.2pt,arc=10pt,
      drop shadow=popink,left=16pt,right=16pt,top=10pt,bottom=10pt,after skip=8pt]
    \poppill{Dans cette leçon}{poppurple}{white}\par\vspace{5pt}
    \popsemi\fontsize{9.3}{12.6}\selectfont\color{popink}#7
  \end{tcolorbox}
  \noindent\begin{minipage}[t]{0.487\linewidth}
    \begin{tcolorbox}[enhanced,colback=popgreen,colframe=popink,boxrule=2.2pt,arc=10pt,
        drop shadow=popink,left=11pt,right=11pt,top=10pt,bottom=10pt,height=3.1cm,valign=center]
      \tcbox[colback=white,colframe=popink,boxrule=1.3pt,arc=5pt,boxsep=0pt,nobeforeafter,
        left=3pt,right=3pt,top=3pt,bottom=3pt,valign=center]{\qrcode[height=1.9cm]{https://play.google.com/store/apps/details?id=com.luvvix.onbuch&hl=fr}}%
      \hspace{8pt}\begin{minipage}[c]{0.5\linewidth}
        {\popdisplay\fontsize{11}{12.5}\selectfont\color{white}\MakeUppercase{Télécharge OnBuch}}\par\vspace{4pt}
        {\raggedright\popsemi\fontsize{8.4}{11}\selectfont\color{white}Gratuit \textbullet\ Google Play\\Tuteur IA, annales, concours, résultats\par}
      \end{minipage}
    \end{tcolorbox}
  \end{minipage}\hfill
  \begin{minipage}[t]{0.487\linewidth}
    \begin{tcolorbox}[enhanced,colback=poppurple,colframe=popink,boxrule=2.2pt,arc=10pt,
        drop shadow=popink,left=11pt,right=11pt,top=10pt,bottom=10pt,height=3.1cm,valign=center]
      \tikz[baseline=-0.7ex]\node[circle,fill=white,draw=popink,line width=1.3pt,minimum size=1.6cm,inner sep=0]{\popdisplay\fontsize{24}{24}\selectfont\color{poppurple}+};%
      \hspace{8pt}\begin{minipage}[c]{0.6\linewidth}
        {\popdisplay\fontsize{11}{12.5}\selectfont\color{white}\MakeUppercase{Passe à OnBuch+}}\par\vspace{4pt}
        {\raggedright\popsemi\fontsize{8.4}{11}\selectfont\color{white}Tous les cours signés, Tuteur IA illimité, fiches PDF — onglet OnBuch+ de l'app\par}
      \end{minipage}
    \end{tcolorbox}
  \end{minipage}
  \vspace{8pt}
  \popcontenu
  \vfill
  \signatureonbuch
  \restoregeometry
  \newpage
}

% Rangée « Ce cours contient » (page de garde)
\newcommand{\poptile}[3]{% #1 couleur #2 gros texte #3 légende
  \begin{tcolorbox}[enhanced,colback=#1,colframe=popink,boxrule=2pt,arc=9pt,shadow={2.5pt}{-2.5pt}{0pt}{popink},
      left=6pt,right=6pt,top=9pt,bottom=9pt,halign=center,valign=center,height=1.75cm,width=\linewidth,nobeforeafter]
    {\popdisplay\fontsize{12.5}{13}\selectfont\color{white}\MakeUppercase{#2}}\par\vspace{3pt}
    {\centering\popsemi\fontsize{7.8}{9.5}\selectfont\color{white}#3\par}
  \end{tcolorbox}}
\newcommand{\popcontenu}{%
  \par\noindent{\popxboldls\fontsize{7.5}{7.5}\selectfont\color{popmuted}CE COURS SIGNÉ CONTIENT}\par\vspace{7pt}
  \noindent\begin{minipage}[t]{0.235\linewidth}\poptile{popblue}{Cours}{complet, illustré, pas à pas}\end{minipage}\hfill
  \begin{minipage}[t]{0.235\linewidth}\poptile{poporange}{Méthodes}{et exercices résolus}\end{minipage}\hfill
  \begin{minipage}[t]{0.235\linewidth}\poptile{poppink}{Exercices}{type examen + corrigés}\end{minipage}\hfill
  \begin{minipage}[t]{0.235\linewidth}\poptile{popgreen}{Fiche bilan}{l'essentiel à retenir}\end{minipage}\par}

% Sommaire
\newtcolorbox{sommaire}{enhanced,colback=white,colframe=popink,boxrule=2.2pt,arc=10pt,
  drop shadow=popink,left=18pt,right=18pt,top=13pt,bottom=13pt,before skip=6pt,after skip=20pt,
  before upper={\poppill{Sommaire}{poppurple}{white}\par\vspace{9pt}\popsemi\fontsize{10.6}{14.5}\selectfont\color{popink}}}

% Signature de fin de cours — poussée tout en bas de la dernière page
\newcommand{\findecours}{%
  \par\vfill\nopagebreak
  \begin{center}\popsquiggle{poporange}\end{center}\vspace{4pt}
  \signatureonbuch
}
% Glyphes absents de Fira Math : redéfinis à partir de glyphes disponibles
\AtBeginDocument{%
  \def\setminus{\mathbin{\backslash}}\let\smallsetminus\setminus
  \def\vdots{\mathord{\vbox{\baselineskip3.2pt\lineskiplimit0pt\kern4pt\hbox{.}\hbox{.}\hbox{.}}}}%
  \def\ddots{\mathinner{\mkern1mu\raise6.5pt\hbox{.}\mkern2mu\raise3.6pt\hbox{.}\mkern2mu\raise0.7pt\hbox{.}\mkern1mu}}%
  \def\triangle{\mathord{\scalebox{0.9}{$\symup{\Delta}$}}}%
  \def\nabla{\mathord{\raisebox{\depth}{\scalebox{1}[-1]{$\symup{\Delta}$}}}}%
  \def\checkmark{\popcheck{popdark}}%
  \def\star{\mathbin{\ast}}\def\bigstar{\mathord{\ast}}%
  \def\diamond{\mathbin{\scalebox{0.6}{$\Diamond$}}}%
  \def\bigcup{\mathop{\mathchoice{\raisebox{-0.3em}{\scalebox{1.7}{$\cup$}}}{\scalebox{1.25}{$\cup$}}{\cup}{\cup}}\displaylimits}%
  \def\bigcap{\mathop{\mathchoice{\raisebox{-0.3em}{\scalebox{1.7}{$\cap$}}}{\scalebox{1.25}{$\cap$}}{\cap}{\cap}}\displaylimits}%
  \def\lesseqqgtr{\mathrel{\lessgtr}}\def\bigoplus{\mathop{\oplus}\displaylimits}%
}
\providecommand{\ssi}{si et seulement si} % abréviation fréquente

%% FIGFIX-BEGIN v5 — correctifs globaux des figures (docs/onbuch-generation/tools/figfix)
\usepackage{adjustbox}\usepackage{etoolbox}\tcbuselibrary{raster}
% toute figure TikZ/pgfplots plus large que la ligne est réduite à la largeur disponible
\BeforeBeginEnvironment{tikzpicture}{\begin{adjustbox}{max width=\linewidth}}
\AfterEndEnvironment{tikzpicture}{\end{adjustbox}}
% légendes pgfplots lisibles par-dessus les courbes
\pgfplotsset{every axis/.append style={legend style={fill=white,fill opacity=0.92,text opacity=1,draw=black!15,inner sep=3pt}}}
% étiquettes d'arêtes (syntaxe edge["..."]) avec fond blanc
\tikzset{every edge quotes/.append style={fill=white,inner sep=1.2pt}}
%% FIGFIX-END
\begin{document}

\pagegarde{Les acides carboxyliques et leurs dérivés}{Chimie}{Tle TI}{Module 1 · Chimie organique}{02}{5 h (cours + TP)}{Cette leçon présente les acides carboxyliques, leur structure, leur nomenclature et leurs propriétés physiques et chimiques. Elle montre ensuite comment leurs dérivés (chlorures d'acyle, anhydrides, esters, amides) permettent de fabriquer des produits industriels majeurs : le nylon 6,6, le tergal et le savon.
\vspace{4pt}
\begin{multicols}{2}\small
\begin{itemize}
\item À la fin de cette leçon, je sais écrire la formule générale d'un acide carboxylique et identifier le groupe carboxyle
\item À la fin de cette leçon, je sais écrire les formules semi-développées des monoacides saturés, insaturés et des polyacides
\item À la fin de cette leçon, je sais nommer un acide carboxylique et ses dérivés (chlorure d'acyle, anhydride, ester, amide)
\item À la fin de cette leçon, je sais expliquer les propriétés physiques des acides carboxyliques par les liaisons hydrogène
\item À la fin de cette leçon, je sais écrire l'équation-bilan de l'action de l'eau sur un acide carboxylique
\item À la fin de cette leçon, je sais écrire les équations-bilans de formation des chlorures d'acyle, des anhydrides et des esters
\end{itemize}\end{multicols}}

\begin{sommaire}
\begin{multicols}{2}
\begin{enumerate}
\item Structure et nomenclature des acides carboxyliques
\item Propriétés physiques et liaisons hydrogène
\item Propriété acide des acides carboxyliques
\item Les dérivés des acides carboxyliques : chlorures d'acyle et anhydrides
\item L'estérification et la synthèse des esters
\item Polycondensation : nylon 6,6 et tergal
\item La saponification : fabrication du savon
\item Travaux pratiques
\item Méthodes et exercices résolus
\item Exercices
\item Corrigés des exercices
\item Fiche bilan
\end{enumerate}
\end{multicols}
\end{sommaire}

\begin{objectifs}
\begin{itemize}
\item À la fin de cette leçon, je sais écrire la formule générale d'un acide carboxylique et identifier le groupe carboxyle
\item À la fin de cette leçon, je sais écrire les formules semi-développées des monoacides saturés, insaturés et des polyacides
\item À la fin de cette leçon, je sais nommer un acide carboxylique et ses dérivés (chlorure d'acyle, anhydride, ester, amide)
\item À la fin de cette leçon, je sais expliquer les propriétés physiques des acides carboxyliques par les liaisons hydrogène
\item À la fin de cette leçon, je sais écrire l'équation-bilan de l'action de l'eau sur un acide carboxylique
\item À la fin de cette leçon, je sais écrire les équations-bilans de formation des chlorures d'acyle, des anhydrides et des esters
\item À la fin de cette leçon, je sais écrire les équations-bilans des polycondensations conduisant au nylon 6,6 et au tergal
\item À la fin de cette leçon, je sais décrire la fabrication d'un savon par saponification et écrire l'équation-bilan correspondante
\end{itemize}
\end{objectifs}

\begin{prerequis}
\begin{itemize}
\item Nomenclature des alcanes, alcènes et alcools (Première)
\item Groupes fonctionnels : hydroxyle, carbonyle ; notion de famille fonctionnelle
\item Réaction d'oxydation ménagée des alcools primaires en acides carboxyliques (leçon précédente)
\item Notion d'équilibre chimique et de réaction limitée (constat d'Avogadro)
\item Écriture et ajustement des équations-bilans en chimie organique
\end{itemize}
\end{prerequis}

\newpage

\coursec{Structure et nomenclature des acides carboxyliques}

\textbf{Situation d'accroche.} Il est 7 h du matin au marché Mokolo, à Yaoundé. La maman de Sandrine négocie, comme chaque semaine, quelques pains de savon artisanal fabriqués à partir d'huile de palme. Pendant ce temps, son frère cadet enfile son maillot de football en tergal et rêve à voix haute d'une corde en nylon, bien solide, pour son filet de pêche à Kribi. Savon, tergal, nylon : trois objets du quotidien camerounais qui semblent n'avoir rien en commun. Et pourtant ! Ils sont tous fabriqués à partir d'une même famille de molécules : les \cle{acides carboxyliques} et leurs dérivés. Comment une seule famille chimique peut-elle donner naissance à des produits aussi différents ? Pour le comprendre, commençons par le commencement : la structure de ces molécules et la manière de les nommer.

\courssub{Le groupe carboxyle : définition et structure}

Ouvre une bouteille de vinaigre : cette odeur piquante, c'est l'acide éthanoïque. Goûte un citron : l'acidité vient de l'acide citrique. Toutes ces molécules partagent un même « badge » structural, un petit groupe d'atomes qui leur confère leur personnalité chimique.

\begin{definition}[Acide carboxylique et groupe carboxyle]
Un \cle{acide carboxylique} est un composé organique oxygéné possédant le \cle{groupe carboxyle} $-\ce{COOH}$. Sa formule générale s'écrit $R-\ce{COOH}$, où R est un atome d'hydrogène ou un groupe alkyle ($C_nH_{2n+1}-$). Le groupe carboxyle est formé d'un \cle{groupe carbonyle} $\ce{C=O}$ et d'un \cle{groupe hydroxyle} $-\ce{O-H}$ portés par le \emph{même} atome de carbone, appelé carbone fonctionnel.
\end{definition}

Regardons ce groupe de près. Le carbone fonctionnel est lié à trois atomes : le groupe R par une liaison simple, l'oxygène du carbonyle par une liaison double, et l'oxygène de l'hydroxyle par une liaison simple. Ce carbone est \cle{trigonal} : il est \cle{hybridé sp}$^2$, ce qui impose une géométrie \emph{plane} autour de lui, avec des angles de liaison voisins de \SI{120}{\degree}. Cette planéité n'est pas un détail : elle rapproche les deux atomes d'oxygène et permet une délocalisation électronique qui explique, comme nous le verrons plus loin dans la leçon, l'acidité de la fonction.

\begin{popfigure}
\begin{center}
\begin{tikzpicture}[scale=1.0]
\node at (0,2.6) {\chemfig{R-C(=[2]O)(-[6]OH)}};
% annotations géométriques
\draw[popblue,thick] (0.05,2.05) arc[start angle=200,end angle=340,radius=0.55];
\node[popblue] at (0.05,1.15) {\small $\approx \SI{120}{\degree}$};
\draw[poporange,thick] (1.15,2.35) arc[start angle=60,end angle=160,radius=0.5];
\node[poporange] at (1.25,3.15) {\small $\approx \SI{120}{\degree}$};
\draw[popgreen!70!black,<->,thick] (-1.6,0.4) -- (2.6,0.4);
\node[popgreen!70!black] at (0.5,0.05) {\small les 4 atomes R, C, O, O sont dans un même plan};
\node[popdark] at (0.5,-0.7) {\small carbone fonctionnel trigonal, hybridation sp$^2$};
\end{tikzpicture}
\legende{Structure du groupe carboxyle : le carbone fonctionnel, hybridé sp$^2$, est trigonal ; le groupe est plan avec des angles voisins de \SI{120}{\degree}.}
\end{center}
\end{popfigure}

\begin{attention}[Ne confonds pas carboxyle et carbonyle !]
Erreur très fréquente au baccalauréat : confondre le groupe \cle{carboxyle} $-\ce{COOH}$ avec le groupe \cle{carbonyle} $\ce{C=O}$ des aldéhydes et des cétones. Dans une cétone, le carbone du $\ce{C=O}$ est lié à \emph{deux} atomes de carbone. Dans un acide carboxylique, le carbone du $\ce{C=O}$ porte \emph{aussi} un groupe $-\ce{OH}$. C'est ce $-\ce{OH}$ voisin qui change tout : il rend l'hydrogène mobile et donne à la molécule son caractère acide. Astuce : « carboxyle = carbonyle + hydroxyle, sur le même carbone ».
\end{attention}

\courssub{Les grandes familles d'acides carboxyliques}

\textbf{Les monoacides saturés.} Si la chaîne carbonée R ne contient que des liaisons simples C$-$C, on parle de monoacide \cle{saturé}. En comptant les atomes, on obtient la formule brute générale :

\begin{aretenir}[Formule générale des monoacides saturés]
Un monoacide carboxylique saturé a pour formule brute générale $C_nH_{2n}O_2$ avec $n \geq 1$, soit encore $C_{n-1}H_{2n-1}-\ce{COOH}$.
\end{aretenir}

Vérifions sur deux exemples : l'acide méthanoïque \ce{HCOOH} ($n=1$ : $C_1H_2O_2$, correct) et l'acide éthanoïque \ce{CH3COOH} ($n=2$ : $C_2H_4O_2$, correct).

\textbf{Les monoacides insaturés.} La chaîne R peut contenir une ou plusieurs doubles liaisons C=C. Deux exemples capitaux pour la suite du programme :
\begin{itemize}
\item l'\cle{acide acrylique} \ce{CH2=CH-COOH} (acide prop-2-énoïque), monomère de certains plastiques et peintures ;
\item l'\cle{acide oléique} \ce{C17H33-COOH}, un acide gras à 18 atomes de carbone présent en abondance dans l'huile de palme : c'est lui qui, par saponification, donnera le savon de la maman de Sandrine (section 7).
\end{itemize}

\textbf{Les polyacides.} Certaines molécules portent \emph{plusieurs} groupes carboxyles. Les plus importants sont les \cle{acides dioïques}, de formule générale $\ce{HOOC}-(\ce{CH2})_n-\ce{COOH}$ :
\begin{itemize}
\item l'\cle{acide oxalique} \ce{HOOC-COOH} ($n=0$), présent dans les feuilles d'oseille ;
\item l'\cle{acide adipique} \ce{HOOC-(CH2)4-COOH} ($n=4$), l'un des deux monomères du nylon 6,6 ;
\item l'\cle{acide téréphtalique} (acide benzène-1,4-dioïque), monomère du tergal.
\end{itemize}
Retiens bien ces deux derniers noms : ce sont eux qui permettent la polycondensation, car chaque molécule possède \emph{deux} fonctions réactives capables de s'enchaîner en longues fibres.

\begin{popfigure}
\begin{center}
\begin{tikzpicture}
\node[align=center] at (0,0) {\chemfig{H-C(=[2]O)-[6]OH}\\[2pt]\small acide méthanoïque};
\node[align=center] at (4.6,0) {\chemfig{CH_3-C(=[2]O)-[6]OH}\\[2pt]\small acide éthanoïque};
\node[align=center] at (9.6,0) {\chemfig{CH_2=CH-C(=[2]O)-[6]OH}\\[2pt]\small acide acrylique};
\node[align=center] at (2.3,-3.4) {\chemfig{HO-C(=[2]O)-CH_2-CH_2-CH_2-CH_2-C(=[2]O)-[6]OH}\\[2pt]\small acide adipique (hexanedioïque)};
\node[align=center] at (9.2,-3.4) {\chemfig{HO-C(=[2]O)-[6]*6(-=-(-C(=[2]O)-[6]OH)=-=)}\\[2pt]\small acide téréphtalique};
\end{tikzpicture}
\legende{Formules semi-développées de quelques acides carboxyliques de référence : monoacides saturés, insaturé et diacides.}
\end{center}
\end{popfigure}

\courssub{Nomenclature des acides carboxyliques}

Nommer correctement une molécule, c'est la première compétence exigée au baccalauréat TI. La bonne nouvelle : la règle est simple et systématique.

\begin{methode}[Nommer un acide carboxylique]
\begin{enumerate}
\item \textbf{Identifier} la chaîne carbonée la plus longue contenant le carbone du groupe carboxyle.
\item \textbf{Numéroter} cette chaîne en commençant obligatoirement par le carbone du carboxyle, qui porte le numéro 1.
\item \textbf{Nommer} la chaîne comme l'alcane correspondant, en remplaçant le « e » final par le suffixe « \cle{-oïque} », et en faisant précéder le tout du mot « \cle{acide} ».
\item \textbf{Indiquer} les ramifications éventuelles avec leur position (le carbone 1 étant celui du $-\ce{COOH}$).
\end{enumerate}
Pour un diacide, on utilise le suffixe « -dioïque » : hexane $\rightarrow$ acide hexane\emph{dioïque}.
\end{methode}

\begin{exemplebox}[Appliquons la méthode]
Nommons la molécule \ce{CH3-CH(CH3)-CH2-COOH}.
\begin{enumerate}
\item Chaîne la plus longue contenant $-\ce{COOH}$ : 4 carbones $\rightarrow$ butane.
\item Numérotation : le carbone du $-\ce{COOH}$ est le C1 ; le groupe méthyle est donc sur le C3.
\item Nom : \textbf{acide 3-méthylbutanoïque}.
\end{enumerate}
Piège classique : numéroter depuis l'autre bout donnerait « 2-méthyl » — faux, car le carbone du carboxyle est \emph{toujours} le C1.
\end{exemplebox}

De nombreux acides carboxyliques sont connus depuis des siècles et portent des noms usuels qu'il faut connaître, car ils reviennent souvent dans les énoncés :

\begin{center}
\begin{tabularx}{\linewidth}{|l|l|X|}
\hline
\rowcolor{popblueL} \textbf{Formule} & \textbf{Nom systématique} & \textbf{Nom usuel et origine} \\
\hline
\ce{HCOOH} & acide méthanoïque & acide formique (fourmis, \emph{formica} en latin) \\
\hline
\ce{CH3COOH} & acide éthanoïque & acide acétique (vinaigre, \emph{acetum}) \\
\hline
\ce{CH3-CH2-COOH} & acide propanoïque & acide propionique \\
\hline
\ce{CH3-(CH2)2-COOH} & acide butanoïque & acide butyrique (beurre rance, \emph{butyrum}) \\
\hline
\ce{HOOC-COOH} & acide éthanedioïque & acide oxalique (oseille) \\
\hline
\ce{HOOC-(CH2)4-COOH} & acide hexanedioïque & acide adipique (monomère du nylon 6,6) \\
\hline
\end{tabularx}
\end{center}

\begin{exoresolu}[Nommer et écrire : entraînement type bac]
\textbf{Énoncé.} (a) Donne le nom systématique de l'acide de formule semi-développée \ce{CH3-CH2-CH(CH3)-COOH}. (b) Écris la formule semi-développée de l'acide 2-méthylpropanoïque, puis vérifie que sa formule brute suit bien $C_nH_{2n}O_2$.
\tcblower
\textbf{Solution.} (a) La chaîne la plus longue contenant le $-\ce{COOH}$ compte 4 carbones (butane). Le carbone du carboxyle est le C1 ; le méthyle est sur le C2. Nom : \textbf{acide 2-méthylbutanoïque}.

(b) Propanoïque $\rightarrow$ chaîne de 3 carbones ; méthyle sur le C2 : \chemfig{CH_3-CH(-[2]CH_3)-C(=[2]O)-[6]OH}, soit \ce{(CH3)2CH-COOH}. Formule brute : \ce{C4H8O2}. Vérification : pour $n=4$, $C_nH_{2n}O_2$ donne $C_4H_8O_2$ : la formule est bien celle d'un monoacide saturé. Remarque : c'est un isomère de l'acide butanoïque \ce{CH3-(CH2)2-COOH}.
\end{exoresolu}

\begin{savaistu}[L'acide des fourmis]
L'acide méthanoïque \ce{HCOOH} doit son nom usuel, acide \emph{formique}, aux fourmis (\emph{formica} en latin) : c'est lui qu'elles injectent lors de leurs piqûres et qui provoque la douleur et la rougeur. Quiconque a travaillé dans un champ de manioc au Cameroun connaît bien ces désagréables brûlures ! Le remède de grand-mère — frotter la piqûre avec de la cendre ou du savon — est chimiquement fondé : cendres et savon sont basiques et neutralisent l'acide. La chimie était déjà là, bien avant le laboratoire.
\end{savaistu}

\begin{aretenir}[L'essentiel de la section]
\begin{itemize}
\item Un acide carboxylique possède le groupe \cle{carboxyle} $-\ce{COOH}$ ; formule générale $R-\ce{COOH}$.
\item Le carbone fonctionnel est trigonal (sp$^2$) : le groupe carboxyle est \emph{plan}.
\item Monoacides saturés : $C_nH_{2n}O_2$ ; il existe aussi des acides insaturés (acrylique, oléique) et des diacides (adipique, téréphtalique).
\item Nomenclature : chaîne la plus longue contenant $-\ce{COOH}$, carbone du carboxyle numéroté 1, suffixe « -oïque » précédé de « acide ».
\end{itemize}
\end{aretenir}

\begin{exercice}[À toi de jouer]
(a) Nomme les acides suivants : \ce{CH3-(CH2)4-COOH} ; \ce{HOOC-(CH2)2-COOH} ; \ce{CH2=CH-COOH}. (b) Écris la formule semi-développée de l'acide 2,2-diméthylpropanoïque et donne sa formule brute. (c) L'acide oléique a pour formule \ce{C17H33-COOH} : explique, sans calcul, pourquoi il ne suit pas la formule $C_nH_{2n}O_2$.
\end{exercice}

\begin{corrige}[Exercice]
(a) Acide hexanoïque ; acide butanedioïque (acide succinique) ; acide prop-2-énoïque (acrylique). (b) \ce{(CH3)3C-COOH}, formule brute \ce{C5H10O2}, conforme à $C_nH_{2n}O_2$ avec $n=5$. (c) \ce{C17H33} contient deux atomes d'hydrogène de moins que le groupe saturé \ce{C17H35} : la chaîne possède une double liaison C=C ; l'acide oléique est donc \emph{insaturé}, d'où l'écart à la formule des monoacides saturés.
\end{corrige}

Maintenant que tu sais reconnaître et nommer ces molécules, une question naturelle se pose : pourquoi l'acide éthanoïque bout-il à \SI{118}{\celsius} alors que des molécules de masse voisine sont gazeuses ? Réponse dans la section suivante, grâce à une liaison discrète mais puissante : la liaison hydrogène.

\coursec{Propriétés physiques et liaisons hydrogène}

Dans la section précédente, nous avons appris à reconnaître le \cle{groupe carboxyle} \chemfig{-C(=[1]O)-[7]OH} et à nommer les acides carboxyliques. Mais pourquoi l'acide éthanoïque du vinaigre est-il liquide à température ordinaire alors que le butane, de masse molaire presque identique, est un gaz ? Pourquoi l'acide méthanoïque se mélange-t-il à l'eau en toutes proportions alors que l'huile de palme refuse de s'y dissoudre ? Toutes ces propriétés physiques s'expliquent par un phénomène unique : la formation de \cle{liaisons hydrogène} entre les molécules. C'est l'objet de cette section.

\courssub{Polarité du groupe carboxyle}

\textbf{Intuition.} Dans une molécule, les électrons des liaisons ne sont pas toujours partagés équitablement. L'atome d'oxygène est très \emph{électronégatif} : il attire les électrons vers lui, comme un aimant attire la limaille de fer.

\begin{propriete}[Polarisation des liaisons du groupe carboxyle]
Dans le groupe carboxyle \ce{-COOH}, l'oxygène est plus électronégatif que le carbone et que l'hydrogène. Il en résulte deux liaisons fortement \cle{polarisées} :
\begin{itemize}
\item la liaison \ce{C=O} : le carbone porte une charge partielle $\delta^+$ et l'oxygène une charge partielle $\delta^-$ ;
\item la liaison \ce{O-H} : l'oxygène porte une charge partielle $\delta^-$ et l'hydrogène une charge partielle $\delta^+$.
\end{itemize}
Le groupe carboxyle possède donc un \cle{moment dipolaire} important : c'est un groupe \emph{polaire}.
\end{propriete}

\begin{popfigure}
\begin{center}
\chemfig{C(-[2]R)(=[@{a}]@{o1}[:30]^{\delta^-}O)(-[:-30]@{o2}O^{\delta^-}-[@{b}]@{h1}H^{\delta^+})}
\chemmove{
\draw[popink,thick,->](a)..controls +(80:5mm) and +(100:5mm)..(o1);
\draw[popink,thick,->](b)..controls +(80:5mm) and +(100:5mm)..(o2);
}
\end{center}
\legende{Polarisation des liaisons du groupe carboxyle : les flèches indiquent le déplacement des doublets d'électrons vers l'oxygène, plus électronégatif.}
\end{popfigure}

Cette polarité a une conséquence capitale : l'hydrogène du groupe \ce{-OH}, chargé positivement ($\delta^+$), va être attiré par les doublets non liants des atomes d'oxygène ($\delta^-$) des molécules voisines. C'est la naissance de la liaison hydrogène.

\courssub{Les liaisons hydrogène et les dimères d'acides carboxyliques}

\begin{definition}[Liaison hydrogène]
Une \cle{liaison hydrogène} est une interaction attractive entre un atome d'hydrogène lié à un atome très électronégatif (O, N) et un \emph{doublet non liant} d'un autre atome électronégatif (O, N) appartenant à une molécule voisine. On la représente par des pointillés : \ce{O-H...O}. Son énergie, de l'ordre de \SI{20}{kJ.mol^{-1}}, est environ \textbf{20 fois plus faible} qu'une liaison covalente, mais bien plus forte que les interactions de Van der Waals.
\end{definition}

Dans un acide carboxylique, chaque molécule possède \emph{à la fois} un donneur de liaison hydrogène (le H du groupe \ce{-OH}) et un accepteur (les doublets de l'oxygène du \ce{C=O}). Deux molécules peuvent donc s'associer par \textbf{deux} liaisons hydrogène simultanées pour former un \cle{dimère} cyclique très stable.

\begin{popfigure}
\begin{center}
{\setchemfig{bond offset=2pt}
\chemfig{
  % Dimère de l'acide acétique (deux molécules liées par deux liaisons hydrogène)
  ?(-[::60,1.5]CH_3-C(=[::60]@{O1}O)(-[::-60]O-@{H1}H))
  (-[::-60,1.5]CH_3-C(=[::-60]@{O2}O)(-[::60]O-@{H2}H))
}
\chemmove{
  \draw[popblue,dashed,thick](H1)--(O2);
  \draw[popblue,dashed,thick](H2)--(O1);
}}
\end{center}
\legende{Dimère de l'acide éthanoïque : deux molécules s'associent par deux liaisons hydrogène (pointillés bleus), formant un cycle à huit centres.}
\end{popfigure}

\begin{propriete}[Existence des dimères]
Les acides carboxyliques s'associent par paires en \cle{dimères} grâce à deux liaisons hydrogène. Cette association est si forte qu'elle persiste \textbf{même en phase vapeur} : à l'état gazeux, l'acide éthanoïque se comporte comme une molécule de masse molaire double ($2 \times 60 = \SI{120}{g.mol^{-1}}$ au lieu de \SI{60}{g.mol^{-1}}).
\end{propriete}

\begin{attention}[La liaison hydrogène n'est PAS une liaison covalente]
Erreur très fréquente au baccalauréat : écrire que la liaison hydrogène est une liaison covalente entre l'hydrogène et l'oxygène. C'est faux ! Une liaison covalente met en commun un doublet d'électrons et vaut environ \SI{400}{kJ.mol^{-1}}. La liaison hydrogène est une simple \textbf{interaction électrostatique intermoléculaire}, environ 20 fois plus faible. Elle se forme et se rompt facilement, mais elle suffit à « coller » les molécules entre elles.
\end{attention}

\courssub{Conséquence : des températures de fusion et d'ébullition élevées}

Pour faire bouillir un liquide, il faut \emph{séparer} les molécules les unes des autres et les envoyer en phase gazeuse. Si les molécules sont retenues par des liaisons hydrogène — et pire, associées en dimères qu'il faut briser — il faut fournir beaucoup plus d'énergie, donc chauffer plus fort. Les températures de fusion et d'ébullition des acides carboxyliques sont donc \textbf{anormalement élevées}.

\begin{methode}[Expliquer une température d'ébullition élevée]
Face à la question « pourquoi cet acide carboxylique bout-il à température si élevée ? », rédige ta réponse en trois étapes :
\begin{enumerate}
\item Citer la polarité du groupe carboxyle (liaisons \ce{C=O} et \ce{O-H} polarisées).
\item Affirmer l'existence de \cle{liaisons hydrogène intermoléculaires}, qui associent les molécules en \emph{dimères} stables.
\item Conclure : la vaporisation exige de rompre ces liaisons hydrogène (et de dissocier les dimères), ce qui demande une énergie importante, d'où une température d'ébullition élevée.
\end{enumerate}
\end{methode}

\begin{exemplebox}[Comparaison à masse molaire voisine]
Comparons trois composés de masses molaires quasi identiques :
\begin{center}
\begin{tabularx}{\linewidth}{|l|X|X|X|}
\hline
\rowcolor{popblueL} \textbf{Composé} & \textbf{Formule} & $M$ (\si{g.mol^{-1}}) & $T_{\text{éb}}$ (\si{\celsius}) \\
\hline
Butane & \ce{C4H10} & 58 & $-0{,}5$ \\
\hline
Propan-1-ol & \ce{C3H7OH} & 60 & 97 \\
\hline
Acide éthanoïque & \ce{CH3COOH} & 60 & 118 \\
\hline
\end{tabularx}
\end{center}
Le butane, molécule \emph{apolaire}, ne forme aucune liaison hydrogène : il bout à \SI{-0,5}{\celsius} et est gazeux à température ordinaire. Le propan-1-ol forme des liaisons hydrogène simples : il bout à \SI{97}{\celsius}. L'acide éthanoïque, associé en \emph{dimères} par deux liaisons hydrogène, bout encore plus haut, à \SI{118}{\celsius}. C'est la preuve expérimentale de l'exceptionnelle cohésion des molécules d'acides carboxyliques.
\end{exemplebox}

\begin{popfigure}
\begin{center}
\begin{tikzpicture}
\begin{axis}[
width=\linewidth, height=7.5cm,
xlabel={Nombre d'atomes de carbone $n$},
ylabel={$T_{\text{éb}}$ (\si{\celsius})},
xmin=0.5, xmax=6.5, ymin=-50, ymax=180,
grid=both, grid style={popline!40},
legend style={at={(0.03,0.97)}, anchor=north west, font=\small, draw=popline},
tick label style={font=\small}, label style={font=\small}
]
\addplot[popmuted, mark=*, thick] coordinates {(1,-161)(2,-89)(3,-42)(4,-0.5)(5,36)(6,69)};
\addlegendentry{Alcanes}
\addplot[popgreen, mark=square*, thick] coordinates {(1,65)(2,78)(3,97)(4,118)(5,138)(6,157)};
\addlegendentry{Alcan-1-ols}
\addplot[popink, mark=triangle*, thick] coordinates {(1,101)(2,118)(3,141)(4,164)};
\addlegendentry{Acides carboxyliques}
\end{axis}
\end{tikzpicture}
\end{center}
\legende{Températures d'ébullition en fonction du nombre d'atomes de carbone. À $n$ égal, on a toujours $T_{\text{éb}}(\text{acide}) > T_{\text{éb}}(\text{alcool}) \gg T_{\text{éb}}(\text{alcane})$, à cause des liaisons hydrogène et des dimères.}
\end{popfigure}

\begin{savaistu}[L'« acide acétique glacial »]
L'acide éthanoïque pur fond à \SI{16,6}{\celsius}. Dans un laboratoire non chauffé, par une matinée fraîche de la saison des pluies à Douala ou à Bafoussam, le flacon d'acide éthanoïque se solidifie littéralement en cristaux ressemblant à de la glace ! C'est pourquoi les chimistes l'ont baptisé \emph{acide acétique glacial}. Cette température de fusion élevée pour une si petite molécule est, encore une fois, la signature des dimères à liaisons hydrogène.
\end{savaistu}

\courssub{Solubilité dans l'eau}

\textbf{Intuition.} L'eau est elle-même un solvant polaire, tissé de liaisons hydrogène. Pour qu'un soluté s'y dissolve, il doit pouvoir « s'insérer » dans ce réseau en formant ses propres liaisons hydrogène avec les molécules d'eau.

\begin{propriete}[Solubilité des acides carboxyliques dans l'eau]
Le groupe carboxyle est \cle{hydrophile} : il forme des liaisons hydrogène avec l'eau, tant par son \ce{O-H} (donneur) que par son \ce{C=O} (accepteur). En revanche, la chaîne carbonée \ce{R-} est \cle{hydrophobe} (apolaire, « graisseuse »). La solubilité résulte de la compétition entre ces deux parties :
\begin{itemize}
\item les acides à chaîne courte, de \ce{C1} à \ce{C4} (méthanoïque, éthanoïque, propanoïque, butanoïque), sont \textbf{miscibles} à l'eau en toutes proportions ;
\item au-delà de \ce{C4}, la partie hydrophobe l'emporte : la solubilité \textbf{décroît rapidement} quand la chaîne s'allonge ;
\item les acides gras à longue chaîne (\ce{C12} à \ce{C18}, comme l'acide palmitique de l'huile de palme) sont pratiquement \textbf{insolubles} dans l'eau.
\end{itemize}
\end{propriete}

\begin{exemplebox}[Quelques valeurs de solubilité à \SI{20}{\celsius}]
\begin{center}
\begin{tabularx}{\linewidth}{|l|X|X|}
\hline
\rowcolor{popgreenL} \textbf{Acide} & \textbf{Nombre de C} & \textbf{Solubilité dans l'eau} \\
\hline
Acide éthanoïque \ce{CH3COOH} & 2 & miscible en toutes proportions \\
\hline
Acide butanoïque \ce{C3H7COOH} & 4 & miscible \\
\hline
Acide pentanoïque \ce{C4H9COOH} & 5 & $\approx \SI{50}{g.L^{-1}}$ \\
\hline
Acide hexanoïque \ce{C5H11COOH} & 6 & $\approx \SI{10}{g.L^{-1}}$ \\
\hline
Acide palmitique \ce{C15H31COOH} & 16 & quasi insoluble \\
\hline
\end{tabularx}
\end{center}
On observe la décroissance brutale de la solubilité dès que la chaîne dépasse quatre atomes de carbone.
\end{exemplebox}

\begin{exoresolu}[Prévoir et justifier une solubilité]
Énoncé. On verse quelques gouttes d'acide propanoïque \ce{C2H5COOH} puis quelques gouttes d'acide octanoïque \ce{C7H15COOH} dans deux tubes contenant chacun \SI{5}{mL} d'eau distillée. On agite. Prévoir l'observation dans chaque tube et justifier.
\tcblower
Solution détaillée.
\begin{enumerate}
\item \textbf{Acide propanoïque} (3 atomes de carbone, chaîne courte) : le mélange est \emph{homogène}, l'acide se dissout totalement. Le groupe carboxyle hydrophile forme des liaisons hydrogène avec l'eau, et la petite chaîne \ce{C2H5-} hydrophobe ne s'y oppose pas suffisamment.
\item \textbf{Acide octanoïque} (8 atomes de carbone) : on observe \emph{deux phases} ; l'acide surnage ou forme des gouttelettes. La longue chaîne \ce{C7H15-}, apolaire et hydrophobe, domine : les liaisons hydrogène du seul groupe \ce{-COOH} ne suffisent plus à entraîner la molécule en solution.
\item \textbf{Conclusion} : la solubilité d'un acide carboxylique dans l'eau décroît quand la longueur de la chaîne carbonée augmente, car la partie hydrophobe l'emporte sur la partie hydrophile.
\end{enumerate}
\end{exoresolu}

\courssub{État physique et odeur à température ordinaire}

À \SI{25}{\celsius}, les acides carboxyliques à chaîne courte (\ce{C1} à \ce{C9}) sont des \textbf{liquides} incolores, visqueux, souvent âcre ; ceux à longue chaîne (à partir d'environ \ce{C10}) sont des \textbf{solides} cireux, comme les acides gras des huiles et des graisses.

Leur odeur est souvent \emph{piquante} ou franchement \emph{désagréable} : l'acide éthanoïque donne l'odeur piquante du vinaigre ; l'acide butanoïque est responsable de l'odeur du \textbf{beurre rance} et de certaines sueurs ; les acides pentanoïque et hexanoïque ont des odeurs de fromage fort ou de chèvre (d'où le nom « acide caproïque », du latin \emph{caper}, le bouc).

\begin{aretenir}[L'essentiel de la section]
\begin{itemize}
\item Le groupe carboxyle est polaire (liaisons \ce{C=O} et \ce{O-H} polarisées).
\item Les acides carboxyliques forment des \cle{liaisons hydrogène} et s'associent en \cle{dimères}, même en phase vapeur.
\item Conséquence : températures de fusion et d'ébullition \emph{élevées} (ex. \SI{118}{\celsius} pour l'acide éthanoïque contre \SI{-0,5}{\celsius} pour le butane de masse voisine).
\item Les acides de \ce{C1} à \ce{C4} sont miscibles à l'eau ; la solubilité décroît quand la chaîne hydrophobe s'allonge.
\item À température ordinaire : liquides (chaîne courte) ou solides (chaîne longue), d'odeur souvent piquante ou désagréable.
\end{itemize}
\end{aretenir}

Ces liaisons hydrogène expliquent aussi le comportement de l'acide carboxylique \emph{face à l'eau} d'un point de vue chimique : l'hydrogène du groupe \ce{-OH}, fortement chargé $\delta^+$, peut être cédé à une molécule d'eau. C'est la \emph{propriété acide}, que nous étudions dans la section suivante.

\coursec{Propriété acide des acides carboxyliques}

Dans la section précédente, nous avons vu que les liaisons hydrogène expliquent les températures d'ébullition élevées et la bonne solubilité des petits acides carboxyliques dans l'eau. Mais lorsqu'on dissout de l'acide éthanoïque dans l'eau, il ne se contente pas de se mélanger : il se passe une véritable \cle{réaction chimique}. C'est elle qui donne au vinaigre son goût piquant, qui fait pétiller le bicarbonate dans une marinade, et qui permet de doser l'acidité du jus de bissap. Cette section est consacrée à cette propriété fondamentale : le caractère \cle{acide} du groupe carboxyle.

\courssub{Action des acides carboxyliques sur l'eau : un équilibre}

\textbf{Intuition d'abord.}\par
Verse un peu de vinaigre dans un verre d'eau et mesure le pH : tu trouveras une valeur voisine de 3. La solution est donc bien acide. Pourtant, si l'acide éthanoïque se transformait \emph{totalement} en ions, le pH serait beaucoup plus bas. Le vinaigre reste « doux » à manipuler : c'est le signe que la réaction avec l'eau est \cle{partielle}. Une grande partie des molécules d'acide reste intacte.

\textbf{La réaction d'ionisation.}\par
Le groupe carboxyle $-\ce{COOH}$ possède un atome d'hydrogène lié à un atome d'oxygène, lui-même voisin d'un groupe carbonyle $\ce{C=O}$ très attracteur d'électrons. La liaison $\ce{O-H}$ est donc fortement polarisée : l'hydrogène peut être cédé sous forme d'ion $\ce{H+}$ à une molécule d'eau, qui joue le rôle de base. On obtient l'équilibre suivant :

\[ \ce{R-COOH + H2O <=> R-COO- + H3O+} \]

L'ion $\ce{R-COO-}$ formé est appelé ion \cle{carboxylate}. La double flèche $\ce{<=>}$ est \emph{essentielle} : elle traduit le fait que la réaction n'est pas totale. À l'équilibre, les quatre espèces coexistent dans la solution.

\begin{definition}[Un acide carboxylique est un acide faible]
Un acide carboxylique $\ce{R-COOH}$ est un \cle{acide faible} : sa réaction avec l'eau est \textbf{partielle} et conduit à un \textbf{équilibre chimique} :
\[ \ce{R-COOH + H2O <=> R-COO- + H3O+} \]
Seule une faible fraction des molécules (quelques pour cent en solution usuelle) cède son proton. La solution contient donc majoritairement des molécules $\ce{R-COOH}$ non dissociées, et une petite quantité d'ions $\ce{R-COO-}$ et $\ce{H3O+}$.
\end{definition}

\begin{methode}[Écrire correctement l'équation d'ionisation d'un acide carboxylique]
\begin{enumerate}
\item \textbf{Identifier le proton acide} : c'est l'hydrogène du groupe $-\ce{COOH}$, et \emph{uniquement} celui-là (jamais un H porté par un carbone).
\item \textbf{Faire perdre $\ce{H+}$} au groupe carboxyle : $-\ce{COOH}$ devient $-\ce{COO-}$. Le reste de la molécule $\ce{R}$ est inchangé.
\item \textbf{Écrire l'accepteur de proton} : l'eau devient $\ce{H3O+}$.
\item \textbf{Vérifier la conservation} : mêmes nombres d'atomes de chaque élément des deux côtés, et même charge totale (ici $0$ à gauche, $(-1)+(+1)=0$ à droite).
\item \textbf{Placer la double flèche} $\ce{<=>}$, car il s'agit d'un équilibre.
\end{enumerate}
\end{methode}

\begin{attention}[Piège classique : la flèche simple]
L'erreur la plus fréquente au baccalauréat est d'écrire :
\[ \ce{R-COOH + H2O -> R-COO- + H3O+} \quad \text{\textbf{FAUX !}} \]
Une flèche simple signifierait une réaction \emph{totale}, ce qui est réservé aux acides forts comme \ce{HCl}. Pour un acide carboxylique, il faut impérativement la \textbf{double flèche} $\ce{<=>}$. Deuxième piège : faire perdre un H de la chaîne carbonée ; seul le H du groupe $-\ce{COOH}$ est acide.
\end{attention}

\courssub{Le couple \ce{R-COOH}/\ce{R-COO-} et la stabilité de l'ion carboxylate}

\textbf{Le couple acide/base.}\par
À tout acide faible est associée sa base conjuguée. Ici, le couple s'écrit \ce{R-COOH}/\ce{R-COO-} : l'acide $\ce{R-COOH}$ cède un proton pour donner sa base conjuguée, l'ion carboxylate $\ce{R-COO-}$. On caractérise la force du couple par sa constante d'acidité $K_a$ ou, plus commodément, par son $pK_a = -\log K_a$. Pour les acides carboxyliques usuels, $pK_a \approx 3$ à $5$ (par exemple $4{,}8$ pour l'acide éthanoïque).

\textbf{Pourquoi l'acide carboxylique est-il beaucoup plus acide qu'un alcool ?}\par
C'est LA question de fond de cette section. Un alcool $\ce{R-OH}$ possède aussi une liaison $\ce{O-H}$ polarisée, pourtant son $pK_a$ vaut environ $16$ : il est des milliards de fois moins acide qu'un acide carboxylique ! La différence vient de la \textbf{stabilité de la base conjuguée formée}.

Quand un alcool perd $\ce{H+}$, il donne un ion alcoolate $\ce{R-O-}$ dont la charge négative est \emph{localisée} sur un seul atome d'oxygène : cet ion est très instable, donc très peu formé. Quand un acide carboxylique perd $\ce{H+}$, l'ion carboxylate $\ce{R-COO-}$ obtenu voit sa charge négative \cle{délocalisée} sur \emph{deux} atomes d'oxygène, grâce au phénomène de \cle{mésomérie} (résonance). Une charge répartie sur deux atomes est beaucoup plus stable qu'une charge concentrée sur un seul : l'ion carboxylate se forme donc facilement, et l'acide est d'autant plus fort.

\begin{popfigure}
\begin{center}
\chemfig{R-C(=[2]O)(-[6]O^{-})} \quad
$\longleftrightarrow$ \quad
\chemfig{R-C(-[2]O^{-})=[6]O} \quad
$\equiv$ \quad
\chemfig{R-C(-[2]O)(-[6]O)} \; avec charge $-1/2$ sur chaque O
\end{center}
\legende{Les deux formes mésomères limites de l'ion carboxylate $\ce{R-COO-}$ (double flèche de résonance $\longleftrightarrow$) et l'hybride de résonance : la charge négative est répartie équitablement sur les deux oxygènes, ce qui stabilise l'ion.}
\end{popfigure}

\begin{aretenir}[L'idée force]
Plus la base conjuguée est \textbf{stable}, plus l'acide est \textbf{fort}. L'ion carboxylate est stabilisé par résonance (charge délocalisée sur deux oxygènes) ; l'ion alcoolate ne l'est pas. C'est pourquoi :
\[ pK_a(\text{acide fort}) < 0 \quad \ll \quad pK_a(\ce{R-COOH}) \approx 3\text{ à }5 \quad \ll \quad pK_a(\ce{R-OH}) \approx 16 \]
Les acides carboxyliques sont donc \textbf{plus acides que les alcools}, mais restent des \textbf{acides faibles} face aux acides forts.
\end{aretenir}

\begin{popfigure}
\begin{center}
\begin{tikzpicture}[scale=1]
% Axe pKa
\draw[-{Stealth[length=3mm]}, thick, popdark] (0,0) -- (13,0) node[below left=2pt and -30pt] {\small $pK_a$};
\foreach \x/\n in {0/0, 1.625/2, 3.25/4, 4.875/6, 6.5/8, 8.125/10, 9.75/12, 11.375/14, 13/16}
  \draw[thick, popdark] (\x,0.12) -- (\x,-0.12) node[below] {\small \n};
% Zone acide fort
\fill[popink!25] (0,0.35) rectangle (1.2,1.15);
\node[align=center, font=\small\bfseries, popink] at (0.6,0.75) {Acide fort\\ \ce{HCl} ($pK_a<0$)};
% Zone acide carboxylique
\fill[poporange!30] (3.6,0.35) rectangle (4.6,1.15);
\node[align=center, font=\small\bfseries, poporange!80!black] at (4.1,0.75) {\ce{R-COOH}\\ $pK_a\approx4{,}8$};
% Zone alcool
\fill[popgreen!30] (12.3,0.35) rectangle (13,1.15);
\node[align=center, font=\small\bfseries, popgreen!60!black] at (12.65,0.75) {\ce{R-OH}\\ $\approx16$};
% Flèche force acide
\draw[-{Stealth[length=3mm]}, very thick, poppurple] (0.3,1.7) -- (12.7,1.7);
\node[font=\small\itshape, poppurple] at (6.5,2.05) {Force de l'acide décroissante (basicité de la base conjuguée croissante)};
\end{tikzpicture}
\end{center}
\legende{Échelle de $pK_a$ : les acides carboxyliques ($pK_a \approx 4{,}8$) se situent entre les acides forts comme \ce{HCl} et les alcools ($pK_a \approx 16$).}
\end{popfigure}

\courssub{Un exemple chiffré pour bien comprendre : le pH du vinaigre dilué}

\begin{exemplebox}[pH d'une solution d'acide éthanoïque à \SI{0,10}{mol.L^{-1}}]
Le couple \ce{CH3COOH}/\ce{CH3COO-} a pour $pK_a = 4{,}8$. Pour une solution de concentration $C = \SI{0,10}{mol.L^{-1}}$, la formule du pH d'un acide faible (approximation valable pour $\tau < 5\,\%$ et $C > \SI{e-4}{mol.L^{-1}}$) donne :
\[ pH = \frac{1}{2}\left(pK_a - \log C\right) = \frac{1}{2}\left(4{,}8 - \log 0{,}10\right) = \frac{1}{2}(4{,}8 + 1) = 2{,}9. \]
Si l'acide était \emph{fort} (dissociation totale), on aurait $pH = -\log C = 1$. La valeur réelle $pH \approx 2{,}9$ confirme que la dissociation est partielle : le taux d'ionisation est
\[ \tau = \frac{[\ce{H3O+}]}{C} = \frac{10^{-2{,}9}}{0{,}10} \approx 1{,}3\,\%. \]
Seulement environ $1$ molécule sur $80$ a cédé son proton : l'acide éthanoïque est bien un acide \textbf{faible}.
\end{exemplebox}

\begin{savaistu}[Le vinaigre, un acide de cuisine]
Le vinaigre consommé au Cameroun — utilisé pour assaisonner les crudités, conserver les piments ou parfumer le poisson braisé — est une solution aqueuse d'acide éthanoïque à environ $6\,\%$ en masse, obtenue par fermentation acétique du vin ou d'autres boissons alcoolisées (les bactéries du genre \emph{Acetobacter} oxydent l'éthanol en acide éthanoïque). Son pH voisin de $2{,}5$ à $3$ suffit à freiner le développement des microbes : c'est un conservateur naturel, connu bien avant la chimie moderne !
\end{savaistu}

\courssub{Réaction avec les bases fortes : une réaction totale}

Si l'acide carboxylique réagit mal avec l'eau (base très faible), il réagit en revanche \textbf{totalement} avec les bases fortes comme l'ion hydroxyde $\ce{OH-}$ (apporté par la soude \ce{NaOH} ou la potasse \ce{KOH}) :

\[ \ce{R-COOH + OH- -> R-COO- + H2O} \]

Cette réaction est \textbf{totale} (flèche simple), \textbf{rapide} et \textbf{exothermique}. Elle s'écrit aussi, avec la soude en solution :
\[ \ce{CH3COOH + (Na+ + OH-) -> (CH3COO- + Na+) + H2O} \]

C'est précisément cette réaction qui constitue le \textbf{principe du dosage} (ou titrage) d'un acide carboxylique par la soude : à l'équivalence, tout l'acide a été transformé en ion carboxylate, et on en déduit la quantité d'acide présente. C'est ainsi qu'on contrôle l'acidité du vinaigre ou des jus de fruits en laboratoire.

\begin{attention}[Ne pas confondre les deux réactions]
\begin{itemize}
\item Avec l'\textbf{eau} : réaction \emph{partielle}, double flèche $\ce{<=>}$, produit $\ce{H3O+}$.
\item Avec une \textbf{base forte} $\ce{OH-}$ : réaction \emph{totale}, flèche simple $\ce{->}$, produit $\ce{H2O}$.
\end{itemize}
Dans les deux cas, l'acide donne l'ion carboxylate $\ce{R-COO-}$ ; ce qui change, c'est la base qui capte le proton et le caractère (limité ou total) de la transformation.
\end{attention}

\courssub{Réaction avec les carbonates et hydrogénocarbonates : un test qualitatif}

Les acides carboxyliques sont assez acides pour décomposer les ions carbonate $\ce{CO3^2-}$ et hydrogénocarbonate $\ce{HCO3-}$, avec un \textbf{dégagement gazeux de dioxyde de carbone} \ce{CO2} (effervescence) :

\[ \ce{R-COOH + HCO3- -> R-COO- + H2O + CO2 ^} \]
\[ \ce{2 R-COOH + CO3^2- -> 2 R-COO- + H2O + CO2 ^} \]

\begin{experience}[Test d'identification d'un acide carboxylique]
\textbf{Protocole} : dans un tube à essais, verser \SI{2}{mL} de la solution à tester, puis ajouter une spatule d'hydrogénocarbonate de sodium \ce{NaHCO3} (bicarbonate de soude). Boucher et recueillir le gaz dégagé dans une solution d'eau de chaux.

\textbf{Observations} : effervescence vive ; le gaz recueilli trouble l'eau de chaux (précipité blanc de \ce{CaCO3}), ce qui prouve qu'il s'agit de \ce{CO2}.

\textbf{Interprétation} : le dégagement de \ce{CO2} est caractéristique d'un acide \emph{assez fort} pour attaquer les carbonates. Les acides carboxyliques donnent ce test ; les alcools et les phénols (plus faibles) ne le donnent pas. Ce test permet donc de distinguer un acide carboxylique d'un alcool.
\end{experience}

\begin{savaistu}[Le bicarbonate en cuisine]
Quand une ménagère ajoute du bicarbonate dans une sauce trop acide ou dans la préparation de certains beignets, l'effervescence observée est exactement cette réaction : l'acide présent (éthanoïque, citrique...) attaque l'ion $\ce{HCO3-}$ et libère du \ce{CO2}, qui fait aussi lever la pâte.
\end{savaistu}

\courssub{Les sels carboxylates : nomenclature}

La réaction d'un acide carboxylique avec une base produit un \cle{sel carboxylate}, composé ionique formé de l'anion $\ce{R-COO-}$ et d'un cation (par exemple \ce{Na+}, \ce{K+}, \ce{Ca^2+}).

\begin{definition}[Nomenclature des sels carboxylates]
Le nom d'un sel carboxylate s'obtient en remplaçant la terminaison « -oïque » de l'acide par « \textbf{-oate de} », suivie du nom du cation :
\begin{center}
\begin{tabularx}{\linewidth}{|l|X|l|}
\hline
\rowcolor{popblueL} \textbf{Acide} & \textbf{Sel} & \textbf{Nom du sel} \\
\hline
\ce{CH3COOH} (éthanoïque) & \ce{CH3COO- + Na+} & éthanoate de sodium \\
\hline
\ce{HCOOH} (méthanoïque) & \ce{HCOO- + K+} & méthanoate de potassium \\
\hline
\ce{C2H5COOH} (propanoïque) & \ce{C2H5COO- + Na+} & propanoate de sodium \\
\hline
\end{tabularx}
\end{center}
\end{definition}

Les sels carboxylates ne sont pas de simples curiosités de nomenclature : ce sont des composés de la vie quotidienne. L'éthanoate de sodium sert de conservateur alimentaire, et surtout, les sels de sodium ou de potassium des acides gras à longue chaîne (comme l'oléate de sodium \ce{C17H33COO- + Na+}) sont... les \textbf{savons} ! Nous verrons dans la dernière section comment la saponification les fabrique à partir des corps gras.

\begin{exoresolu}[Dosage de l'acide éthanoïque d'un vinaigre]
\textbf{Énoncé.} On prélève $V_a = \SI{10,0}{mL}$ d'une solution de vinaigre dilué et on la dose par une solution de soude de concentration $C_b = \SI{0,10}{mol.L^{-1}}$. L'équivalence est atteinte pour $V_{bE} = \SI{12,0}{mL}$.
\begin{enumerate}
\item Écrire l'équation-bilan de la réaction de dosage et préciser ses caractéristiques.
\item Calculer la concentration $C_a$ en acide éthanoïque du vinaigre dilué.
\end{enumerate}
\tcblower
\textbf{Solution détaillée.}
\begin{enumerate}
\item L'acide éthanoïque réagit avec les ions hydroxyde selon :
\[ \ce{CH3COOH + OH- -> CH3COO- + H2O} \]
Cette réaction est \textbf{totale}, \textbf{rapide} et \textbf{unique} : elle convient parfaitement à un dosage. (On notera la flèche simple, contrairement à l'action sur l'eau.)
\item À l'équivalence, les réactifs ont été mélangés dans les proportions stœchiométriques ($1$ pour $1$) :
\[ n_a = n_{bE} \quad \Longrightarrow \quad C_a \times V_a = C_b \times V_{bE} \]
\[ C_a = \frac{C_b \times V_{bE}}{V_a} = \frac{\SI{0,10}{mol.L^{-1}} \times \SI{12,0}{mL}}{\SI{10,0}{mL}} = \SI{0,12}{mol.L^{-1}}. \]
\end{enumerate}
Le vinaigre dilué a donc une concentration en acide éthanoïque de $\boxed{C_a = \SI{0,12}{mol.L^{-1}}}$.
\end{exoresolu}

\begin{aretenir}[L'essentiel de la section]
\begin{itemize}
\item Un acide carboxylique est un \textbf{acide faible} : $\ce{R-COOH + H2O <=> R-COO- + H3O+}$ (double flèche obligatoire).
\item Le couple \ce{R-COOH}/\ce{R-COO-} a un $pK_a$ de $3$ à $5$ ; l'ion carboxylate est \textbf{stabilisé par résonance}, ce qui rend l'acide bien plus fort qu'un alcool ($pK_a \approx 16$).
\item Avec une base forte : $\ce{R-COOH + OH- -> R-COO- + H2O}$, réaction \textbf{totale}, base du \textbf{dosage}.
\item Avec les (hydrogéno)carbonates : dégagement de \ce{CO2}, test qualitatif de la fonction acide carboxylique.
\item Les sels obtenus se nomment « \textbf{-oate de} + cation » : éthanoate de sodium, propanoate de potassium...
\end{itemize}
\end{aretenir}

Maintenant que tu maîtrises le comportement acide du groupe carboxyle, une question naturelle se pose : peut-on \emph{remplacer} le groupe $-\ce{OH}$ du carboxyle par d'autres groupes pour obtenir de nouvelles familles de composés ? C'est exactement l'objet de la prochaine section : les \textbf{dérivés des acides carboxyliques} — chlorures d'acyle et anhydrides — dont la réactivité remarquable ouvrira la voie à la synthèse des esters, des polyesters et des polyamides.

\coursec{Les dérivés des acides carboxyliques : chlorures d'acyle et anhydrides}

Dans la section précédente, nous avons vu que le groupe carboxyle $-\ce{COOH}$ donne aux acides carboxyliques une \emph{propriété acide} : ils cèdent leur proton à l'eau, mais seulement partiellement. Cette molécule a cependant un autre secret : son groupe $-\ce{OH}$ peut être \emph{remplacé} par un atome ou un groupe d'atomes plus réactif. C'est exactement ce que fait un savonnier de Foumban ou un ingénieur de la SONARA lorsqu'il veut fabriquer un ester rapidement : au lieu d'utiliser l'acide carboxylique, lent et timide, il le transforme d'abord en un dérivé beaucoup plus « nerveux » : un \cle{chlorure d'acyle} ou un \cle{anhydride d'acide}. C'est toute la stratégie de cette section : \textbf{rendre l'acide plus réactif pour mieux le faire travailler}.

\courssub{Principe général : substituer le groupe $-\ce{OH}$ du carboxyle}

\begin{definition}[Dérivés des acides carboxyliques]
Un \cle{dérivé d'acide carboxylique} est un composé obtenu en \textbf{substituant le groupe hydroxyle $-\ce{OH}$} du groupe carboxyle par un autre atome ou groupe d'atomes, tout en \textbf{conservant le groupe carbonyle $\ce{C=O}$}. On obtient ainsi la famille générale $\ce{R-CO-Z}$ :
\begin{itemize}
\item si $Z = \ce{Cl}$ : un \cle{chlorure d'acyle}, de formule générale $R-COCl$ ;
\item si $Z = -O-CO-R'$ : un \cle{anhydride d'acide}, de formule générale $(R-CO)_2O$ (anhydride symétrique) ou $R-CO-O-CO-R'$ (anhydride mixte) ;
\item si $Z = -O-R'$ : un \cle{ester} (étudié dans la section suivante) ;
\item si $Z = -NH_2$ : un \cle{amide}.
\end{itemize}
\end{definition}

L'idée intuitive est simple : le groupe $-\ce{OH}$ est un « mauvais partant ». Il s'accroche fortement au carbone du carboxyle, ce qui rend les réactions de l'acide (comme l'estérification) \emph{lentes} et \emph{limitées} par un équilibre. En le remplaçant par un atome de chlore ou par un groupe acyloxy, on crée une liaison $\ce{C-Z}$ beaucoup plus fragile : le dérivé devient un réactif « pressé » de réagir, qui cède facilement son fragment $\ce{R-CO}-$ (appelé \cle{groupe acyle}) à un partenaire comme un alcool ou une amine.

\begin{aretenir}[L'atome commun : le groupe acyle]
Dans tous les dérivés $\ce{R-CO-Z}$, le fragment $\ce{R-CO}-$ est le \cle{groupe acyle}. Il se nomme en remplaçant la terminaison « oïque » de l'acide par « \textbf{yle} » : éthanoïque $\rightarrow$ éthano\textbf{yle} ($\ce{CH3-CO}-$), propanoïque $\rightarrow$ propano\textbf{yle} ($\ce{CH3-CH2-CO}-$), benzoïque $\rightarrow$ benzo\textbf{yle} ($\ce{C6H5-CO}-$).
\end{aretenir}

\courssub{Les chlorures d'acyle}

\textbf{Structure et nomenclature.}\par

Un chlorure d'acyle possède le groupe caractéristique $-\ce{COCl}$, c'est-à-dire un carbonyle dont le carbone porte un atome de chlore : \chemfig{R-C(=[1]O)-[7]Cl}.

Pour le nommer, on écrit « \textbf{chlorure de} » suivi du nom du groupe acyle : on remplace la terminaison « oïque » de l'acide par « \textbf{yle} ».

\begin{center}
\begin{tabularx}{\linewidth}{|l|l|X|}
\hline
\rowcolor{popblueL} \textbf{Acide de départ} & \textbf{Chlorure d'acyle} & \textbf{Nom (IUPAC)} \\
\hline
\ce{H-COOH} (méthanoïque) & \ce{H-COCl} & chlorure de méthanoyle \\
\hline
\ce{CH3-COOH} (éthanoïque) & \ce{CH3-COCl} & chlorure d'éthanoyle \\
\hline
\ce{CH3-CH2-COOH} (propanoïque) & \ce{CH3-CH2-COCl} & chlorure de propanoyle \\
\hline
\ce{C6H5-COOH} (benzoïque) & \ce{C6H5-COCl} & chlorure de benzoyle \\
\hline
\end{tabularx}
\end{center}

\textbf{Préparation : action d'un agent chlorurant sur l'acide.}\par

Pour transformer l'acide en chlorure d'acyle, on utilise un \emph{agent chlorurant} qui arrache le $-\ce{OH}$ et le remplace par un chlore. Trois réactifs classiques sont au programme :

\textbf{(a) Le pentachlorure de phosphore \ce{PCl5}} (solide) :
\[ \ce{R-COOH + PCl5 -> R-COCl + POCl3 + HCl} \]

\textbf{(b) Le trichlorure de phosphore \ce{PCl3}} (liquide) :
\[ \ce{3 R-COOH + PCl3 -> 3 R-COCl + H3PO3} \]

\textbf{(c) Le chlorure de thionyle \ce{SOCl2}} (liquide) — \emph{le réactif le plus pratique} :
\[ \ce{R-COOH + SOCl2 -> R-COCl + SO2 + HCl} \]

Cette dernière réaction est préférée en laboratoire car les deux sous-produits, \ce{SO2} et \ce{HCl}, sont \textbf{gazeux} : ils s'échappent du milieu réactionnel, ce qui déplace l'équilibre et rend la réaction \textbf{totale}, tout en simplifiant la purification du chlorure d'acyle. On chauffe généralement à reflux, en milieu parfaitement anhydre.

\begin{attention}[Erreur fréquente : oublier les produits secondaires]
Lorsqu'on te demande l'équation-bilan de la préparation d'un chlorure d'acyle, n'écris \textbf{jamais} $\ce{R-COOH + SOCl2 -> R-COCl}$ seul : c'est \emph{faux} car la matière n'est pas conservée ! Avec \ce{SOCl2}, il se forme obligatoirement \ce{SO2} et \ce{HCl}. Avec \ce{PCl5}, il se forme \ce{POCl3} et \ce{HCl}. Avec \ce{PCl3}, il se forme \ce{H3PO3} (acide phosphoreux) et il faut 3 molécules d'acide pour 1 \ce{PCl3}. Vérifie toujours la conservation des atomes de chlore, de soufre et de phosphore.
\end{attention}

\begin{attention}[Sécurité : les chlorures d'acyle s'hydrolysent violemment]
Les chlorures d'acyle réagissent \textbf{violemment avec l'eau}, même avec la simple humidité de l'air :
\[ \ce{R-COCl + H2O -> R-COOH + HCl} \]
Cette hydrolyse dégage du chlorure d'hydrogène gazeux, corrosif et suffocant. Toute manipulation de chlorure d'acyle se fait donc \textbf{sous hotte ventilée}, avec de la verrerie \textbf{parfaitement sèche}, loin de toute trace d'eau, et avec gants et lunettes. C'est aussi pourquoi on ne conserve jamais un flacon de chlorure d'acyle ouvert à l'air libre.
\end{attention}

\begin{savaistu}[Pourquoi le chlorure d'éthanoyle « fume » à l'air]
Ouvre (en imagination !) un flacon de chlorure d'éthanoyle \ce{CH3COCl} dans un laboratoire humide de Douala : tu verrais immédiatement des « fumées » blanches s'en échapper. Ce ne sont pas des vapeurs du liquide lui-même, mais le produit de sa réaction avec l'humidité de l'air : le \ce{HCl} gazeux formé se combine à la vapeur d'eau pour donner de fines gouttelettes d'acide chlorhydrique en suspension — un brouillard acide. C'est ce phénomène spectaculaire qui a valu à ces composés leur réputation de produits « fumants ». Les anciens chimistes s'en servaient d'ailleurs comme test qualitatif : un liquide qui fume à l'air et dont les vapeurs font virer au rouge un papier pH humide est très probablement un chlorure d'acyle.
\end{savaistu}

\courssub{Les anhydrides d'acide}

\textbf{Structure et nomenclature.}\par

Un anhydride d'acide résulte de la \textbf{condensation de deux molécules d'acide carboxylique avec élimination d'une molécule d'eau} (d'où son nom : « an-hydride », c'est-à-dire « sans eau »). Les deux groupes acyle sont reliés par un atome d'oxygène : \chemfig{R-C(=[1]O)-[7]O-C(=[1]O)-[7]R'} que l'on écrit plus simplement $\ce{(R-CO)2O}$ lorsque les deux groupes sont identiques (anhydride symétrique).

Pour le nommer, on écrit « \textbf{anhydride} » suivi du nom de l'acide dont il dérive : l'anhydride de l'acide éthanoïque est l'\textbf{anhydride éthanoïque} $\ce{(CH3-CO)2O}$ ; celui de l'acide propanoïque est l'\textbf{anhydride propanoïque}. Si les deux acides sont différents, on cite les deux noms par ordre alphabétique : \ce{CH3-CO-O-CO-CH2-CH3} est l'anhydride éthanoïque propanoïque (anhydride \emph{mixte}).

\textbf{Préparation : déshydratation de deux molécules d'acide.}\par

On chauffe l'acide carboxylique en présence d'un agent déshydratant puissant, le \textbf{décaoxyde de tétraphosphore \ce{P4O10}} (parfois noté \ce{P2O5} par simplification), qui « piège » l'eau formée :
\[ \ce{2 R-COOH ->[\ce{P4O10},\ \text{chauffage}] (R-CO)2O + H2O} \]

\begin{methode}[Équilibrer la formation d'un anhydride]
Pour écrire correctement l'équation de formation d'un anhydride, procède en trois étapes :
\begin{enumerate}
\item Écris \textbf{deux molécules} d'acide carboxylique côte à côte : $\ce{R-COOH + HOOC-R}$.
\item Repère le $-\ce{OH}$ de l'une et le $-\ce{H}$ du $-\ce{COOH}$ de l'autre : ce sont eux qui formeront l'eau $\ce{H2O}$.
\item Relie les deux fragments restants $\ce{R-CO}-$ par l'atome d'oxygène restant : tu obtiens $\ce{(R-CO)2O}$, et il reste exactement \textbf{une molécule d'eau}. Vérifie : 2 acides $\rightarrow$ 1 anhydride + 1 eau.
\end{enumerate}
\end{methode}

\begin{exemplebox}[Formation de l'anhydride éthanoïque]
À partir de l'acide éthanoïque (acide acétique, présent dans le vinaigre) :
\[ \ce{2 CH3-COOH ->[\ce{P4O10},\ \Delta] CH3-CO-O-CO-CH3 + H2O} \]
Vérification de la conservation : à gauche, $\ce{C4H8O4}$ ; à droite, l'anhydride $\ce{C4H6O3}$ plus l'eau $\ce{H2O}$ donnent bien $\ce{C4H8O4}$. L'équation est équilibrée.
\end{exemplebox}

\begin{popfigure}
\begin{center}
\chemfig{CH_3-C(=[1]O)-[7]OH} \hspace{0.4cm} $+$ \hspace{0.2cm} \chemfig{Cl-S(=[1]O)(=[7]O)-Cl} \hspace{0.4cm} $\longrightarrow$ \hspace{0.4cm} \chemfig{CH_3-C(=[1]O)-[7]Cl} \hspace{0.4cm} $+$ \hspace{0.2cm} $\ce{SO2}\uparrow$ \hspace{0.2cm} $+$ \hspace{0.2cm} $\ce{HCl}\uparrow$

\vspace{0.8cm}

\chemfig{CH_3-C(=[1]O)-[7]OH} \hspace{0.4cm} $+$ \hspace{0.2cm} \chemfig{HO-C(=[2]O)-[6]CH_3} \hspace{0.4cm} $\xrightarrow{\ce{P4O10},\ \Delta}$ \hspace{0.4cm} \chemfig{CH_3-C(=[1]O)-[7]O-C(=[1]O)-[7]CH_3} \hspace{0.4cm} $+$ \hspace{0.2cm} $\ce{H2O}$
\end{center}
\legende{En haut : préparation du chlorure d'éthanoyle par action du chlorure de thionyle sur l'acide éthanoïque (sous-produits gazeux \ce{SO2} et \ce{HCl}). En bas : formation de l'anhydride éthanoïque par déshydratation de deux molécules d'acide éthanoïque.}
\end{popfigure}

\courssub{Réactivité comparée : pourquoi ces dérivés sont-ils si utiles ?}

L'hydrolyse est un excellent test de réactivité : plus un dérivé réagit vite avec l'eau, plus il est réactif. Or :
\begin{itemize}
\item l'\textbf{acide carboxylique} ne subit pratiquement pas d'hydrolyse (c'est lui le produit stable !) et ses réactions (estérification) sont lentes et limitées par un équilibre ;
\item l'\textbf{anhydride d'acide} s'hydrolyse lentement à froid, rapidement à chaud : $\ce{(R-CO)2O + H2O -> 2 R-COOH}$ ;
\item le \textbf{chlorure d'acyle} s'hydrolyse \textbf{violemment dès la température ambiante} : $\ce{R-COCl + H2O -> R-COOH + HCl}$.
\end{itemize}

\begin{popfigure}
\begin{center}
\begin{tikzpicture}[scale=0.9]
% Axes
\draw[->, thick, popdark] (0,0) -- (0,4.5) node[above, font=\small\bfseries] {Réactivité croissante};
\draw[->, thick, popdark] (0,0) -- (10.5,0);
% Acide carboxylique
\node[draw=popblue, thick, fill=popblueL, minimum width=2.5cm, minimum height=1.1cm, anchor=south west, font=\small\bfseries] (acide) at (0.4,0.2) {Acide carboxylique};
\node[font=\footnotesize, anchor=north] at (acide.south) {$\ce{R-COOH}$};
\node[font=\footnotesize, anchor=north] at ($(acide.south)+(0,-0.3)$) {Réactions lentes, limitées};
% Anhydride
\node[draw=poporange, thick, fill=poporangeL, minimum width=2.5cm, minimum height=2.3cm, anchor=south west, font=\small\bfseries] (anhyd) at (3.6,0.2) {Anhydride d'acide};
\node[font=\footnotesize, anchor=north] at (anhyd.south) {$\ce{(R-CO)2O}$};
\node[font=\footnotesize, anchor=north] at ($(anhyd.south)+(0,-0.3)$) {Hydrolyse lente à froid};
\node[font=\footnotesize, anchor=north] at ($(anhyd.south)+(0,-0.6)$) {Rapide à chaud};
% Chlorure d'acyle
\node[draw=poppink, thick, fill=poppinkL, minimum width=2.5cm, minimum height=3.5cm, anchor=south west, font=\small\bfseries] (chlor) at (6.8,0.2) {Chlorure d'acyle};
\node[font=\footnotesize, anchor=north] at (chlor.south) {$\ce{R-COCl}$};
\node[font=\footnotesize, anchor=north] at ($(chlor.south)+(0,-0.3)$) {Hydrolyse violente};
\node[font=\footnotesize, anchor=north] at ($(chlor.south)+(0,-0.6)$) {Dès l'ambiante};
\node[font=\footnotesize, anchor=north] at ($(chlor.south)+(0,-0.9)$) {Dégagement de \ce{HCl}};
\node[font=\footnotesize\itshape, anchor=north] at ($(chlor.south)+(0,-1.3)$) {Manipulation sous hotte};
% Flèches
\draw[->, very thick, popgreen] (acide.east) -- (anhyd.west) node[midway, above, font=\tiny] {+ réactif};
\draw[->, very thick, popgreen] (anhyd.east) -- (chlor.west) node[midway, above, font=\tiny] {+ réactif};
\end{tikzpicture}
\end{center}
\legende{Hiérarchie de réactivité : acide carboxylique $<$ anhydride d'acide $<$ chlorure d'acyle. Plus le dérivé est réactif, plus ses réactions avec l'eau, les alcools ou les amines sont rapides et totales.}
\end{popfigure}

\begin{aretenir}[L'intérêt stratégique des dérivés]
Les chlorures d'acyle et les anhydrides sont des \textbf{intermédiaires de synthèse} précieux : leurs réactions avec les alcools (pour donner des esters) et avec les amines ou l'ammoniac (pour donner des amides) sont \textbf{rapides} et \textbf{totales}, alors que la même réaction à partir de l'acide serait lente et limitée. C'est le principe de la synthèse industrielle de l'aspirine (à partir de l'anhydride éthanoïque) et de nombreux parfums et médicaments.
\end{aretenir}

\begin{exoresolu}[Préparation du chlorure de propanoyle]
\textbf{Énoncé.} Un technicien de laboratoire veut préparer du chlorure de propanoyle \ce{CH3-CH2-COCl} à partir de \SI{14,8}{g} d'acide propanoïque et d'un excès de chlorure de thionyle \ce{SOCl2}. (1) Écris l'équation-bilan de la réaction. (2) Calcule la masse maximale de chlorure de propanoyle obtenue si la réaction est totale. (3) Cite deux précautions expérimentales indispensables. Données : $M(\ce{C3H6O2}) = \SI{74}{g.mol^{-1}}$ ; $M(\ce{C3H5ClO}) = \SI{92,5}{g.mol^{-1}}$.
\tcblower
\textbf{Solution.} (1) L'action de \ce{SOCl2} sur l'acide propanoïque s'écrit, sans oublier les sous-produits gazeux :
\[ \ce{CH3-CH2-COOH + SOCl2 -> CH3-CH2-COCl + SO2 ^ + HCl ^} \]
(2) La réaction étant totale et \ce{SOCl2} en excès, l'acide est le réactif limitant. D'après les coefficients stœchiométriques, $n(\text{chlorure}) = n(\text{acide})$ :
\[ n(\text{acide}) = \frac{m}{M} = \frac{14,8}{74} = \SI{0,20}{mol} \]
\[ m(\text{chlorure}) = n \times M = 0,20 \times 92,5 = \SI{18,5}{g} \]
On peut donc obtenir au maximum \SI{18,5}{g} de chlorure de propanoyle.
(3) Précautions : travailler \textbf{sous hotte} (dégagement de \ce{SO2} et \ce{HCl}, gaz toxiques et corrosifs) et utiliser de la \textbf{verrerie parfaitement sèche}, car le chlorure d'acyle s'hydrolyse violemment au contact de l'eau, ce qui détruirait le produit attendu.
\end{exoresolu}

\begin{exercice}[Anhydride et nomenclature]
(1) Écris l'équation-bilan de la formation de l'anhydride propanoïque à partir de l'acide propanoïque, en précisant les conditions. (2) Nomme les composés suivants : \ce{CH3-CH2-CH2-COCl} ; \ce{(H-CO)2O}. (3) Classe par ordre de réactivité croissante vis-à-vis de l'eau : anhydride éthanoïque, acide éthanoïque, chlorure d'éthanoyle. Justifie en une phrase.
\end{exercice}

\begin{corrige}[Exercice n]
(1) $\ce{2 CH3-CH2-COOH ->[\ce{P4O10},\ \Delta] (CH3-CH2-CO)2O + H2O}$ : deux molécules d'acide propanoïque perdent une molécule d'eau sous l'action du déshydratant \ce{P4O10} à chaud.
(2) \ce{CH3-CH2-CH2-COCl} : chlorure de butanoyle (acide butanoïque $\rightarrow$ butano\textbf{yle}). \ce{(H-CO)2O} : anhydride méthanoïque.
(3) Acide éthanoïque $<$ anhydride éthanoïque $<$ chlorure d'éthanoyle : le chlorure s'hydrolyse violemment à froid, l'anhydride lentement à froid, et l'acide est le produit stable de ces hydrolyses.
\end{corrige}

En résumé, tu sais désormais transformer un acide carboxylique « paresseux » en un dérivé « pressé » : le chlorure d'acyle (avec \ce{PCl5}, \ce{PCl3} ou surtout \ce{SOCl2}) et l'anhydride d'acide (par déshydratation de deux acides). Ces deux intermédiaires réactifs vont maintenant nous servir d'armes redoutables : dans la section suivante, nous les ferons réagir avec les alcools pour fabriquer des \textbf{esters} rapidement et totalement — parfums, arômes et aspirine ne sont plus très loin !

\coursec{L'estérification et la synthèse des esters}

Dans la section précédente, nous avons préparé le terrain : tu sais désormais qu'un acide carboxylique peut être « activé » en chlorure d'acyle ou en anhydride d'acide. Pourquoi une telle activation ? La réponse tient en un mot : l'\cle{estérification}. La réaction directe entre un acide carboxylique et un alcool est certes possible, mais elle est d'une lenteur décourageante et s'arrête bien avant d'être terminée. Les dérivés d'acides vont nous offrir un raccourci rapide et total vers les esters. Mais commençons par le commencement : la réaction directe elle-même.

\courssub{Qu'est-ce que l'estérification ?}

\textbf{Une situation concrète pour démarrer.}\par
Tu as certainement déjà remarqué l'odeur agréable des bananes mûres sur les étals du marché Mboppi à Douala, ou le parfum des ananas de la plantation de Penja. Ces odeurs ne sont pas dues aux sucres des fruits, mais à des molécules organiques appelées \cle{esters}. Or, les industriels qui fabriquent les arômes alimentaires, les parfums ou les solvants ne récoltent pas ces molécules dans les fruits : ils les synthétisent en laboratoire en faisant réagir un \cle{acide carboxylique} avec un \cle{alcool}. Cette réaction s'appelle l'estérification.

\begin{definition}[L'estérification]
L'\cle{estérification} est la réaction chimique entre un \textbf{acide carboxylique} et un \textbf{alcool}, qui produit un \textbf{ester} et de l'\textbf{eau}. C'est une réaction \textbf{lente}, \textbf{athermique} (elle ne dégage ni n'absorbe pratiquement de chaleur) et \textbf{limitée} par la réaction inverse, l'hydrolyse de l'ester : c'est un équilibre chimique.
\end{definition}

\textbf{L'équation générale.}\par
En notant $\mathrm{R}$ le groupe alkyle (ou hydrogène) porté par l'acide et $\mathrm{R'}$ le groupe alkyle de l'alcool, l'équation générale s'écrit :

\[ \ce{R-COOH + R^{\prime}-OH <=> R-COO-R^{\prime} + H2O} \]

Regarde bien cette équation : le groupe $-\mathrm{OH}$ de l'acide et l'hydrogène $\mathrm{H}$ du groupe hydroxyle de l'alcool s'associent pour former l'eau, tandis que le reste de l'acide ($\mathrm{R-CO}-$) se lie à l'oxygène de l'alcool ($-\mathrm{O-R'}$) pour former l'ester. C'est exactement l'inverse de l'hydrolyse d'un ester.

\begin{popfigure}
\begin{center}
\chemfig{\text{R}-C(=[2]O)(-[6,,,,popblue]@{oh}OH)} \hspace{0,4cm} $+$ \hspace{0,4cm} \chemfig{@{ha}H-[,,,,poporange]O-\text{R'}} \hspace{0,4cm} $\rightleftharpoons$ \hspace{0,4cm} \chemfig{\text{R}-C(=[2]O)(-[6,,,,popblue]O-[,,,,poporange]\text{R'})} \hspace{0,4cm} $+$ \hspace{0,4cm} \chemfig{H_2O}
\chemmove{\draw[popblue,thick,->](oh) .. controls +(south:1.2cm) and +(south west:1.2cm) .. (ha);}
\end{center}
\legende{Équation générale de l'estérification. En bleu : le groupe $-\mathrm{OH}$ fourni par l'acide ; en orange : le groupe $-\mathrm{O-R'}$ fourni par l'alcool. La flèche courbe montre que le $\mathrm{OH}$ de l'acide et le $\mathrm{H}$ de l'alcool s'éliminent ensemble sous forme d'eau.}
\end{popfigure}

\textbf{Un exemple fondamental : l'éthanoate d'éthyle.}\par
Faisons réagir l'acide éthanoïque (l'acide du vinaigre) avec l'éthanol, en présence de quelques gouttes d'acide sulfurique concentré comme catalyseur :

\[ \ce{CH3COOH + C2H5OH <=> CH3COOC2H5 + H2O} \]

On obtient l'\cle{éthanoate d'éthyle}, un ester à l'odeur fruitée utilisé comme solvant dans les colles et les vernis.

\courssub{Les trois caractéristiques de la réaction}

\begin{aretenir}[Trois adjectifs à retenir]
L'estérification directe est :
\begin{itemize}
\item \textbf{lente} : à température ambiante, il faudrait des mois pour atteindre l'équilibre ; même à chaud avec catalyseur, il faut plusieurs heures ;
\item \textbf{athermique} : la variation d'enthalpie est quasi nulle ; la température accélère la réaction mais ne déplace pas l'équilibre ;
\item \textbf{limitée} : c'est un équilibre. L'ester formé réagit avec l'eau (hydrolyse) pour régénérer l'acide et l'alcool. On n'obtient jamais un rendement de \SI{100}{\percent}.
\end{itemize}
\end{aretenir}

\textbf{Pourquoi la réaction est-elle limitée ?}\par
L'eau produite par l'estérification n'est pas un spectateur innocent : elle attaque l'ester formé selon la réaction inverse, l'\cle{hydrolyse} :

\[ \ce{CH3COOC2H5 + H2O <=> CH3COOH + C2H5OH} \]

Au bout d'un certain temps, la vitesse d'estérification (qui consomme les réactifs) devient égale à la vitesse d'hydrolyse (qui les régénère) : les quantités n'évoluent plus. On a atteint un \textbf{équilibre chimique dynamique}.

\textbf{La limite d'équilibre dépend de la classe de l'alcool.}\par
Pour un mélange équimolaire (1 mol d'acide + 1 mol d'alcool), le rendement maximal à l'équilibre vaut environ :
\begin{itemize}
\item \SI{67}{\percent} (soit $\frac{2}{3}$) avec un \textbf{alcool primaire} ;
\item \SI{60}{\percent} avec un \textbf{alcool secondaire} ;
\item seulement \SI{5}{\percent} environ avec un \textbf{alcool tertiaire}.
\end{itemize}

\begin{exemplebox}[Le fameux « deux tiers »]
Partons de $n_0 = \SI{1}{mol}$ d'acide éthanoïque et $n_0 = \SI{1}{mol}$ d'éthanol (alcool primaire). Dressons le tableau d'avancement à l'équilibre, avec $x_{eq}$ l'avancement final :

\begin{center}
\begin{tabularx}{\linewidth}{|l|X|X|X|X|}
\hline
\rowcolor{popblueL}
État & \ce{CH3COOH} & \ce{C2H5OH} & \ce{CH3COOC2H5} & \ce{H2O} \\
\hline
Initial (mol) & 1 & 1 & 0 & 0 \\
\hline
Équilibre (mol) & $1-x_{eq}$ & $1-x_{eq}$ & $x_{eq}$ & $x_{eq}$ \\
\hline
\end{tabularx}
\end{center}

La constante d'équilibre vaut $K = \dfrac{x_{eq}^2}{(1-x_{eq})^2} = 4$, d'où $\dfrac{x_{eq}}{1-x_{eq}} = 2$ et $x_{eq} = \dfrac{2}{3} \approx \SI{0,67}{mol}$. Le rendement maximal est donc de \SI{67}{\percent} : il restera toujours $\frac{1}{3}$ de mole d'acide et d'alcool non transformés.
\end{exemplebox}

\begin{popfigure}
\begin{center}
\begin{tikzpicture}
\begin{axis}[
  width=0.85\linewidth, height=6.2cm,
  xlabel={Temps $t$ (h)}, ylabel={Quantité d'ester formé $n$ (mol)},
  xmin=0, xmax=10, ymin=0, ymax=1.15,
  xtick={0,2,4,6,8,10}, ytick={0, 0.33, 0.67, 1},
  yticklabels={0, {0,33}, {0,67}, {1,00}},
  grid=major, grid style={popline!40},
  axis lines=left, label style={text=popdark},
]
\addplot[popcurve,domain=0:10,samples=200]{0.67*(1-exp(-0.55*x))};
\addplot[poporange, dashed, thick, domain=0:10]{0.67};
\addplot[popmuted, dotted, thick, domain=0:10]{1};
\node[popblue, anchor=south west] at (axis cs:5.2,0.72) {\small limite d'équilibre : $n_{eq} = \SI{0,67}{mol}$};
\node[popmuted, anchor=south west] at (axis cs:5.2,1.02) {\small rendement \SI{100}{\percent} (jamais atteint)};
\node[popblue, anchor=west] at (axis cs:0.3,0.25) {\small montée lente};
\end{axis}
\end{tikzpicture}
\end{center}
\legende{Évolution de la quantité d'ester au cours du temps pour 1 mol d'acide éthanoïque et 1 mol d'éthanol. La courbe monte lentement (réaction lente) puis plafonne à \SI{0,67}{mol} : l'équilibre est atteint, la réaction semble s'arrêter alors qu'elle se poursuit dans les deux sens à la même vitesse.}
\end{popfigure}

\courssub{Comment améliorer le rendement ?}

Puisque l'estérification est un équilibre, on applique la \textbf{loi de Le Chatelier} (déplacement d'équilibre) : pour pousser la réaction vers la formation d'ester, il faut soit ajouter un réactif en excès, soit retirer un produit au fur et à mesure.

\begin{methode}[Augmenter le rendement d'une estérification]
\begin{enumerate}
\item \textbf{Utiliser un excès de l'un des réactifs} (le moins cher, souvent l'alcool) : l'équilibre se déplace vers la droite. Avec un large excès d'alcool, le rendement par rapport à l'acide peut dépasser \SI{90}{\percent}.
\item \textbf{Éliminer l'eau} au fur et à mesure de sa formation (appareil de Dean-Stark, ou acide sulfurique concentré qui absorbe l'eau) : la réaction inverse devient impossible.
\item \textbf{Éliminer l'ester par distillation} s'il est le composé le plus volatil du mélange (cas de l'éthanoate d'éthyle, $T_{eb} = \SI{77}{\celsius}$).
\item \textbf{Ajouter un catalyseur} : quelques gouttes d'acide sulfurique concentré (ions \ce{H+}) ou de chlorure d'hydrogène gazeux. Attention : le catalyseur \emph{accélère} l'atteinte de l'équilibre mais ne change \emph{pas} la limite.
\item \textbf{Chauffer} (chauffage à reflux) : la température augmente la vitesse, sans déplacer l'équilibre (réaction athermique).
\end{enumerate}
\end{methode}

\begin{attention}[Catalyseur et rendement : ne pas confondre]
Erreur classique au baccalauréat : affirmer que « l'acide sulfurique augmente le rendement ». Faux ! Le catalyseur \ce{H+} permet d'atteindre l'équilibre \emph{plus vite}, mais la composition finale à l'équilibre reste la même. Seuls l'excès d'un réactif ou l'élimination d'un produit augmentent réellement le rendement.
\end{attention}

\courssub{Les synthèses rapides et totales des esters}

Voici enfin la récompense de la section précédente : les dérivés d'acides réagissent avec les alcools de façon \textbf{rapide} et \textbf{totale} (réaction non réversible), ce qui donne des rendements excellents.

\textbf{Avec un chlorure d'acyle.}\par
Le chlorure d'acyle est très réactif ; la réaction est vive, exothermique, et dégage du chlorure d'hydrogène gazeux :

\[ \text{R-COCl} + \text{R'-OH} \rightarrow \text{R-COO-R'} + \text{HCl} \]

Exemple : le chlorure d'éthanoyle réagit violemment avec l'éthanol :

\[ \ce{CH3COCl + C2H5OH -> CH3COOC2H5 + HCl} \]

\textbf{Avec un anhydride d'acide.}\par
La réaction est plus douce, mais tout aussi totale :

\[ \mathrm{(R-CO)_2O + R'\text{-}OH \rightarrow R\text{-}COO\text{-}R' + R\text{-}COOH} \]

Exemple industriel majeur : l'anhydride éthanoïque réagit avec l'\textbf{acide salicylique} (dont le groupe hydroxyle est phénolique) pour donner l'\textbf{aspirine} (acide acétylsalicylique). Autre exemple :

\[ \ce{(CH3CO)2O + C2H5OH -> CH3COOC2H5 + CH3COOH} \]

\begin{aretenir}[Comparer les trois voies de synthèse d'un ester]
\begin{center}
\begin{tabularx}{\linewidth}{|l|X|X|}
\hline
\rowcolor{popblueL}
Réactif & Vitesse & Caractère \\
\hline
Acide carboxylique & lente & limitée (équilibre) \\
\hline
Chlorure d'acyle & très rapide (vive) & totale \\
\hline
Anhydride d'acide & rapide & totale \\
\hline
\end{tabularx}
\end{center}
En pratique industrielle, on préfère l'anhydride au chlorure (moins corrosif, pas de dégagement de \ce{HCl}).
\end{aretenir}

\courssub{Nomenclature des esters}

Le nom d'un ester comporte \textbf{deux mots} : le premier décrit la partie venant de l'acide, le second la partie venant de l'alcool.

\begin{methode}[Nommer un ester]
\begin{enumerate}
\item \textbf{Identifier la partie acide} $\mathrm{R-COO}-$ : compter la chaîne carbonée incluant le carbone du groupe ester, nommer l'alcane correspondant, remplacer la terminaison « -ane » par « \textbf{-oate} » (ex. éthanoate, propanoate, butanoate).
\item \textbf{Identifier le groupe alkyle} $-\mathrm{R'}$ venant de l'alcool : méthyle, éthyle, propyle...
\item \textbf{Assembler} : « alkanoate \textbf{de} alkyle », sans oublier le « de ».
\end{enumerate}
Exemple : \ce{CH3COO-C2H5} = éthan\textbf{oate} d'\textbf{éthyle} ; \ce{HCOO-CH3} = méthanoate de méthyle ; \ce{CH3CH2COO-CH3} = propanoate de méthyle.
\end{methode}

\begin{attention}[L'erreur qui coûte des points : inverser les deux parties]
Dans \ce{CH3COO-C2H5}, c'est le groupe \ce{CH3CO}- qui vient de l'\textbf{acide} (éthanoïque) et le groupe \ce{C2H5} qui vient de l'\textbf{alcool} (éthanol). Le nom est donc éthanoate d'éthyle, et jamais « éthanoate de méthyle ». Astuce infaillible : \textbf{le carbone du groupe $\mathrm{C=O}$ appartient toujours à la partie acide}. Autre point subtil prouvé par les expériences de marquage isotopique : \textbf{l'atome d'oxygène de la liaison ester provient de l'alcool}, pas de l'acide. C'est bien le $\mathrm{OH}$ de l'acide qui part dans l'eau.
\end{attention}

\courssub{Les esters autour de nous}

\begin{savaistu}[L'odeur de la banane du marché]
L'odeur caractéristique de la banane mûre, celle qui flotte sur les marchés de Douala ou de Yaoundé, est due à un ester : l'\textbf{éthanoate de 3-méthylbutyle} (encore appelé éthanoate d'isoamyle), de formule \ce{CH3COO-CH2-CH2-CH(CH3)2}. L'industrie alimentaire le synthétise pour aromatiser bonbons, sirops et yaourts « goût banane ». De même, l'éthanoate d'éthyle évoque la pomme, le butanoate d'éthyle l'ananas et le méthanoate d'éthyle le rhum. Les esters sont partout : parfums, arômes, solvants, plastifiants... et même dans ton assiette, car les \textbf{corps gras} comme l'huile de palme rouge sont des \textbf{triesters} du glycérol (triglycérides) — nous les retrouverons dans la section sur la saponification.
\end{savaistu}

\begin{exoresolu}[Rendement d'une estérification et déplacement d'équilibre]
\textbf{Énoncé.} On chauffe à reflux un mélange de $\SI{0,50}{mol}$ d'acide éthanoïque et de $\SI{0,50}{mol}$ d'éthanol en présence de quelques gouttes d'acide sulfurique concentré. À l'équilibre, on dose l'acide restant : il en reste $\SI{0,17}{mol}$.
\begin{enumerate}
\item Écrire l'équation-bilan de la réaction et nommer l'ester formé.
\item Déterminer l'avancement à l'équilibre et le rendement de la réaction.
\item On recommence avec $\SI{0,50}{mol}$ d'acide et $\SI{2,0}{mol}$ d'éthanol. Sans calcul, dire si le rendement (par rapport à l'acide) augmente ou diminue, et justifier.
\end{enumerate}
\tcblower
\textbf{Solution détaillée.}
\begin{enumerate}
\item L'acide éthanoïque réagit avec l'éthanol selon :
\[ \ce{CH3COOH + C2H5OH <=> CH3COOC2H5 + H2O} \]
L'ester obtenu est l'\textbf{éthanoate d'éthyle}. La réaction est lente, athermique et limitée ; l'acide sulfurique joue le rôle de catalyseur.
\item Dressons le tableau d'avancement (en mol) :

\begin{center}
\begin{tabularx}{\linewidth}{|l|X|X|X|X|}
\hline
\rowcolor{popblueL}
État & \ce{CH3COOH} & \ce{C2H5OH} & \ce{CH3COOC2H5} & \ce{H2O} \\
\hline
Initial & 0,50 & 0,50 & 0 & 0 \\
\hline
Équilibre & $0,50 - x_{eq}$ & $0,50 - x_{eq}$ & $x_{eq}$ & $x_{eq}$ \\
\hline
\end{tabularx}
\end{center}

Il reste $\SI{0,17}{mol}$ d'acide, donc $0,50 - x_{eq} = 0,17$, soit $x_{eq} = \SI{0,33}{mol}$. Le rendement est :
\[ \eta = \frac{x_{eq}}{x_{max}} = \frac{0,33}{0,50} = 0,66 \approx \SI{66}{\percent} \]
On retrouve la limite classique de $\frac{2}{3}$ pour un mélange équimolaire avec un alcool primaire.
\item En introduisant un large excès d'éthanol, on déplace l'équilibre vers la formation de l'ester (loi de Le Chatelier). Le rendement \textbf{par rapport à l'acide augmente} et peut dépasser \SI{90}{\percent}. Remarque : le catalyseur, lui, ne changerait rien au rendement final, il ne ferait qu'accélérer l'atteinte de l'équilibre.
\end{enumerate}
\end{exoresolu}

\begin{exercice}[À toi de jouer]
\begin{enumerate}
\item Écris l'équation-bilan de la réaction entre l'acide propanoïque \ce{CH3CH2COOH} et le méthanol \ce{CH3OH}. Nomme l'ester obtenu.
\item Nomme les esters suivants : \ce{HCOO-C2H5} ; \ce{CH3CH2CH2COO-CH3}.
\item Pourquoi la synthèse de l'aspirine utilise-t-elle l'anhydride éthanoïque plutôt que l'acide éthanoïque ?
\end{enumerate}
\end{exercice}

\begin{corrige}[Exercice : À toi de jouer]
\begin{enumerate}
\item $\ce{CH3CH2COOH + CH3OH <=> CH3CH2COOCH3 + H2O}$. L'ester est le \textbf{propanoate de méthyle} (la partie acide compte 3 carbones : propanoate ; l'alcool à 1 carbone donne méthyle).
\item \ce{HCOO-C2H5} : \textbf{méthanoate d'éthyle} ; \ce{CH3CH2CH2COO-CH3} : \textbf{butanoate de méthyle} (4 carbones côté acide, 1 carbone côté alcool).
\item La réaction entre un anhydride et un alcool est \textbf{rapide et totale}, alors que la réaction directe avec l'acide éthanoïque serait lente et limitée à environ \SI{67}{\percent}. L'anhydride garantit donc un bien meilleur rendement industriel.
\end{enumerate}
\end{corrige}

Tu maîtrises maintenant la formation des esters, lente et équilibrée avec un acide, rapide et totale avec ses dérivés. Cette idée d'enchaîner des réactions d'estérification va prendre une ampleur spectaculaire dans la section suivante : en reliant des centaines de molécules par des liaisons ester ou amide, on fabrique les \textbf{polyesters} et les \textbf{polyamides} — le tergal de tes vêtements et le fameux nylon 6,6.

\coursec{Polycondensation : nylon 6,6 et tergal}

Dans la section précédente, nous avons vu qu'un acide carboxylique et un alcool peuvent s'unir par \cle{estérification} pour former un ester, avec élimination d'une molécule d'eau. Que se passerait-il si, au lieu d'utiliser un acide à une seule fonction et un alcool à une seule fonction, on faisait réagir un \emph{diacide} avec un \emph{diol} ? Le premier groupe $-COOH$ réagit avec un $-OH$, puis le second groupe $-COOH$ de la même molécule réagit avec un autre diol, et ainsi de suite : les molécules s'enchaînent comme les wagons d'un train pour former une très longue chaîne, un \cle{polymère}. C'est exactement le principe de la fabrication du tergal de nos pagnes et du nylon de nos filets de pêche. Cette réaction porte un nom : la \cle{polycondensation}.

\courssub{La polycondensation : définition et principe}

\begin{definition}[Polycondensation]
Une \cle{polycondensation} est une réaction de \cle{polymérisation} au cours de laquelle des monomères s'assemblent par réactions successives, avec \textbf{élimination d'une petite molécule} (eau \ce{H2O}, chlorure d'hydrogène \ce{HCl}, méthanol\ldots) à chaque étape de la croissance de la chaîne.
\end{definition}

Pour qu'une polycondensation soit possible, une condition est indispensable :

\begin{propriete}[Condition sur les monomères]
Chaque monomère doit porter \textbf{au moins deux groupes fonctionnels} réactifs (on dit qu'il est \cle{bifonctionnel}) : diacide, diol, diamine, ou molécule mixte portant à la fois $-COOH$ et $-OH$ (ou $-NH_2$). Un monomère monofonctionnel (comme l'éthanol ou l'acide éthanoïque) ne peut former qu'un seul lien : la réaction s'arrête immédiatement et aucune chaîne ne se forme.
\end{propriete}

L'intuition est simple : imagine des perles percées de \emph{deux} trous. On peut les enfiler indéfiniment sur un fil. Une perle percée d'\emph{un seul} trou termine le collier. De même, un monomère bifonctionnel peut se lier à deux voisins, ce qui permet à la chaîne de grandir sans fin.

\begin{attention}[Polycondensation ou polyaddition ?]
Ne confonds pas les deux grands types de polymérisation :
\begin{itemize}
\item la \cle{polycondensation} : assemblage de monomères bifonctionnels \textbf{avec élimination d'une petite molécule} (\ce{H2O}, \ce{HCl}\ldots). Exemples : nylon, tergal ;
\item la \cle{polyaddition} : simple addition de monomères portant une double liaison \ce{C=C}, \textbf{sans aucun sous-produit}. Exemple : $n\,\ce{CH2=CH2} \rightarrow [-\ce{CH2-CH2}-]_n$ (polyéthylène).
\end{itemize}
Au Baccalauréat, la présence d'une petite molécule éliminée dans l'équation-bilan est la signature d'une polycondensation.
\end{attention}

\begin{popfigure}
\begin{center}
\begin{tikzpicture}[scale=1, every node/.style={font=\small}]
% Monomères bifonctionnels
\draw[fill=popblueL, draw=popblue, thick, rounded corners=2pt] (0,0) rectangle (1.6,0.7);
\node at (0.8,0.35) {$A-A$};
\draw[fill=poporangeL, draw=poporange, thick, rounded corners=2pt] (2.2,0) rectangle (3.8,0.7);
\node at (3.0,0.35) {$B-B$};
\draw[fill=popblueL, draw=popblue, thick, rounded corners=2pt] (4.4,0) rectangle (6.0,0.7);
\node at (5.2,0.35) {$A-A$};
\draw[fill=poporangeL, draw=poporange, thick, rounded corners=2pt] (6.6,0) rectangle (8.2,0.7);
\node at (7.4,0.35) {$B-B$};
\node[below=2pt, text=popmuted] at (4.1,0) {monomères bifonctionnels};
% Flèche
\draw[-{Stealth[length=3mm]}, very thick, popdark] (8.6,0.35) -- (9.8,0.35);
\node[above, text=popdark] at (9.2,0.35) {chaleur};
% Chaîne polymère
\draw[fill=popblueL, draw=popblue, thick, rounded corners=2pt] (10.2,0) rectangle (11.4,0.7);
\draw[fill=poporangeL, draw=poporange, thick, rounded corners=2pt] (11.4,0) rectangle (12.6,0.7);
\draw[fill=popblueL, draw=popblue, thick, rounded corners=2pt] (12.6,0) rectangle (13.8,0.7);
\node at (12.0,0.35) {$-A-B-A-$};
\node[right, text=popmuted] at (13.9,0.35) {$\cdots$};
\node[below=2pt, text=popmuted] at (12.0,0) {chaîne polymère};
% Petites molécules
\node[text=popgreen, font=\small\bfseries] at (12.0,-1.1) {$+\;n\,\ce{H2O}$ (petites molécules éliminées)};
\end{tikzpicture}
\legende{Principe de la polycondensation : des monomères bifonctionnels s'enchaînent en éliminant une petite molécule à chaque liaison formée.}
\end{center}
\end{popfigure}

\courssub{Le nylon 6,6 : un polyamide}

\textbf{Les monomères.} Le nylon 6,6 est obtenu par polycondensation de deux monomères :
\begin{itemize}
\item l'\cle{acide hexanedioïque} (ou \emph{acide adipique}), un diacide à 6 atomes de carbone : $\ce{HOOC-(CH2)4-COOH}$ ;
\item l'\cle{hexaméthylènediamine} (ou hexane-1,6-diamine), une diamine à 6 atomes de carbone : $\ce{H2N-(CH2)6-NH2}$.
\end{itemize}

\textbf{La réaction.} Chaque groupe carboxyle $-COOH$ du diacide réagit avec un groupe amine $-NH_2$ de la diamine : il se forme une \cle{liaison amide} $-CO-NH-$ avec élimination d'une molécule d'eau. Comme chaque monomère porte \emph{deux} fonctions, la chaîne grandit dans les deux sens :

\[ n\,\ce{HOOC-(CH2)4-COOH} + n\,\ce{H2N-(CH2)6-NH2} \xrightarrow{\text{chaleur}} \left[-\ce{CO-(CH2)4-CO-NH-(CH2)6-NH}-\right]_n + 2n\,\ce{H2O} \]

\begin{popfigure}
\begin{center}
$n$ \chemfig{HOOC-[:-30]CH_2-[:30]CH_2-[:-30]CH_2-[:30]CH_2-COOH}
\quad $+\;n$ \chemfig{H_2N-[:-30]CH_2-[:30]CH_2-[:-30]CH_2-[:30]CH_2-[:-30]CH_2-[:30]CH_2-NH_2}
\quad $\xrightarrow{\;\Delta\;}$
\vspace{8pt}

\chemfig{-[2]C(=[1]O)-[:-30]CH_2-[:30]CH_2-[:-30]CH_2-[:30]CH_2-C(=[1]O)-NH-[:-30]CH_2-[:30]CH_2-[:-30]CH_2-[:30]CH_2-[:-30]CH_2-[:30]CH_2-NH-}
\quad $+\;2n\,\ce{H2O}$
\legende{Formation du nylon 6,6 : l'unité répétée (motif) contient deux liaisons amide $-CO-NH-$ ; $n$ est l'indice de polymérisation.}
\end{center}
\end{popfigure}

\begin{attention}[Le piège du nombre de molécules d'eau]
L'erreur la plus fréquente au Baccalauréat est d'écrire $n\,\ce{H2O}$ au lieu de $2n\,\ce{H2O}$. Raisonne ainsi : chacun des $n$ diacides possède \textbf{deux} groupes $-COOH$, et chacun des $n$ diamines possède \textbf{deux} groupes $-NH_2$. Il se forme donc $2n$ liaisons amide, et chaque liaison amide élimine \textbf{une} molécule d'eau : au total, $2n\,\ce{H2O}$. Vérifie toujours : \emph{nombre de liaisons formées = nombre de petites molécules éliminées}.
\end{attention}

\begin{savaistu}[Pourquoi « 6,6 » ?]
Le nom « nylon 6,6 » vient tout simplement du nombre d'atomes de carbone de chaque monomère : le premier chiffre correspond à la \textbf{diamine} (6 carbones dans $\ce{H2N-(CH2)6-NH2}$), le second au \textbf{diacide} (6 carbones dans $\ce{HOOC-(CH2)4-COOH}$). Il existe aussi le nylon 6, obtenu à partir d'un seul monomère à 6 carbones, le caprolactame. Inventé en 1935 par Wallace Carothers chez DuPont, le nylon a révolutionné le textile. Au Cameroun, les filets de pêche des pêcheurs de Limbe, Kribi et Douala, ainsi que de nombreuses cordes, sont en nylon : il est léger, très résistant à la traction et ne pourrit pas dans l'eau de mer, contrairement aux fibres végétales traditionnelles.
\end{savaistu}

\courssub{Le tergal : un polyester}

\textbf{Les monomères.} Le tergal (appelé aussi \emph{Dacron} ou \emph{PET}, polyéthylène téréphtalate) est obtenu par polycondensation de :
\begin{itemize}
\item l'\cle{acide téréphtalique} (acide benzène-1,4-dicarboxylique) : $\ce{HOOC-C6H4-COOH}$, un diacide aromatique dont les deux groupes $-COOH$ sont en positions \emph{para} sur le cycle benzénique ;
\item l'\cle{éthane-1,2-diol} (ou \emph{éthylèneglycol}) : $\ce{HO-CH2-CH2-OH}$.
\end{itemize}

\textbf{La réaction.} C'est une estérification répétée : chaque groupe $-COOH$ réagit avec un groupe $-OH$ pour former une \cle{liaison ester} $-CO-O-$, avec élimination d'eau. La réaction est conduite à chaud (vers $\SI{250}{\celsius}$) en présence d'un catalyseur (souvent un dérivé d'antimoine ou de titane) :

\[ n\,\ce{HOOC-C6H4-COOH} + n\,\ce{HO-CH2-CH2-OH} \xrightarrow[\text{catalyseur}]{\Delta} \left[-\ce{CO-C6H4-CO-O-CH2-CH2-O}-\right]_n + 2n\,\ce{H2O} \]

\begin{popfigure}
\begin{center}
$n$ \chemfig{HOOC-[:30]*6(---(-COOH)---=)}
\quad $+\;n$ \chemfig{HO-[:-30]CH_2-[:30]CH_2-OH}
\quad $\xrightarrow[\text{catalyseur}]{\;\Delta\;}$
\vspace{8pt}

\chemfig{-[:30]C(=[1]O)-[:-30]*6(---(-C(=[1]O)-O-[:-30]CH_2-[:30]CH_2-O-)---=)}
\quad $+\;2n\,\ce{H2O}$
\legende{Formation du tergal (PET) : le motif encadré contient deux liaisons ester $-CO-O-$.}
\end{center}
\end{popfigure}

\begin{methode}[Écrire le motif d'un polycondensat]
Pour obtenir sans faute le motif et l'équation-bilan d'une polycondensation :
\begin{enumerate}
\item \textbf{Identifier} les deux fonctions réactives de chaque monomère (ex. $-COOH$ et $-NH_2$, ou $-COOH$ et $-OH$).
\item \textbf{Écrire} les deux monomères côte à côte, en plaçant face à face les fonctions qui vont réagir.
\item \textbf{Éliminer} une molécule d'eau entre les deux fonctions : $-COOH$ perd $-OH$, l'amine $-NH_2$ (ou l'alcool $-OH$) perd $-H$.
\item \textbf{Former} la liaison (amide $-CO-NH-$ ou ester $-CO-O-$) et prolonger la chaîne de part et d'autre par des tirets.
\item \textbf{Encadrer} l'unité répétée (le motif) et l'affecter de l'indice $n$.
\item \textbf{Équilibrer} en écrivant $2n\,\ce{H2O}$ si chaque monomère est bifonctionnel (deux liaisons formées par couple de monomères).
\end{enumerate}
\end{methode}

\courssub{Polyester ou polyamide : comment les distinguer ?}

\begin{aretenir}[Liaison caractéristique du polymère]
C'est la \textbf{nature de la liaison} créée lors de la polycondensation qui classe le polymère :
\begin{itemize}
\item si la liaison formée est une liaison \cle{ester} $-CO-O-$ (diacide + diol), le polymère est un \cle{polyester} : c'est le cas du \textbf{tergal} ;
\item si la liaison formée est une liaison \cle{amide} $-CO-NH-$ (diacide + diamine), le polymère est un \cle{polyamide} : c'est le cas du \textbf{nylon 6,6}.
\end{itemize}
Dans les deux cas, il s'agit d'une polycondensation avec élimination d'eau.
\end{aretenir}

\begin{center}
\begin{tabularx}{\linewidth}{|l|X|X|}
\hline
\rowcolor{popblueL} \textbf{Critère} & \textbf{Nylon 6,6 (polyamide)} & \textbf{Tergal (polyester)} \\
\hline
Monomères & acide hexanedioïque + hexaméthylènediamine & acide téréphtalique + éthane-1,2-diol \\
\hline
Liaison formée & amide $-CO-NH-$ & ester $-CO-O-$ \\
\hline
Petite molécule éliminée & $2n\,\ce{H2O}$ & $2n\,\ce{H2O}$ \\
\hline
Applications & cordes, filets de pêche, bas, pièces techniques & fibres textiles (pagnes), toiles, bouteilles plastiques \\
\hline
\end{tabularx}
\end{center}

Les applications sont partout autour de nous : les pagnes et vêtements en tergal vendus dans nos marchés (le tergal se froisse peu et sèche vite), les filets de pêche en nylon du littoral camerounais, les cordes, les ceintures de sécurité, et même les bouteilles d'eau en PET, qui est le même polymère que le tergal, moulé au lieu d'être étiré en fibre.

\begin{exoresolu}[Masse de motif et indice de polymérisation du nylon 6,6]
\textbf{Énoncé.} Le motif du nylon 6,6 a pour formule brute $\ce{C12H22N2O2}$. Une fibre de nylon a une masse molaire macromoléculaire moyenne $M = \SI{1,13e5}{g.mol^{-1}}$.
\begin{enumerate}
\item Calculer la masse molaire $M_0$ du motif.
\item En déduire l'indice de polymérisation moyen $n$ de cette fibre.
\item Quelle masse d'eau est éliminée lors de la formation d'une mole de ce polymère ?
\end{enumerate}
On donne : $M(\ce{C}) = \SI{12}{g.mol^{-1}}$, $M(\ce{H}) = \SI{1}{g.mol^{-1}}$, $M(\ce{N}) = \SI{14}{g.mol^{-1}}$, $M(\ce{O}) = \SI{16}{g.mol^{-1}}$, $M(\ce{H2O}) = \SI{18}{g.mol^{-1}}$.
\tcblower
\textbf{Solution détaillée.}
\begin{enumerate}
\item Masse molaire du motif :
\[ M_0 = 12\,M(\ce{C}) + 22\,M(\ce{H}) + 2\,M(\ce{N}) + 2\,M(\ce{O}) = 12\times 12 + 22\times 1 + 2\times 14 + 2\times 16 \]
\[ M_0 = 144 + 22 + 28 + 32 = \SI{226}{g.mol^{-1}}. \]
\item L'indice de polymérisation est le rapport de la masse du polymère à la masse du motif :
\[ n = \frac{M}{M_0} = \frac{\num{1,13e5}}{226} = 500. \]
La chaîne compte donc en moyenne \textbf{500 motifs}, soit 500 molécules de diacide et 500 molécules de diamine assemblées.
\item D'après l'équation-bilan, la formation de $n$ motifs élimine $2n$ molécules d'eau. Pour une mole de polymère, il s'élimine donc $2n = \SI{1000}{mol}$ d'eau, soit :
\[ m(\ce{H2O}) = 2n \times M(\ce{H2O}) = 1000 \times 18 = \SI{1,8e4}{g} = \SI{18}{kg}. \]
\end{enumerate}
\textbf{Remarque.} Ce calcul illustre bien la règle : pour le nylon 6,6, on élimine \emph{deux} moles d'eau par mole de motif formé, car chaque monomère est bifonctionnel.
\end{exoresolu}

\begin{exercice}[À toi de jouer]
Le tergal est obtenu à partir de l'acide téréphtalique $\ce{HOOC-C6H4-COOH}$ et de l'éthane-1,2-diol $\ce{HO-CH2-CH2-OH}$.
\begin{enumerate}
\item Écrire les formules semi-développées des deux monomères et entourer les fonctions réactives.
\item Écrire l'équation-bilan de la polycondensation et encadrer le motif.
\item Le tergal est-il un polyamide ou un polyester ? Justifier.
\item Le motif du tergal a pour formule brute $\ce{C10H8O4}$. Calculer sa masse molaire, puis l'indice de polymérisation d'un échantillon de masse molaire moyenne $\SI{1,92e4}{g.mol^{-1}}$.
\end{enumerate}
\end{exercice}

\begin{corrige}[Exercice : tergal]
\begin{enumerate}
\item Diacide : \chemfig{HOOC-[:30]*6(---(-COOH)---=)} (deux fonctions acide carboxylique) ; diol : \chemfig{HO-CH_2-CH_2-OH} (deux fonctions alcool).
\item $n\,\ce{HOOC-C6H4-COOH} + n\,\ce{HO-CH2-CH2-OH} \rightarrow [-\ce{CO-C6H4-CO-O-CH2-CH2-O}-]_n + 2n\,\ce{H2O}$.
\item C'est un \textbf{polyester}, car la liaison formée entre les monomères est une liaison \textbf{ester} $-CO-O-$.
\item $M_0 = 10\times 12 + 8\times 1 + 4\times 16 = 120 + 8 + 64 = \SI{192}{g.mol^{-1}}$ ; $n = \dfrac{\num{1,92e4}}{192} = 100$. La chaîne compte en moyenne 100 motifs.
\end{enumerate}
\end{corrige}

\begin{aretenir}[L'essentiel de la section]
\begin{itemize}
\item Une \cle{polycondensation} est une polymérisation de monomères \textbf{bifonctionnels} avec élimination d'une petite molécule (souvent \ce{H2O}) à chaque étape.
\item \textbf{Nylon 6,6} = acide hexanedioïque + hexaméthylènediamine $\rightarrow$ \textbf{polyamide} (liaisons $-CO-NH-$) $+\;2n\,\ce{H2O}$.
\item \textbf{Tergal} = acide téréphtalique + éthane-1,2-diol $\rightarrow$ \textbf{polyester} (liaisons $-CO-O-$) $+\;2n\,\ce{H2O}$.
\item Le nom « 6,6 » indique les 6 atomes de carbone de chaque monomère.
\item Ne jamais confondre avec la \cle{polyaddition}, qui ne produit aucun sous-produit.
\end{itemize}
\end{aretenir}

\coursec{La saponification : fabrication du savon}

Dans la section précédente, nous avons vu que la polycondensation permet d'assembler de petites molécules en immenses chaînes (nylon 6,6, tergal) en éliminant de l'eau. Nous allons maintenant faire le chemin inverse : \textbf{casser} des esters pour obtenir des produits utiles. Cette réaction de rupture, réalisée en milieu basique, porte un nom que tu connais sans le savoir depuis toujours : la \cle{saponification}. C'est elle qui transforme l'huile de palme de nos cuisines ou le beurre de karité du Nord-Cameroun en savon. C'est aussi la réaction qui clôt notre étude des dérivés d'acides carboxyliques — et l'une des plus anciennes réactions chimiques maîtrisées par l'humanité.

\courssub{Les corps gras : des triglycérides}

\textbf{Le glycérol et les acides gras}\par

Ouvre une bouteille d'huile d'arachide, observe l'huile de palme rouge utilisée pour le « ndolé » ou le beurre de karité : tous ces corps gras ont la même architecture moléculaire. Ce sont des \cle{triglycérides}, c'est-à-dire des \textbf{triesters} formés à partir :

\begin{itemize}
\item d'une molécule d'alcool à trois fonctions hydroxyle : le \cle{glycérol} (ou propane-1,2,3-triol), de formule \ce{HO-CH2-CH(OH)-CH2-OH} ;
\item de trois molécules d'\cle{acides gras}, c'est-à-dire des acides carboxyliques à longue chaîne carbonée (12 à 22 atomes de carbone), saturés ou insaturés, de formule générale $R-COOH$.
\end{itemize}

Quelques acides gras célèbres :

\begin{center}
\begin{tabularx}{\linewidth}{|l|X|l|}
\hline
\rowcolor{popblueL} \textbf{Acide gras} & \textbf{Formule semi-développée} & \textbf{Source} \\
\hline
Acide palmitique & \ce{CH3-(CH2)14-COOH} ($C_{16}$ saturé) & huile de palme \\
\hline
Acide oléique & \ce{CH3-(CH2)7-CH=CH-(CH2)7-COOH} ($C_{18}$, 1 insaturation) & huile d'olive, arachide \\
\hline
Acide stéarique & \ce{CH3-(CH2)16-COOH} ($C_{18}$ saturé) & beurre de karité \\
\hline
\end{tabularx}
\end{center}

\begin{definition}[Triglycéride]
Un \cle{triglycéride} est le triester obtenu par estérification des trois fonctions hydroxyle du glycérol par trois acides gras. Si les trois acides gras sont identiques, on parle de triglycéride \emph{simple} (palmitine, oléine, stéarine) ; s'ils sont différents, de triglycéride \emph{mixte} — le cas le plus fréquent dans les huiles naturelles.
\end{definition}

\begin{popfigure}
\[
\text{Palmitine : } \chemfig{CH_2(-[2]O-C(=[2]O)-[6]C_{15}H_{31})-CH(-[2]O-C(=[2]O)-[6]C_{15}H_{31})-CH_2(-[6]O-C(=[2]O)-[6]C_{15}H_{31})}
\]
\legende{Formule semi-développée de la palmitine, triglycéride simple du glycérol et de trois molécules d'acide palmitique $C_{15}H_{31}-COOH$. On repère le squelette du glycérol (à gauche) et les trois groupes ester $-COO-$.}
\end{popfigure}

\courssub{Définition de la saponification}

Tu as appris à la section 5 que l'hydrolyse acide d'un ester est une réaction \emph{limitée} : c'est l'inverse de l'estérification, et un équilibre s'établit. Pour casser \emph{totalement} un ester, il faut changer de stratégie : utiliser une \textbf{base forte} (soude \ce{NaOH} ou potasse \ce{KOH}) et chauffer.

\begin{definition}[Saponification]
La \cle{saponification} est l'hydrolyse basique d'un ester, c'est-à-dire l'action à chaud d'une base forte (\ce{NaOH} ou \ce{KOH}) sur un ester. Elle produit un \textbf{ion carboxylate} (sous forme de sel de sodium ou de potassium : le \cle{savon}) et un \textbf{alcool}. Contrairement à l'hydrolyse acide, elle est \textbf{totale} (non limitée par un équilibre), car l'ion carboxylate formé est une base faible qui ne réagit pratiquement pas avec l'alcool.
\end{definition}

Équation générale avec un ester simple :

\[ \mathrm{R{-}CO{-}O{-}R'} + \mathrm{NaOH} \xrightarrow{\text{chauffage}} \mathrm{R{-}CO{-}O^-}\,\mathrm{Na^+} + \mathrm{R'{-}OH} \]

Et avec un triglycéride, les \textbf{trois} fonctions ester sont hydrolysées : il faut donc \textbf{3 moles de soude} par mole de triglycéride, et l'alcool libéré est le glycérol :

\[ \ce{\text{triglycéride} + 3 NaOH ->[\Delta] 3 R-COO^- Na^+ + \text{glycérol}} \]

\begin{popfigure}
\[
\chemfig{CH_2(-[2]O-C(=[2]O)-[6]C_{15}H_{31})-CH(-[2]O-C(=[2]O)-[6]C_{15}H_{31})-CH_2(-[6]O-C(=[2]O)-[6]C_{15}H_{31})}
\; \ce{+ 3 NaOH ->[\Delta]} \;
3\, \chemfig{C_{15}H_{31}-C(=[2]O)-[6]O^{-}\,Na^{+}}
\; \ce{+} \;
\chemfig{CH_2(-[2]OH)-CH(-[2]OH)-CH_2(-[6]OH)}
\]
\legende{Saponification de la palmitine : 1 mole de triglycéride + 3 moles de soude donnent 3 moles de palmitate de sodium (savon) et 1 mole de glycérol. Réaction totale, réalisée à chaud.}
\end{popfigure}

\begin{methode}[Équilibrer une équation de saponification d'un triglycéride]
\begin{enumerate}
\item Repérer le nombre de fonctions ester du corps gras : un triglycéride en contient \textbf{3}.
\item Écrire les réactifs : 1 triglycéride + 3 \ce{NaOH} (ou 3 \ce{KOH}).
\item Écrire les produits : 3 ions carboxylate $R-COO^-$ (associés à \ce{Na+} ou \ce{K+}) + 1 glycérol.
\item Vérifier la conservation des atomes : le squelette $C_3H_5$ du glycérol se retrouve intact dans le produit ; chaque groupe $R-CO$ se retrouve dans un ion carboxylate.
\item Ne pas oublier de préciser les conditions : \textbf{chauffage} ($\Delta$) et soude \textbf{concentrée}.
\end{enumerate}
\end{methode}

\begin{attention}[Ne pas confondre saponification et hydrolyse acide]
C'est l'erreur classique du baccalauréat !
\begin{itemize}
\item \textbf{Hydrolyse acide} d'un ester : $ \mathrm{R{-}COO{-}R'} + H_2O \rightleftharpoons \mathrm{R{-}COOH} + \mathrm{R'{-}OH} $. Réaction \textbf{lente} et \textbf{limitée} (équilibre) ; elle produit un \textbf{acide carboxylique}.
\item \textbf{Saponification} : $ \mathrm{R{-}COO{-}R'} + NaOH \rightarrow \mathrm{R{-}COO^-}\,Na^+ + \mathrm{R'{-}OH} $. Réaction \textbf{totale} ; elle produit un \textbf{ion carboxylate} (savon).
\end{itemize}
Retenir : en milieu basique, l'acide carboxylique ne peut pas exister tel quel — il est immédiatement transformé en ion carboxylate, ce qui déplace la réaction jusqu'à épuisement de l'ester.
\end{attention}

\courssub{Fabrication artisanale du savon : protocole}

\begin{experience}[Préparation d'un savon à l'huile de palme]
\textbf{Dispositif et protocole :}
\begin{enumerate}
\item \textbf{Chauffage} : dans un ballon (ou une marmite, en artisanal), introduire l'huile de palme et une solution concentrée de soude en excès. Chauffer au reflux (ou à feu doux) en agitant pendant environ 30 minutes. Le mélange devient épais et pâteux : c'est la \emph{pâte à savon}.
\item \textbf{Relargage} : verser la pâte chaude dans une solution \textbf{saturée de chlorure de sodium} \ce{NaCl}. Le savon, très peu soluble dans l'eau salée, \textbf{précipite} en surface : on dit qu'il « relargue ». Le glycérol, lui, reste dissous dans la phase aqueuse avec l'excès de soude et le sel.
\item \textbf{Décantation et lavage} : recueillir le savon à l'aide d'une spatule (ou par filtration), le laver plusieurs fois à l'eau froide salée pour éliminer la soude résiduelle, puis le presser et le laisser sécher.
\end{enumerate}
\textbf{Observations :} l'huile, initialement insoluble dans la soude, disparaît progressivement ; après relargage, une masse solide jaunâtre (ou rougeâtre avec l'huile de palme rouge) surnage.
\textbf{Interprétation :} la réaction totale a consommé le triglycéride ; le relargage exploite la faible solubilité du savon dans l'eau salée pour le séparer du glycérol.
\end{experience}

\begin{popfigure}
\begin{tikzpicture}[font=\small, >=Stealth]
% Étape 1 : chauffage
\draw[thick, popblue] (0,0) -- (0,1.6) -- (1.6,1.6) -- (1.6,0);
\draw[thick, popblue] (0,0) arc[start angle=180, end angle=360, x radius=0.8, y radius=0.35];
\fill[poporangeL] (0,0.15) -- (0,1.0) -- (1.6,1.0) -- (1.6,0.15) arc[start angle=0, end angle=-180, x radius=0.8, y radius=0.3] -- cycle;
\draw[decorate, decoration={coil, amplitude=2pt, segment length=6pt}, poppink] (0.3,-0.5) -- (1.3,-0.5);
\node[below, popdark] at (0.8,-0.6) {\textbf{1. Chauffage}};
\node[below, popmuted] at (0.8,-1.05) {huile + soude, $\Delta$};
% flèche
\draw[->, very thick, poppurple] (2.0,0.7) -- (3.0,0.7);
% Étape 2 : relargage
\draw[thick, popblue] (3.4,0) -- (3.4,1.6) -- (5.0,1.6) -- (5.0,0);
\draw[thick, popblue] (3.4,0) arc[start angle=180, end angle=360, x radius=0.8, y radius=0.35];
\fill[popblueL] (3.4,0.15) -- (3.4,0.9) -- (5.0,0.9) -- (5.0,0.15) arc[start angle=0, end angle=-180, x radius=0.8, y radius=0.3] -- cycle;
\fill[popgoldL] (3.4,0.9) -- (3.4,1.25) -- (5.0,1.25) -- (5.0,0.9) -- cycle;
\node[popdark] at (4.2,1.07) {\scriptsize savon};
\node[popmuted] at (4.2,0.5) {\scriptsize eau salée};
\node[below, popdark] at (4.2,-0.6) {\textbf{2. Relargage}};
\node[below, popmuted] at (4.2,-1.05) {ajout de \ce{NaCl} saturé};
% flèche
\draw[->, very thick, poppurple] (5.4,0.7) -- (6.4,0.7);
% Étape 3 : décantation / savon
\draw[thick, popblue] (6.8,0) -- (6.8,1.6) -- (8.4,1.6) -- (8.4,0);
\draw[thick, popblue] (6.8,0) arc[start angle=180, end angle=360, x radius=0.8, y radius=0.35];
\fill[popblueL] (6.8,0.15) -- (6.8,0.8) -- (8.4,0.8) -- (8.4,0.15) arc[start angle=0, end angle=-180, x radius=0.8, y radius=0.3] -- cycle;
\fill[popgold] (7.0,0.85) rectangle (8.2,1.35);
\node[white] at (7.6,1.1) {\scriptsize savon};
\node[below, popdark] at (7.6,-0.6) {\textbf{3. Décantation}};
\node[below, popmuted] at (7.6,-1.05) {lavage, séchage};
\end{tikzpicture}
\legende{Les trois grandes étapes de la fabrication du savon : chauffage du corps gras avec la soude, relargage dans l'eau salée (le savon précipite en surface), puis décantation, lavage et séchage.}
\end{popfigure}

\begin{attention}[Sécurité et pièges pratiques]
\begin{itemize}
\item La soude concentrée est \textbf{corrosive} : gants, lunettes et blouse obligatoires. En cas de contact, rincer abondamment à l'eau.
\item Le savon obtenu ne doit pas être utilisé immédiatement : il peut contenir de la \textbf{soude en excès}, dangereuse pour la peau. Le lavage et le séchage (plusieurs semaines pour le savon artisanal) sont indispensables.
\item Piège de raisonnement : ne pas écrire \ce{R-COOH} comme produit — en milieu basique, c'est l'ion \ce{R-COO-} qui se forme.
\end{itemize}
\end{attention}

\courssub{Pourquoi le savon lave-t-il ? La propriété détersive}

L'ion carboxylate à longue chaîne, par exemple l'oléate de sodium \ce{CH3-(CH2)7-CH=CH-(CH2)7-COO^- Na+}, possède une structure \textbf{amphiphile} (qui « aime les deux ») :

\begin{itemize}
\item une \textbf{tête hydrophile} : le groupe ionique \ce{-COO-}, qui aime l'eau (il y est soluble grâce aux interactions ion-dipôle et aux liaisons hydrogène) ;
\item une \textbf{queue lipophile} (ou hydrophobe) : la longue chaîne carbonée $R$, qui fuit l'eau mais se lie aux graisses.
\end{itemize}

Plongés dans l'eau, les ions carboxylate se regroupent en \cle{micelles} : sphères dont les queues hydrophobes se cachent à l'intérieur (emprisonnant les salissures grasses) tandis que les têtes hydrophiles restent en contact avec l'eau. Les taches de graisse sont ainsi « emballées » et entraînées au rinçage : c'est la \textbf{propriété détersive} du savon.

\begin{popfigure}
\begin{tikzpicture}[font=\small]
% Ion carboxylate isolé
\draw[very thick, popgreen] (0,0.9) .. controls (0.5,1.1) and (1.0,0.7) .. (1.5,0.9) .. controls (1.9,1.05) and (2.2,0.8) .. (2.5,0.9);
\fill[popink] (2.75,0.9) circle (0.28);
\node[white] at (2.75,0.9) {\scriptsize \ce{COO-}};
\node[below, popgreen] at (1.2,0.55) {queue lipophile $R$};
\node[below, popink] at (2.9,0.45) {tête hydrophile};
% Micelle
\begin{scope}[xshift=6.2cm, yshift=1.6cm]
\fill[poporangeL] (0,0) circle (0.85);
\foreach \a in {0,30,...,330}{
  \draw[thick, popgreen] (\a:0.85) -- (\a:1.45);
  \fill[popink] (\a:1.6) circle (0.13);
}
\node[popdark] at (0,0) {\scriptsize graisse};
\node[popmuted] at (0,-2.05) {micelle dans l'eau};
\end{scope}
\end{tikzpicture}
\legende{À gauche : un ion carboxylate, avec sa queue carbonée lipophile et sa tête ionique hydrophile. À droite : une micelle — les queues capturent la graisse au centre, les têtes \ce{COO-} restent au contact de l'eau.}
\end{popfigure}

\begin{attention}[Limite des savons : les eaux dures]
Dans une eau « dure », riche en ions calcium \ce{Ca^2+} ou magnésium \ce{Mg^2+}, les savons \textbf{précipitent} sous forme de carboxylates de calcium ou de magnésium insolubles :
\[ \ce{2 R-COO^- + Ca^2+ -> (R-COO)2Ca v} \]
Ce dépôt (le « tartre » des lavabos) rend le savon inefficace : il ne mousse plus. C'est pourquoi on utilise dans ces régions des \textbf{détergents} synthétiques (alkylbenzènesulfonates), dont les sels de calcium restent solubles.
\end{attention}

\courssub{Exemple chiffré : rendement en savon}

\begin{exoresolu}[Saponification de l'oléine]
On saponifie à chaud $m = \SI{442}{g}$ d'oléine, triglycéride de l'acide oléique de masse molaire $M(\text{oléine}) \approx \SI{884}{g.mol^{-1}}$, par un excès de soude. On obtient de l'oléate de sodium \ce{C17H33-COO^-Na+} de masse molaire $M(\text{savon}) \approx \SI{304}{g.mol^{-1}}$. Calculer la masse de savon obtenue, la réaction étant totale.
\tcblower
\textbf{Étape 1 — Quantité de triglycéride :}
\[ n(\text{oléine}) = \frac{m}{M} = \frac{442}{884} = \SI{0,50}{mol} \]
\textbf{Étape 2 — Stœchiométrie :} d'après l'équation-bilan, 1 mole de triglycéride donne \textbf{3 moles} de savon :
\[ n(\text{savon}) = 3 \times 0{,}50 = \SI{1,5}{mol} \]
\textbf{Étape 3 — Masse de savon :}
\[ m(\text{savon}) = n \times M = 1{,}5 \times 304 = \SI{456}{g} \]
\textbf{Conclusion :} la saponification de \SI{442}{g} d'oléine fournit environ \SI{456}{g} d'oléate de sodium (et \SI{46}{g} de glycérol, sous-produit valorisable en pharmacie et en cosmétique). Remarque : la masse de savon dépasse celle du corps gras de départ car le sodium s'est substitué à une partie du squelette — c'est tout à fait normal.
\end{exoresolu}

\begin{savaistu}[Le savon noir et le savon de karité du Nord-Cameroun]
Bien avant les savonneries industrielles, les femmes du Nord-Cameroun fabriquaient le « savon noir » en chauffant des huiles végétales ou du beurre de karité avec de la \textbf{potasse} \ce{KOH}, obtenue par lixiviation des cendres de bois (on filtre l'eau à travers les cendres, puis on concentre la solution). C'est exactement la réaction de saponification étudiée ici : la potasse joue le rôle de base forte, et les savons de potassium obtenus sont plus mous et plus solubles que les savons de sodium. Le glycérol, non séparé dans le procédé artisanal, confère à ces savons leurs propriétés adoucissantes pour la peau — un savoir-faire traditionnel qui est aussi de la très bonne chimie !
\end{savaistu}

\begin{aretenir}[L'essentiel de la saponification]
\begin{itemize}
\item Les corps gras sont des \textbf{triglycérides} : triesters du glycérol et de trois acides gras $R-COOH$.
\item La \textbf{saponification} est l'hydrolyse basique d'un ester, à chaud : elle est \textbf{totale} et produit un \textbf{carboxylate de sodium ou de potassium (savon)} et un alcool.
\item Bilan type : 1 triglycéride + 3 \ce{NaOH} $\rightarrow$ 3 savons + 1 glycérol.
\item Fabrication : chauffage avec la soude, \textbf{relargage} au \ce{NaCl} saturé, décantation, lavage, séchage.
\item Le savon lave grâce à sa structure \textbf{amphiphile} : tête hydrophile \ce{COO-}, queue lipophile $R$, formant des \textbf{micelles}.
\item Limite : les savons précipitent dans les \textbf{eaux dures} (\ce{Ca^2+}, \ce{Mg^2+}).
\end{itemize}
\end{aretenir}

\newpage

\coursec{Travaux pratiques}

\courssub{Objectifs du TP}

Ce travail pratique te place dans la peau d'un artisan savonnier, comme ceux que l'on rencontre dans les quartiers de Douala ou de Bafoussam qui transforment l'huile de palme en savon ménager. À l'issue de la séance, tu seras capable de :

\begin{itemize}
\item réaliser la \cle{saponification} d'un corps gras (triglycéride) par la soude concentrée, à chaud ;
\item isoler le savon formé par \cle{relargage} dans une solution saturée de chlorure de sodium ;
\item vérifier expérimentalement les propriétés du savon : caractère basique, pouvoir moussant, influence de la dureté de l'eau ;
\item écrire l'équation-bilan de la réaction et identifier les produits formés (savon et glycérol).
\end{itemize}

\begin{aretenir}[Rappel théorique]
La saponification est l'action d'une base forte (\ce{NaOH} ou \ce{KOH}) sur un ester d'acide gras et de glycérol (triglycéride). C'est une réaction \textbf{totale} mais \textbf{lente}, réalisée à chaud :
\[ \ce{C3H5(COOR)3 + 3 NaOH -> 3 RCOONa + C3H5(OH)3} \]
\emph{R} représente la chaîne carbonée de l'acide gras (ex. \ce{C17H33} pour l'acide oléique). Les produits sont le \textbf{savon} (carboxylate de sodium) et le \textbf{glycérol} (propane-1,2,3-triol).
\end{aretenir}

\courssub{Matériel et produits}

\begin{center}
\begin{tabularx}{\linewidth}{|l|X|}
\hline
\rowcolor{popblueL} \textbf{Matériel} & \textbf{Produits} \\
\hline
Bécher de \SI{250}{mL} & Huile de palme rouge (ou huile d'arachide raffinée) : \SI{20}{mL} \\
\hline
Casserole ou bécher de \SI{500}{mL} (bain-marie) & Solution de soude \ce{NaOH} à environ \SI{10}{mol.L^{-1}} : \SI{20}{mL} \\
\hline
Agitateur en verre (ou spatule en bois) & Solution saturée de chlorure de sodium \ce{NaCl} : \SI{100}{mL} \\
\hline
Thermomètre, pince en bois & Eau distillée froide \\
\hline
Tamis fin ou tissu propre (pagne découpé) & Eau de robinet et eau distillée (pour les tests) \\
\hline
Papier pH ou pH-mètre, tubes à essais & Éthanol éventuel (quelques mL, accélérateur facultatif) \\
\hline
Gants, lunettes de protection, blouse & \\
\hline
\end{tabularx}
\end{center}

\begin{savaistu}[Alternatives locales]
Si le laboratoire ne dispose pas de soude en pastilles, on peut utiliser la \textbf{lessive de cendres} (filtrat concentré de cendres de bois, riche en potasse \ce{K2CO3}/\ce{KOH}) : c'est la méthode traditionnelle de fabrication du savon noir au village. Le savon obtenu est alors un savon \emph{mou} (carboxylate de potassium). De même, un réchaud à gaz ou même un feu de bois contrôlé peut remplacer la plaque chauffante, et un pagne en coton propre remplace avantageusement le papier filtre.
\end{savaistu}

\courssub{Sécurité}

\begin{attention}[La soude concentrée est dangereuse]
La solution de soude à \SI{10}{mol.L^{-1}} est \textbf{corrosive} (pictogramme : deux gouttes attaquant une main et une surface métallique). Elle provoque des brûlures chimiques graves de la peau et des yeux.
\begin{itemize}
\item Porte obligatoirement \textbf{gants, lunettes et blouse} pendant toute la manipulation.
\item Verse la soude \textbf{lentement} et sans éclaboussures ; ne jamais pencher le visage au-dessus du bécher.
\item En cas de contact avec la peau : rincer \textbf{abondamment à l'eau froide} pendant plusieurs minutes, puis prévenir l'enseignant.
\item Le chauffage se fait \textbf{au bain-marie} et non à la flamme directe, pour éviter les projections de pâte chaude et toute inflammation de l'huile.
\item Ne jamais goûter les produits ; se laver les mains en fin de séance.
\end{itemize}
\end{attention}

\courssub{Schéma du dispositif}

\begin{popfigure}
\begin{center}
\begin{tikzpicture}[scale=0.95, every node/.style={font=\small}]
% Bain-marie (casserole)
\draw[thick, popdark] (-3.2,0) -- (-3.2,2.2) -- (3.2,2.2) -- (3.2,0) -- cycle;
\draw[popblue!50, fill=popblue!15] (-3.1,0.05) rectangle (3.1,1.6);
\node[popblue] at (2.2,0.35) {eau du bain-marie};
% Bécher
\draw[thick, popdark] (-1.2,0.3) -- (-1.2,3.2) -- (1.2,3.2) -- (1.2,0.3);
\draw[thick, popdark] (-1.2,3.2) -- (-1.45,3.45);
\draw[thick, popdark] (1.2,3.2) -- (1.45,3.45);
% Pâte huile + soude
\draw[fill=poporange!45, draw=poporange] (-1.15,0.35) rectangle (1.15,1.9);
\node[align=center] at (0,1.1) {huile de palme\\ $+$ soude};
% Agitateur
\draw[thick, popdark, rotate=15] (0.35,0.5) rectangle (0.55,4.3);
\node[rotate=15, anchor=south] at (0.85,4.35) {agitateur};
% Thermomètre
\draw[thick, poppurple, rotate=-12] (-0.9,0.6) rectangle (-0.7,4.0);
\fill[poppink] (-0.83,0.75) circle (0.12);
% Chauffage
\foreach \x in {-2.4,-1.2,0,1.2,2.4}{
  \draw[decorate, decoration={snake, amplitude=0.12cm, segment length=0.35cm}, poppink, thick] (\x,-0.75) -- (\x,-0.05);
}
\draw[thick, popdark] (-3.4,-0.75) -- (3.4,-0.75);
\node at (0,-1.15) {plaque chauffante (ou réchaud) --- feu doux};
% Vapeur
\foreach \x in {-0.5,0.2,0.9}{
  \draw[decorate, decoration={snake, amplitude=0.08cm, segment length=0.25cm}, popmuted] (\x,3.6) -- (\x,4.4);
}
\end{tikzpicture}
\end{center}
\legende{Dispositif de saponification au bain-marie : le bécher contenant l'huile et la soude est chauffé doucement, sous agitation constante.}
\end{popfigure}

\courssub{Protocole}

\begin{methode}[Fabrication du savon d'huile de palme]
\begin{enumerate}
\item \textbf{Préparation du mélange.} Introduis dans le bécher de \SI{250}{mL} environ \SI{20}{mL} d'huile de palme, puis ajoute avec précaution \SI{20}{mL} de solution de soude à \SI{10}{mol.L^{-1}}. (Facultatif : ajoute \SI{5}{mL} d'éthanol pour homogénéiser le mélange et accélérer la réaction.)
\item \textbf{Chauffage.} Place le bécher au bain-marie et chauffe doucement (température du bain : environ \SI{80}{\celsius}) pendant \textbf{20 à 30 minutes}, en \textbf{agitant sans interruption} avec la spatule. Le mélange s'épaissit progressivement.
\item \textbf{Fin de réaction.} Arrête le chauffage lorsque tu obtiens une \textbf{pâte épaisse et homogène}, de couleur orangée (l'huile de palme colore naturellement le savon).
\item \textbf{Relargage.} Verse la pâte chaude dans un bécher contenant \SI{100}{mL} de solution saturée de \ce{NaCl}. Agite doucement : le savon, \textbf{insoluble dans l'eau salée}, précipite et surnage.
\item \textbf{Filtration et lavage.} Filtre sur le tamis ou le tissu. Lave le savon recueilli avec un peu d'\textbf{eau distillée froide} pour éliminer l'excès de soude, puis presse-le dans le tissu.
\item \textbf{Séchage.} Dépose le savon sur du papier absorbant et laisse sécher à l'air libre (quelques heures à quelques jours).
\item \textbf{Tests.} 
\begin{itemize}
\item Dissous un fragment de savon dans un tube à essais contenant de l'eau distillée ; agite et observe la mousse.
\item Mesure le pH de cette solution avec le papier pH.
\item Prépare deux tubes identiques, l'un avec de l'\textbf{eau distillée}, l'autre avec de l'\textbf{eau de robinet} ; ajoute la même quantité de savon, agite de la même façon et compare la hauteur de mousse.
\end{itemize}
\end{enumerate}
\end{methode}

\courssub{Observations et résultats attendus}

\begin{center}
\begin{tabularx}{\linewidth}{|X|l|}
\hline
\rowcolor{popblueL} \textbf{Étape / test} & \textbf{Observation attendue} \\
\hline
Chauffage huile + soude & Le mélange, d'abord hétérogène (deux phases), devient une pâte épaisse et homogène après 20--30 min \\
\hline
Relargage dans \ce{NaCl} saturé & Un dépôt solide orangé se forme et surnage à la surface du liquide \\
\hline
Masse de savon sec obtenue & environ 15 à \SI{20}{g} \\
\hline
Dissolution dans l'eau distillée & Solution trouble, formation d'une mousse abondante et persistante \\
\hline
pH de la solution de savon & pH $\approx 9$ à $10$ (solution basique) \\
\hline
Moussage en eau de robinet & Mousse nettement moins abondante qu'en eau distillée ; aspect parfois trouble \\
\hline
\end{tabularx}
\end{center}

\courssub{Exploitation}

\begin{exoresolu}[Questions d'exploitation du TP]
\textbf{1.} Écris l'équation-bilan de la saponification du triglycéride de l'acide oléique (oléine) par la soude. \par
\textbf{2.} Pourquoi chauffe-t-on le mélange et pourquoi agite-t-on constamment ? \par
\textbf{3.} Quel est le rôle du relargage dans la solution saturée de \ce{NaCl} ? \par
\textbf{4.} Pourquoi le pH de la solution de savon est-il supérieur à 7 ? \par
\textbf{5.} Explique pourquoi le savon mousse moins bien dans l'eau de robinet que dans l'eau distillée.
\tcblower
\textbf{1.} L'oléine est le triester du glycérol et de l'acide oléique (\ce{C17H33COOH}). La soude rompt les trois liaisons ester :
\[ \ce{C3H5(COOC17H33)3 + 3 NaOH -> 3 C17H33COONa + C3H5(OH)3} \]
Les produits sont l'\textbf{oléate de sodium} (le savon) et le \textbf{glycérol} (propane-1,2,3-triol), qui reste dissous dans l'eau salée. \par
\textbf{2.} La saponification est une réaction \textbf{lente} : le chauffage augmente la vitesse de réaction (facteur cinétique). L'agitation maintient le contact entre l'huile et la soude, deux phases peu miscibles au début, et homogénéise la température. \par
\textbf{3.} Le savon est \textbf{très peu soluble dans l'eau salée} : l'excès d'ions \ce{Na+} diminue la solubilité des carboxylates (effet d'ion commun / relargage) et provoque sa précipitation. Le relargage permet donc de \textbf{séparer le savon} du glycérol et de l'excès de soude, qui restent en solution. \par
\textbf{4.} L'ion carboxylate \ce{C17H33COO-} est la \textbf{base conjuguée} d'un acide faible (l'acide oléique). Il réagit partiellement avec l'eau (hydrolyse basique) :
\[ \ce{C17H33COO- + H2O <=> C17H33COOH + OH-} \]
La formation d'ions hydroxyde \ce{OH-} rend la solution basique (pH $\approx 9$ à $10$). \par
\textbf{5.} L'eau de robinet est une \textbf{eau dure} : elle contient des ions \ce{Ca^2+} et \ce{Mg^2+}. Ces ions forment avec les carboxylates des sels \textbf{insolubles} (précipité grisâtre, « savon de chaux ») qui ne moussent pas :
\[ \ce{2 C17H33COO- + Ca^2+ -> (C17H33COO)2Ca (s)} \]
Une partie du savon est ainsi consommée inutilement : la mousse diminue. L'eau distillée, exempte de ces ions, permet un moussage maximal.
\end{exoresolu}

\courssub{Conclusion}

Au cours de ce TP, tu as transformé un corps gras naturel — l'huile de palme — en \cle{savon} par saponification à chaud avec la soude concentrée. La réaction, \textbf{totale mais lente}, a produit un mélange de carboxylates de sodium et de glycérol. Le \textbf{relargage} dans l'eau salée a permis d'isoler le savon par précipitation, puis le lavage à l'eau distillée froide a éliminé l'excès de soude. Les tests ont confirmé les propriétés essentielles du savon : une solution \textbf{basique} (pH $\approx 9$ à $10$), un bon \textbf{pouvoir moussant} en eau douce, fortement réduit en eau dure à cause des ions \ce{Ca^2+} et \ce{Mg^2+}. Cette démarche reproduit exactement le procédé des savonneries artisanales camerounaises et constitue une application industrielle majeure des réactions des dérivés d'acides carboxyliques étudiées dans cette leçon.

\newpage

\coursec{Méthodes et exercices résolus}

\courssub{Méthode 1 : Écrire la formule semi-développée d'un acide carboxylique}

\begin{methode}[Formules semi-développées des acides carboxyliques]
\begin{enumerate}
\item Repérer le \cle{groupe carboxyle} $-\mathrm{COOH}$, toujours placé en \textbf{extrémité} de chaîne : le carbone du carboxyle porte le numéro 1.
\item Construire la chaîne carbonée principale en comptant le carbone du groupe $-\mathrm{COOH}$ dans le nombre total d'atomes de carbone.
\item Pour un \cle{monoacide saturé}, appliquer la formule générale $C_nH_{2n+1}COOH$ (soit $C_{n+1}H_{2(n+1)}O_2$).
\item Pour un \cle{acide insaturé}, placer la double liaison $\mathrm{C=C}$ à la position indiquée et retirer deux hydrogènes par double liaison.
\item Pour un \cle{polyacide} (diacide), placer un groupe $-\mathrm{COOH}$ à chaque extrémité : $\mathrm{HOOC-(CH_2)}_n\mathrm{-COOH}$.
\item Vérifier la \textbf{tétravalence} de chaque carbone : chaque carbone doit former exactement 4 liaisons.
\end{enumerate}
\textbf{Réflexe Bac :} un acide à $N$ atomes de carbone au total possède $N-1$ carbones dans le reste alkyle R de R-COOH.
\end{methode}

\begin{exoresolu}[Écriture de formules semi-développées]
Écrire les formules semi-développées des composés suivants :
\begin{enumerate}
\item l'acide butanoïque ;
\item l'acide prop-2-énoïque (acide acrylique) ;
\item l'acide hexanedioïque (acide adipique), utilisé dans la fabrication du nylon 6,6.
\end{enumerate}
\tcblower
\textbf{1. Acide butanoïque.} Le suffixe « butan » indique une chaîne de \textbf{4 atomes de carbone au total}, carbone du carboxyle compris. Le reste alkyle compte donc $4-1=3$ carbones :
\[ \chemfig{CH_3-CH_2-CH_2-C(=[1]O)-[7]OH} \qquad \text{soit} \qquad \ce{CH3-CH2-CH2-COOH} \]
Vérification avec la formule générale des monoacides saturés : $n=3$ donne $C_3H_7COOH$, soit $C_4H_8O_2$. Correct.

\textbf{2. Acide prop-2-énoïque.} Chaîne de 3 carbones ; la double liaison est entre C2 et C3 (le carboxyle porte le numéro 1) :
\[ \chemfig{CH_2=CH-C(=[1]O)-[7]OH} \qquad \text{soit} \qquad \ce{CH2=CH-COOH} \]
C'est l'acide \cle{acrylique}, monomère de nombreux polymères (peintures, vernis).

\textbf{3. Acide hexanedioïque.} Le préfixe « hexan » indique 6 carbones ; « dioïque » indique \textbf{deux} groupes carboxyles, nécessairement aux deux extrémités C1 et C6. Entre eux : $6-2=4$ groupes $-\ce{CH2}-$ :
\[ \chemfig{HO-C(=[1]O)-[7]CH_2-CH_2-CH_2-CH_2-C(=[1]O)-[7]OH} \qquad \text{soit} \qquad \ce{HOOC-(CH2)4-COOH} \]
\textbf{Conclusion :} acide butanoïque : \ce{CH3-CH2-CH2-COOH} ; acide prop-2-énoïque : \ce{CH2=CH-COOH} ; acide hexanedioïque : \ce{HOOC-(CH2)4-COOH}.
\end{exoresolu}

\courssub{Méthode 2 : Nommer les acides carboxyliques et leurs dérivés}

\begin{methode}[Nomenclature des acides carboxyliques et dérivés]
\begin{enumerate}
\item \textbf{Acide carboxylique} : choisir la chaîne la plus longue contenant le groupe $-\mathrm{COOH}$ ; le carbone du carboxyle est le C1. Nom = nom de l'alcane correspondant dont on remplace le « e » final par « \textbf{oïque} », précédé du mot « acide » : \emph{acide + alcane-e + oïque}.
\item Numéroter les ramifications à partir du carboxyle (ex. acide 2-méthylpropanoïque).
\item \textbf{Chlorure d'acyle} R-COCl : « \emph{chlorure de} + nom de l'acyle » ; la terminaison « oïque » devient « \textbf{oyle} » : chlorure d'éthanoyle, chlorure de propanoyle.
\item \textbf{Anhydride d'acide} R-CO-O-CO-R' : « \emph{anhydride} + nom de l'acide » dont on supprime le mot « acide » : anhydride éthanoïque ; si les deux groupes diffèrent, citer les deux par ordre alphabétique (anhydride éthanoïque propanoïque).
\item \textbf{Ester} R-COO-R' : « \emph{alkanoate (nom de l'acide, oïque $\to$ oate) de} + nom du groupe alkyle R' » : éthanoate d'éthyle.
\item \textbf{Amide} R-CONH$_2$ : terminaison « oïque » $\to$ « \textbf{amide} » : éthanamide.
\end{enumerate}
\textbf{Réflexe Bac :} toujours vérifier que la chaîne choisie contient bien le carbone du carboxyle.
\end{methode}

\begin{exoresolu}[Nommer et identifier des dérivés d'acides]
\begin{enumerate}
\item Nommer les composés : \ce{CH3-CH(CH3)-COOH} ; \ce{CH3-CH2-COCl} ; \ce{CH3-CO-O-CO-CH3} ; \ce{CH3-COO-CH2-CH3}.
\item Écrire la formule semi-développée du propanoate de méthyle et du chlorure de 2-méthylpropanoyle.
\end{enumerate}
\tcblower
\textbf{1. Nomenclature.}
\begin{itemize}
\item \ce{CH3-CH(CH3)-COOH} : chaîne principale de 3 carbones contenant le carboxyle (acide propanoïque), avec un méthyle sur C2 : \textbf{acide 2-méthylpropanoïque}.
\item \ce{CH3-CH2-COCl} : dérive de l'acide propanoïque : \textbf{chlorure de propanoyle}.
\item \ce{CH3-CO-O-CO-CH3} : deux groupes éthanoyle identiques : \textbf{anhydride éthanoïque}.
\item \ce{CH3-COO-CH2-CH3} : dérive de l'acide éthanoïque, groupe alkyle éthyle : \textbf{éthanoate d'éthyle}.
\end{itemize}
\textbf{2. Formules.}
\begin{itemize}
\item Propanoate de méthyle : acide propanoïque \ce{CH3-CH2-COOH} + groupe méthyle : \ce{CH3-CH2-COO-CH3}.
\item Chlorure de 2-méthylpropanoyle : \ce{CH3-CH(CH3)-COCl}, soit $\chemfig{CH_3-CH(-[2]CH_3)-C(=[1]O)-[7]Cl}$.
\end{itemize}
\textbf{Conclusion :} la clé est d'identifier le « bloc » acide (R-CO) qui donne la racine du nom, puis le « bloc » substituant (Cl, OR', OCOR', NH$_2$) qui fixe la fonction.
\end{exoresolu}

\courssub{Méthode 3 : Expliquer les propriétés physiques par les liaisons hydrogène}

\begin{methode}[Liaisons hydrogène et propriétés physiques]
\begin{enumerate}
\item Repérer dans la molécule un hydrogène porté par un atome très électronégatif (O ou N) : c'est le cas du groupe $-\mathrm{OH}$ du carboxyle.
\item Rappeler qu'une \cle{liaison hydrogène} s'établit entre cet H (déficitaire en électrons, $\delta+$) et le doublet d'un atome électronégatif (O, $\delta-$) d'une molécule voisine.
\item En déduire les conséquences :
\begin{itemize}
\item \textbf{températures de fusion et d'ébullition élevées} : il faut fournir de l'énergie pour rompre les liaisons hydrogène (les acides forment même des \textbf{dimères} à deux liaisons hydrogène) ;
\item \textbf{solubilité dans l'eau} des acides à chaîne courte ($\leq 4$ carbones) : liaisons hydrogène avec l'eau ;
\item \textbf{solubilité décroissante} quand la chaîne carbonée (partie hydrophobe) s'allonge.
\end{itemize}
\item Comparer toujours à masse molaire voisine avec un composé ne formant pas de liaisons hydrogène (alcane, ester).
\end{enumerate}
\textbf{Réflexe Bac :} « plus de liaisons hydrogène $\Rightarrow$ plus de cohésion $\Rightarrow$ $T_{eb}$ plus élevée ».
\end{methode}

\begin{exoresolu}[Comparaison de températures d'ébullition]
On donne les températures d'ébullition : propane $\theta_{eb} = \SI{-42}{\celsius}$ ; éthanal $\theta_{eb} = \SI{21}{\celsius}$ ; éthanol $\theta_{eb} = \SI{78}{\celsius}$ ; acide méthanoïque $\theta_{eb} = \SI{101}{\celsius}$. Ces quatre composés ont des masses molaires voisines (44 à \SI{46}{g.mol^{-1}}).
\begin{enumerate}
\item Expliquer l'augmentation régulière des températures d'ébullition.
\item L'acide méthanoïque est miscible à l'eau en toutes proportions. Justifier.
\end{enumerate}
\tcblower
\textbf{1. Interprétation.}
\begin{itemize}
\item \textbf{Propane} \ce{C3H8} : molécule apolaire ; seules de faibles interactions de Van der Waals assurent la cohésion : ébullition très basse (\SI{-42}{\celsius}).
\item \textbf{Éthanal} \ce{CH3CHO} : molécule \textbf{polaire} (liaison C=O) mais sans H lié à O : interactions dipôle-dipôle, plus fortes que Van der Waals : $\theta_{eb} = \SI{21}{\celsius}$.
\item \textbf{Éthanol} \ce{C2H5OH} : possède un groupe $-\mathrm{OH}$ : des \cle{liaisons hydrogène} s'établissent entre molécules : cohésion renforcée, $\theta_{eb} = \SI{78}{\celsius}$.
\item \textbf{Acide méthanoïque} \ce{HCOOH} : le groupe carboxyle contient à la fois C=O (accepteur) et O-H (donneur) : les molécules s'associent en \textbf{dimères} par \emph{deux} liaisons hydrogène, formant des « pseudo-molécules » de masse double :
\[ \chemfig{H-C(=[1]O)(-[7]O-H@{a})} \quad \cdots \quad \text{dimère cyclique à 2 liaisons H} \]
D'où la température d'ébullition la plus élevée : \SI{101}{\celsius}.
\end{itemize}
\textbf{2. Miscibilité avec l'eau.} L'acide méthanoïque établit des liaisons hydrogène \textbf{avec les molécules d'eau} (donneur via O-H, accepteur via C=O). Sa chaîne carbonée hydrophobe est quasi inexistante (un seul carbone) : rien ne s'oppose au mélange, d'où une miscibilité totale.
\textbf{Conclusion :} les liaisons hydrogène, intermoléculaires (ébullition) ou avec l'eau (solubilité), expliquent quantitativement les propriétés physiques des acides carboxyliques.
\end{exoresolu}

\courssub{Méthode 4 : Exploiter le caractère acide (réaction avec l'eau, pH, $K_A$)}

\begin{methode}[Acide carboxylique et eau : équation, $K_A$, pH]
\begin{enumerate}
\item Écrire l'équation de la réaction, \textbf{limitée} (équilibre), avec l'eau :
\[ \ce{RCOOH + H2O <=> RCOO- + H3O+} \]
\item Identifier les couples : \ce{RCOOH}/\ce{RCOO-} et \ce{H3O+}/\ce{H2O}.
\item Écrire la constante d'acidité : $K_A = \dfrac{[\ce{RCOO-}][\ce{H3O+}]}{[\ce{RCOOH}]}$ et $pK_A = -\log K_A$.
\item Pour comparer deux acides : \textbf{plus $pK_A$ est petit (plus $K_A$ est grand), plus l'acide est fort} (plus dissocié).
\item Pour un calcul de pH d'acide faible peu dilué : $[\ce{H3O+}] \approx \sqrt{K_A \cdot C}$, soit $pH \approx \tfrac{1}{2}\left(pK_A - \log C\right)$, à valider par l'hypothèse $[\ce{H3O+}] \ll C$.
\end{enumerate}
\textbf{Réflexe Bac :} un acide carboxylique est un acide \textbf{faible} : la flèche est toujours \ce{<=>}, jamais \ce{->}.
\end{methode}

\begin{exoresolu}[pH d'une solution d'acide éthanoïque]
Le vinaigre de table consommé à Douala est une solution d'acide éthanoïque \ce{CH3COOH}. On prépare une solution S de concentration $C = \SI{1,0e-2}{mol.L^{-1}}$. On donne $pK_A(\ce{CH3COOH}/\ce{CH3COO-}) = 4,8$.
\begin{enumerate}
\item Écrire l'équation de la réaction de l'acide éthanoïque avec l'eau.
\item Calculer le pH de la solution S et vérifier l'approximation utilisée.
\item Comparer la force de l'acide éthanoïque à celle de l'acide méthanoïque ($pK_A = 3,8$).
\end{enumerate}
\tcblower
\textbf{1. Équation de la réaction.}
\[ \ce{CH3COOH + H2O <=> CH3COO- + H3O+} \]
Couples mis en jeu : \ce{CH3COOH}/\ce{CH3COO-} et \ce{H3O+}/\ce{H2O}.

\textbf{2. Calcul du pH.} L'acide est faible et peu dissocié ; dressons le tableau d'avancement volumique (en \si{mol.L^{-1}}) :
\begin{center}
\begin{tabularx}{\linewidth}{|l|X|X|X|X|}
\hline
\rowcolor{popblueL}
État & \ce{CH3COOH} & \ce{H2O} & \ce{CH3COO-} & \ce{H3O+} \\
\hline
Initial & $C$ & excès & $0$ & $0$ \\
\hline
Équilibre & $C - x$ & excès & $x$ & $x$ \\
\hline
\end{tabularx}
\end{center}
\[
K_A = \frac{[\ce{CH3COO-}][\ce{H3O+}]}{[\ce{CH3COOH}]} = \frac{x^2}{C - x} \approx \frac{x^2}{C}
\quad\Rightarrow\quad
x = [\ce{H3O+}] \approx \sqrt{K_A \cdot C}
\]
Application numérique : $K_A = 10^{-4,8} = \SI{1,58e-5}{}$ :
\[
[\ce{H3O+}] = \sqrt{1,58\times 10^{-5} \times 1,0\times 10^{-2}} = \sqrt{1,58\times 10^{-7}} = \SI{4,0e-4}{mol.L^{-1}}
\]
\[
\boxed{pH = -\log(4,0\times 10^{-4}) = 3,4}
\]
Vérification : $[\ce{H3O+}]/C = \dfrac{4,0\times 10^{-4}}{1,0\times 10^{-2}} = 4\ \% \ll 1$ : l'approximation $C - x \approx C$ est \textbf{justifiée} (acide peu dissocié).

\textbf{3. Comparaison.} $pK_A(\ce{HCOOH}) = 3,8 < pK_A(\ce{CH3COOH}) = 4,8$. À concentration égale, l'acide méthanoïque est \textbf{plus dissocié}, donc \textbf{plus fort} que l'acide éthanoïque.
\textbf{Conclusion :} la solution S a un pH de 3,4 ; l'acide éthanoïque est un acide faible, moins fort que l'acide méthanoïque.
\end{exoresolu}

\courssub{Méthode 5 : Écrire la formation des dérivés (chlorure d'acyle, anhydride) et la synthèse des esters}

\begin{methode}[Équations de formation des dérivés et des esters]
\begin{enumerate}
\item \textbf{Chlorure d'acyle} à partir de l'acide : réactif chlorurant $\ce{SOCl2}$ (le plus pratique) ou $\ce{PCl5}$ :
\[ \mathrm{RCOOH} + \ce{SOCl2} \rightarrow \mathrm{RCOCl} + \ce{SO2} + \ce{HCl} \]
\item \textbf{Anhydride} par déshydratation de deux acides (chauffage, $\ce{P4O10}$) ou action d'un chlorure d'acyle sur un carboxylate :
\[ \mathrm{RCOOH} + \mathrm{R'COOH} \xrightarrow{\ce{P4O10}} \mathrm{R-CO-O-CO-R'} + \ce{H2O} \]
\item \textbf{Estérification directe} (acide + alcool) : \textbf{lente, limitée} (équilibre), catalysée par $\ce{H+}$ :
\[ \mathrm{RCOOH} + \mathrm{R'OH} \rightleftharpoons \mathrm{RCOOR'} + \ce{H2O} \]
\item \textbf{Estérification indirecte} via chlorure d'acyle ou anhydride : \textbf{rapide et totale} :
\[ \mathrm{RCOCl} + \mathrm{R'OH} \rightarrow \mathrm{RCOOR'} + \ce{HCl} \]
\[ \mathrm{R-CO-O-CO-R} + \mathrm{R'OH} \rightarrow \mathrm{RCOOR'} + \mathrm{RCOOH} \]
\item Vérifier la conservation des atomes de chaque côté de l'équation (bilan matière).
\end{enumerate}
\textbf{Réflexe Bac :} pour un rendement élevé en ester, le jury attend le passage par le \cle{chlorure d'acyle} ou l'\cle{anhydride}, pas l'estérification directe.
\end{methode}

\begin{exoresolu}[Synthèse de l'éthanoate de butyle]
L'éthanoate de butyle est un ester à odeur de banane, utilisé comme arôme. On veut le préparer efficacement au laboratoire.
\begin{enumerate}
\item Écrire l'équation de sa formation par estérification directe et donner deux caractéristiques de cette réaction.
\item Écrire deux équations de synthèse rapide et totale de cet ester, à partir : a) du chlorure d'éthanoyle ; b) de l'anhydride éthanoïque.
\item On fait réagir $n_1 = \SI{0,20}{mol}$ d'anhydride éthanoïque avec un excès de butan-1-ol. La réaction étant totale, calculer la masse d'ester obtenue. Donnée : $M(\text{ester}) = \SI{116}{g.mol^{-1}}$.
\end{enumerate}
\tcblower
\textbf{1. Estérification directe.}
\[ \ce{CH3COOH + CH3-CH2-CH2-CH2OH <=> CH3COO-CH2-CH2-CH2-CH3 + H2O} \]
Caractéristiques : réaction \textbf{lente}, \textbf{limitée} (équilibre, rendement d'environ $67\ \%$ avec un alcool primaire en mélange équimolaire) et \textbf{athermique}.

\textbf{2. Synthèses rapides et totales.}

a) Avec le chlorure d'éthanoyle :
\[ \ce{CH3COCl + CH3CH2CH2CH2OH -> CH3COOCH2CH2CH2CH3 + HCl} \]

b) Avec l'anhydride éthanoïque :
\[ \ce{CH3-CO-O-CO-CH3 + CH3CH2CH2CH2OH -> CH3COOCH2CH2CH2CH3 + CH3COOH} \]
Ces deux réactions sont \textbf{rapides} et \textbf{totales} (flèche simple \ce{->}).

\textbf{3. Masse d'ester.} D'après l'équation b), les coefficients stœchiométriques valent 1 : 1 mole d'anhydride donne 1 mole d'ester. Le butan-1-ol étant en excès, l'anhydride est le réactif limitant :
\begin{center}
\begin{tabularx}{\linewidth}{|l|X|X|X|X|}
\hline
\rowcolor{popblueL}
 & anhydride & butan-1-ol & ester & acide \\
\hline
Initial (mol) & $0,20$ & excès & $0$ & $0$ \\
\hline
Final (mol) & $0$ & excès & $0,20$ & $0,20$ \\
\hline
\end{tabularx}
\end{center}
\[
n(\text{ester}) = 0,20\ \text{mol} \quad\Rightarrow\quad m = n \times M = 0,20 \times 116 = \boxed{\SI{23,2}{g}}
\]
\textbf{Conclusion :} la synthèse totale à partir de l'anhydride éthanoïque fournit \SI{23,2}{g} d'éthanoate de butyle, alors que l'estérification directe n'en donnerait qu'environ $0,67 \times 23,2 \approx \SI{15,5}{g}$ en mélange équimolaire.
\end{exoresolu}

\courssub{Méthode 6 : Écrire une équation de polycondensation (nylon 6,6, tergal)}

\begin{methode}[Polycondensation : méthode de rédaction]
\begin{enumerate}
\item Identifier les deux \cle{monomères} : chacun doit porter \textbf{deux groupes fonctionnels} (diacide + diamine pour un polyamide ; diacide + diol pour un polyester).
\item Repérer la réaction d'assemblage : formation d'une liaison \textbf{amide} $-\mathrm{CO-NH}-$ (avec élimination de $\ce{H2O}$) ou \textbf{ester} $-\mathrm{CO-O}-$ (avec élimination de $\ce{H2O}$).
\item Écrire l'équation-bilan avec $n$ motifs de chaque monomère et compter les molécules d'eau éliminées : pour $n$ + $n$ monomères formant une chaîne, on élimine environ $2n$ molécules d'eau ($2n-1$ exactement, souvent arrondi à $2n$ au lycée).
\item Encadrer le \cle{motif} (unité répétitive) entre parenthèses avec l'indice $n$.
\item Vérifier la conservation des atomes.
\end{enumerate}
\textbf{Réflexe Bac :} « polycondensation = polymérisation + élimination d'une petite molécule (eau) » ; pas de double liaison consommée, contrairement à la polyaddition.
\end{methode}

\begin{exoresolu}[Le nylon 6,6 et le tergal]
\begin{enumerate}
\item Le nylon 6,6 est obtenu par polycondensation de l'acide hexanedioïque \ce{HOOC-(CH2)4-COOH} et de l'hexane-1,6-diamine \ce{H2N-(CH2)6-NH2}. Écrire l'équation-bilan de la réaction et encadrer le motif du polymère.
\item Le tergal (Dacron) est obtenu à partir de l'acide benzène-1,4-dicarboxylique (acide téréphtalique) \ce{HOOC-C6H4-COOH} et de l'éthane-1,2-diol \ce{HO-CH2-CH2-OH}. Écrire l'équation-bilan et préciser la nature des liaisons formées.
\item Pourquoi le nom « nylon 6,6 » ?
\end{enumerate}
\tcblower
\textbf{1. Nylon 6,6.} Chaque groupe $-\ce{COOH}$ réagit avec un groupe $-\ce{NH2}$ en éliminant une molécule d'eau (liaison \textbf{amide}) :
\[
n\,\ce{HOOC-(CH2)4-COOH} + n\,\ce{H2N-(CH2)6-NH2}
\;\ce{->}\;
\ce{HO-[\,CO-(CH2)4-CO-NH-(CH2)6-NH\,]_n-H} + 2n\,\ce{H2O}
\]
Le \cle{motif} est : \ce{-CO-(CH2)4-CO-NH-(CH2)6-NH-}. Vérification : à gauche $n(6\mathrm{C}) + n(6\mathrm{C}) = 12n$ carbones ; le motif en contient $4+2+6=12$, répété $n$ fois : bilan correct.

\textbf{2. Tergal.} Chaque groupe $-\ce{COOH}$ réagit avec un groupe $-\ce{OH}$ (liaison \textbf{ester}) :
\[
n\,\ce{HOOC-C6H4-COOH} + n\,\ce{HO-CH2-CH2-OH}
\;\ce{->}\;
\ce{HO-[\,CO-C6H4-CO-O-CH2-CH2-O\,]_n-H} + 2n\,\ce{H2O}
\]
Les liaisons formées sont des liaisons \textbf{ester} : le tergal est un \cle{polyester} (PET), utilisé pour les fibres textiles et les bouteilles plastiques.

\textbf{3. Origine du nom.} « 6,6 » indique le nombre d'atomes de carbone de chaque monomère : \textbf{6} carbones pour la diamine \ce{H2N-(CH2)6-NH2} et \textbf{6} carbones pour le diacide \ce{HOOC-(CH2)4-COOH}.
\textbf{Conclusion :} nylon 6,6 = polyamide (liaisons $-\mathrm{CO-NH}-$) ; tergal = polyester (liaisons $-\mathrm{CO-O}-$) ; dans les deux cas, $2n$ molécules d'eau sont éliminées : c'est bien une \cle{polycondensation}.
\end{exoresolu}

\courssub{Méthode 7 : La saponification — fabrication du savon}

\begin{methode}[Équation et bilan de la saponification]
\begin{enumerate}
\item Identifier le \cle{corps gras} : un \textbf{triester du glycérol} (triglycéride) formé de trois acides gras R-COOH.
\item Écrire l'équation avec la soude (ou la potasse) à \textbf{chaud} :
\[ \text{triglycéride} + 3\,\ce{NaOH} \;\ce{->}\; 3\,\ce{RCOONa} \text{ (savon)} + \text{glycérol} \]
\item Caractériser la réaction : \textbf{totale} et \textbf{rapide à chaud} (contrairement à l'hydrolyse acide, limitée).
\item Pour les calculs : 1 mol de triglycéride $\to$ \textbf{3 mol de savon} et 1 mol de glycérol ; dresser un tableau d'avancement.
\item Décrire le procédé : chauffage du mélange huile + soude concentrée, puis \textbf{relargage} (ajout de sel NaCl concentré : le savon, insoluble dans l'eau salée, précipite), filtration, lavage, séchage.
\end{enumerate}
\textbf{Réflexe Bac :} avec la potasse \ce{KOH} on obtient un savon \emph{mou} (pâtes à savon) ; avec la soude \ce{NaOH}, un savon \emph{dur}.
\end{methode}

\begin{exoresolu}[Savonnerie artisanale à partir d'huile de palme]
Une coopérative de femmes de la région du Littoral fabrique du savon à partir d'huile de palme. On assimile l'huile utilisée à l'oléine, triglycéride de l'acide oléique \ce{C17H33-COOH}, de formule \ce{(C17H33-COO)3-C3H5} et de masse molaire $M_T = \SI{884}{g.mol^{-1}}$. On fait réagir à chaud $m_T = \SI{442}{g}$ d'oléine avec un excès de soude. On donne $M(\ce{C17H33-COONa}) = \SI{304}{g.mol^{-1}}$.
\begin{enumerate}
\item Écrire l'équation-bilan de la saponification et nommer les produits.
\item Calculer la masse de savon obtenue, la réaction étant totale.
\item Citer le produit secondaire valorisable et une de ses utilisations.
\end{enumerate}
\tcblower
\textbf{1. Équation-bilan.}
\[
\ce{(C17H33-COO)3-C3H5} + 3\,\ce{NaOH} \;\ce{->}\; 3\,\ce{C17H33-COONa} + \ce{C3H5(OH)3}
\]
Produits : l'\textbf{oléate de sodium} \ce{C17H33-COONa}, qui est le \cle{savon}, et le \textbf{glycérol} (propane-1,2,3-triol) \ce{C3H5(OH)3}. Vérification du bilan : 3 Na à gauche, 3 Na dans les 3 molécules de savon ; les 3 chaînes \ce{C17H33COO} sont conservées : équation correcte.

\textbf{2. Masse de savon.} Quantité de triglycéride :
\[
n_T = \frac{m_T}{M_T} = \frac{442}{884} = \SI{0,500}{mol}
\]
Tableau d'avancement (la soude est en excès) :
\begin{center}
\begin{tabularx}{\linewidth}{|l|X|X|X|X|}
\hline
\rowcolor{popblueL}
 & oléine & \ce{NaOH} & savon & glycérol \\
\hline
Initial (mol) & $0,500$ & excès & $0$ & $0$ \\
\hline
Final (mol) & $0$ & excès & $3 \times 0,500 = 1,50$ & $0,500$ \\
\hline
\end{tabularx}
\end{center}
\[
m_{\text{savon}} = n_{\text{savon}} \times M_{\text{savon}} = 1,50 \times 304 = \boxed{\SI{456}{g}}
\]

\textbf{3. Produit secondaire.} Le \textbf{glycérol} (glycérine) est récupéré après relargage ; il est utilisé en \textbf{pharmacie} (sirops, suppositoires), en \textbf{cosmétique} (crèmes hydratantes) et dans l'industrie alimentaire.
\textbf{Conclusion :} à partir de \SI{442}{g} d'huile de palme, la coopérative obtient \SI{456}{g} de savon et \SI{46}{g} de glycérol ($0,500 \times 92$), une production entièrement valorisable.
\end{exoresolu}

\begin{attention}[Les 5 erreurs qui coûtent des points au Bac]
\begin{enumerate}
\item \textbf{Oublier le caractère limité des réactions.} L'action de l'eau sur un acide carboxylique et l'estérification directe s'écrivent avec une double flèche \ce{<=>} ; seules les réactions des chlorures d'acyle, des anhydrides et la saponification sont totales (\ce{->}). Inverser les flèches est sanctionné.
\item \textbf{Mal compter les carbones dans la nomenclature.} Le carbone du groupe $-\mathrm{COOH}$ \emph{appartient} à la chaîne principale et porte le numéro 1 : \ce{CH3-CH2-COOH} est l'acide \emph{propan}oïque (3 C), pas « éthanoïque ».
\item \textbf{Confondre estérification et saponification.} L'estérification (acide + alcool) est lente et limitée ; la saponification (ester + soude) est totale et donne un \textbf{carboxylate de sodium} (savon) + un alcool, jamais l'inverse.
\item \textbf{Se tromper de stœchiométrie dans la saponification.} Un triglycéride consomme \textbf{3} moles de soude et produit \textbf{3} moles de savon pour 1 mole de glycérol. Oublier le facteur 3 fausse tout le calcul de masse.
\item \textbf{Écrire un motif de polycondensation faux.} Le motif du nylon 6,6 contient les \emph{deux} monomères : \ce{-CO-(CH2)4-CO-NH-(CH2)6-NH-}. N'oublie pas les $2n$ molécules d'eau éliminées, et vérifie toujours la conservation des atomes entre réactifs et motif.
\end{enumerate}
\end{attention}

\newpage

\coursec{Exercices}

\courssub{Je vérifie mes connaissances}

\begin{exercice}[QCM : les familles de composés oxygénés]
Pour chaque question, une seule réponse est exacte. Entoure la bonne lettre.
\begin{enumerate}
\item La formule générale d'un acide carboxylique saturé à chaîne linéaire est :
\begin{enumerate}
\item[$\square$~a.] $C_nH_{2n}O$ \hfill $\square$~b. $C_nH_{2n}O_2$ \hfill $\square$~c. $C_nH_{2n+2}O_2$
\end{enumerate}
\item Le groupe caractéristique de l'acide éthanoïque est :
\begin{enumerate}
\item[$\square$~a.] le groupe hydroxyle \hfill $\square$~b. le groupe carbonyle \hfill $\square$~c. le groupe carboxyle
\end{enumerate}
\item Le composé de formule \ce{CH3-CO-Cl} est :
\begin{enumerate}
\item[$\square$~a.] un ester \hfill $\square$~b. un chlorure d'acyle \hfill $\square$~c. un anhydride d'acide
\end{enumerate}
\item La réaction de formation d'un savon à partir d'un corps gras et de la soude s'appelle :
\begin{enumerate}
\item[$\square$~a.] l'estérification \hfill $\square$~b. la polycondensation \hfill $\square$~c. la saponification
\end{enumerate}
\item Le nylon 6,6 est :
\begin{enumerate}
\item[$\square$~a.] un polyester \hfill $\square$~b. un polyamide \hfill $\square$~c. un polyalcool
\end{enumerate}
\end{enumerate}
\end{exercice}

\begin{exercice}[Vrai ou faux ? Justifie ta réponse]
Réponds par \textbf{vrai} ou \textbf{faux}, puis justifie en une ou deux phrases.
\begin{enumerate}
\item Les acides carboxyliques sont des acides forts : ils se dissocient totalement dans l'eau.
\item Les acides carboxyliques de petite masse molaire sont solubles dans l'eau grâce aux liaisons hydrogène.
\item La réaction d'estérification entre un acide carboxylique et un alcool est rapide et totale.
\item La réaction d'un chlorure d'acyle avec un alcool donne un ester et du chlorure d'hydrogène.
\item La saponification est une réaction lente mais totale.
\end{enumerate}
\end{exercice}

\begin{exercice}[Texte à trous]
Complète le texte suivant avec les mots ou expressions de la liste ci-dessous (certains mots sont des intrus) :

\smallskip
\textbf{Liste :} carboxyle -- hydrogène -- totale -- limitée -- ester -- amide -- polycondensation -- hydrolyse -- soude -- glycérol -- éthanol.

\smallskip
Les acides carboxyliques possèdent le groupe \ldots\ldots\ldots\ldots, formé d'un groupe carbonyle et d'un groupe hydroxyle portés par le même atome de carbone. Leurs points d'ébullition élevés s'expliquent par l'existence de liaisons \ldots\ldots\ldots\ldots entre les molécules. La réaction d'un acide carboxylique avec un alcool conduit à un \ldots\ldots\ldots\ldots et à de l'eau : c'est une réaction lente et \ldots\ldots\ldots\ldots par un équilibre. Le nylon 6,6 est obtenu par réaction de \ldots\ldots\ldots\ldots entre un diacide et une diamine ; les monomères sont reliés par des liaisons \ldots\ldots\ldots\ldots. La fabrication du savon utilise un corps gras et de la \ldots\ldots\ldots\ldots à chaud.
\end{exercice}

\begin{exercice}[Définitions essentielles]
Définis en une ou deux phrases chacun des termes suivants, en donnant à chaque fois un exemple avec une équation-bilan :
\begin{enumerate}
\item Acide carboxylique ;
\item Chlorure d'acyle ;
\item Anhydride d'acide ;
\item Polycondensation ;
\item Saponification.
\end{enumerate}
\end{exercice}

\courssub{Je m'entraîne}

\begin{exercice}[Nommer et reconnaître les familles]
Pour chacun des composés suivants, donne le nom officiel et précise la famille à laquelle il appartient (acide carboxylique, chlorure d'acyle, anhydride d'acide, ester) :
\begin{enumerate}
\item \ce{CH3-CH2-COOH} ;
\item \ce{CH3-CO-Cl} ;
\item \ce{CH3-CO-O-CO-CH3} ;
\item \ce{CH3-COO-CH2-CH3} ;
\item \ce{HOOC-(CH2)4-COOH} ;
\item \ce{CH3-CH2-CH2-COO-CH3}.
\end{enumerate}
\end{exercice}

\begin{exercice}[Isomères de formule brute \ce{C5H10O2}]
On considère les acides carboxyliques de formule brute \ce{C5H10O2}.
\begin{enumerate}
\item Écris les formules semi-développées de tous ces acides.
\item Nomme chacun d'eux.
\item Parmi eux, lequel possède une chaîne carbonée ramifiée avec deux ramifications ?
\end{enumerate}
\end{exercice}

\begin{exercice}[Liaisons hydrogène et températures d'ébullition]
L'acide propanoïque ($M = \SI{74}{g.mol^{-1}}$, $T_{eb} = \SI{141}{\celsius}$) bout à une température nettement plus élevée que le butan-1-ol ($M = \SI{74}{g.mol^{-1}}$, $T_{eb} = \SI{117}{\celsius}$), alors que les deux composés ont des masses molaires quasi identiques.
\begin{enumerate}
\item Écris les formules semi-développées des deux composés.
\item Explique, en t'appuyant sur les liaisons hydrogène, l'écart observé entre les températures d'ébullition. On précisera en particulier l'existence de \emph{dimères} pour l'acide carboxylique.
\item L'acide propanoïque est-il soluble dans l'eau ? Justifie.
\end{enumerate}
\end{exercice}

\begin{exercice}[Acide éthanoïque en solution aqueuse]
Le vinaigre est une solution aqueuse d'acide éthanoïque \ce{CH3COOH}. On donne : $pK_a(\ce{CH3COOH}/\ce{CH3COO-}) = 4,8$.
\begin{enumerate}
\item Écris l'équation-bilan de la réaction de l'acide éthanoïque avec l'eau. Cette réaction est-elle totale ? Justifie.
\item Écris l'équation-bilan de la réaction de l'acide éthanoïque avec une solution d'hydroxyde de sodium (soude).
\item On prépare une solution d'acide éthanoïque de concentration $C = \SI{0,10}{mol.L^{-1}}$. En admettant que l'acide est peu dissocié, montre que le pH de la solution vaut approximativement :
\[ pH \approx \dfrac{1}{2}\left(pK_a - \log C\right) \]
puis calcule sa valeur.
\end{enumerate}
\end{exercice}

\courssub{Je me prépare au Bac}

\begin{exercice}[Synthèse de l'éthanoate d'éthyle]
\textit{(D'après Baccalauréat TI -- session 2019)}

L'éthanoate d'éthyle est un solvant utilisé dans les colles et les vernis. Au laboratoire, on souhaite le préparer de manière rapide et totale.

\textbf{Données :} masses molaires atomiques (en \si{g.mol^{-1}}) : $M(\ce{C}) = 12$ ; $M(\ce{H}) = 1$ ; $M(\ce{O}) = 16$ ; $M(\ce{Cl}) = 35,5$. Masse molaire de l'éthanoate d'éthyle : $M = \SI{88}{g.mol^{-1}}$.

\begin{enumerate}
\item Écris la formule semi-développée et le nom du chlorure d'éthanoyle.
\item Écris l'équation-bilan de la réaction entre le chlorure d'éthanoyle et l'éthanol. Nomme les produits formés.
\item Cite deux avantages de cette méthode de synthèse par rapport à l'estérification directe de l'acide éthanoïque par l'éthanol.
\item On fait réagir $m = \SI{7,85}{g}$ de chlorure d'éthanoyle avec un excès d'éthanol. La réaction est supposée totale.
\begin{enumerate}
\item Calcule la quantité de matière de chlorure d'éthanoyle utilisée.
\item Construis le tableau d'avancement de la réaction.
\item Déduis-en la masse d'ester obtenue.
\end{enumerate}
\item En réalité, on recueille $\SI{7,5}{g}$ d'ester. Calcule le rendement de la synthèse.
\end{enumerate}
\end{exercice}

\begin{exercice}[Le nylon 6,6, une fibre de polyamide]
\textit{(D'après Baccalauréat TI -- session 2021)}

Le nylon 6,6 est une fibre textile synthétique obtenue par polycondensation de l'acide hexanedioïque (acide adipique) \ce{HOOC-(CH2)4-COOH} et de l'hexane-1,6-diamine \ce{H2N-(CH2)6-NH2}.

\textbf{Données :} masses molaires atomiques (en \si{g.mol^{-1}}) : $M(\ce{C}) = 12$ ; $M(\ce{H}) = 1$ ; $M(\ce{O}) = 16$ ; $M(\ce{N}) = 14$. Masse molaire du motif du nylon 6,6 : $M_{motif} = \SI{226}{g.mol^{-1}}$.

\begin{enumerate}
\item Nomme les deux groupes caractéristiques présents dans les monomères et précise pourquoi ces monomères sont dits \emph{bifonctionnels}.
\item Écris l'équation-bilan de la réaction de polycondensation conduisant au nylon 6,6, en faisant apparaître $n$ motifs. Quel petit composé est éliminé au cours de la réaction ?
\item Encadre le motif du polymère et entoure la liaison caractéristique formée. Comment s'appelle cette liaison ?
\item Vérifie, par le calcul, que la masse molaire du motif vaut bien $\SI{226}{g.mol^{-1}}$.
\item Une usine textile produit un échantillon de nylon 6,6 contenant en moyenne $n = 500$ motifs par macromolécule. Calcule la masse molaire moyenne de cette macromolécule.
\item Le tergal est un autre polymère de polycondensation, obtenu à partir de l'acide benzène-1,4-dicarboxylique (acide téréphtalique) \ce{HOOC-C6H4-COOH} et de l'éthane-1,2-diol \ce{HO-CH2-CH2-OH}.
\begin{enumerate}
\item Écris l'équation-bilan de formation du tergal.
\item Quelle différence de nature de liaison distingue le tergal du nylon 6,6 ? À quelle famille de polymères appartient chacun d'eux ?
\end{enumerate}
\end{enumerate}
\end{exercice}

\begin{exercice}[Fabrication artisanale du savon à partir de l'huile de palme]
\textit{(D'après Baccalauréat TI -- session 2022)}

Dans une savonnerie artisanale de la région du Littoral, on fabrique du savon en chauffant de l'huile de palme avec une solution concentrée d'hydroxyde de sodium. On assimile l'huile de palme à la \textbf{palmitine}, triglycéride de formule :

\[ \chemfig{CH_2(-[2]O-C(=[1]O)-C_{15}H_{31})-CH(-[2]O-C(=[1]O)-C_{15}H_{31})-CH_2(-[6]O-C(=[7]O)-C_{15}H_{31})} \]

\textbf{Données :} masse molaire de la palmitine : $M_P = \SI{806}{g.mol^{-1}}$ ; masse molaire du palmitate de sodium (savon) \ce{C15H31COONa} : $M_S = \SI{278}{g.mol^{-1}}$ ; masse molaire du glycérol : $M_G = \SI{92}{g.mol^{-1}}$.

\begin{enumerate}
\item La palmitine est un triester du glycérol (propane-1,2,3-triol) et de l'acide hexadécanoïque (acide palmitique) \ce{C15H31COOH}. Écris la formule semi-développée du glycérol.
\item Écris l'équation-bilan de la réaction entre la palmitine et l'hydroxyde de sodium en excès, à chaud. Nomme tous les produits obtenus.
\item Comment s'appelle cette réaction ? Cite deux de ses caractéristiques.
\item On traite $m = \SI{100}{g}$ d'huile de palme par un excès de soude.
\begin{enumerate}
\item Calcule la quantité de matière de palmitine mise en jeu.
\item À l'aide d'un tableau d'avancement, détermine la masse de savon et la masse de glycérol obtenues, la réaction étant totale.
\end{enumerate}
\item Pourquoi utilise-t-on la soude en excès et à chaud dans la pratique ?
\item Le savon obtenu est \emph{insoluble} dans l'eau salée. Nomme l'opération qui consiste à ajouter du chlorure de sodium concentré pour faire précipiter le savon, puis cite une utilisation possible du glycérol récupéré comme sous-produit.
\end{enumerate}
\end{exercice}

\newpage

\coursec{Corrigés des exercices}

\begin{corrige}[Exercice 1 — QCM : les familles de composés oxygénés]
\begin{enumerate}
\item \textbf{Réponse b.} Un acide carboxylique saturé à chaîne linéaire a pour formule générale $C_nH_{2n}O_2$ (exemple : \ce{CH3COOH} = \ce{C2H4O2}). La formule $C_nH_{2n}O$ correspond aux aldéhydes et cétones saturés.
\item \textbf{Réponse c.} L'acide éthanoïque \ce{CH3-COOH} possède le groupe \cle{carboxyle} \ce{-COOH}, formé d'un groupe carbonyle \ce{C=O} et d'un groupe hydroxyle \ce{-OH} portés par le même carbone.
\item \textbf{Réponse b.} \ce{CH3-CO-Cl} est le chlorure d'éthanoyle : un \cle{chlorure d'acyle}, de formule générale \ce{R-CO-Cl}.
\item \textbf{Réponse c.} La \cle{saponification} est l'action à chaud d'un corps gras (triester) avec la soude (ou la potasse) : elle donne un savon et du glycérol.
\item \textbf{Réponse b.} Le nylon 6,6 est un \cle{polyamide} : ses motifs sont reliés par des liaisons amide \ce{-CO-NH-}.
\end{enumerate}
\end{corrige}

\begin{corrige}[Exercice 2 — Vrai ou faux ?]
\begin{enumerate}
\item \textbf{Faux.} Les acides carboxyliques sont des \cle{acides faibles} : leur réaction avec l'eau est partielle (équilibre), par exemple \ce{CH3COOH + H2O <=> CH3COO- + H3O+}. Seule une faible fraction des molécules est dissociée.
\item \textbf{Vrai.} Le groupe carboxyle peut former des \cle{liaisons hydrogène} avec les molécules d'eau (donneur et accepteur de liaison H). Les acides à courte chaîne carbonée (moins de 5 atomes de carbone environ) sont donc solubles, voire miscibles, dans l'eau.
\item \textbf{Faux.} L'estérification directe entre un acide carboxylique et un alcool est \textbf{lente} et \textbf{limitée} par un équilibre (elle est aussi athermique). C'est pourquoi on préfère utiliser un chlorure d'acyle ou un anhydride d'acide pour obtenir une réaction rapide et totale.
\item \textbf{Vrai.} $R-CO-Cl + R'-OH \rightarrow R-COO-R' + \ce{HCl}$ : on obtient un ester et du chlorure d'hydrogène. La réaction est rapide et totale.
\item \textbf{Vrai.} La saponification est une réaction \textbf{lente} (d'où le chauffage à reflux) mais \textbf{totale}, car l'ion carboxylate formé est une base trop faible pour réagir en sens inverse avec le glycérol.
\end{enumerate}
\end{corrige}

\begin{corrige}[Exercice 3 — Texte à trous]
Les acides carboxyliques possèdent le groupe \textbf{carboxyle}, formé d'un groupe carbonyle et d'un groupe hydroxyle portés par le même atome de carbone. Leurs points d'ébullition élevés s'expliquent par l'existence de liaisons \textbf{hydrogène} entre les molécules. La réaction d'un acide carboxylique avec un alcool conduit à un \textbf{ester} et à de l'eau : c'est une réaction lente et \textbf{limitée} par un équilibre. Le nylon 6,6 est obtenu par réaction de \textbf{polycondensation} entre un diacide et une diamine ; les monomères sont reliés par des liaisons \textbf{amide}. La fabrication du savon utilise un corps gras et de la \textbf{soude} à chaud.

\smallskip
\emph{Intrus non utilisés :} totale, hydrolyse, glycérol, éthanol.
\end{corrige}

\begin{corrige}[Exercice 4 — Définitions essentielles]
\begin{enumerate}
\item \textbf{Acide carboxylique :} composé organique possédant le groupe caractéristique carboxyle \ce{-COOH}. Formule générale $R\text{-}COOH$ où R est un groupe alkyle ou un atome d'hydrogène. Exemple : l'acide éthanoïque \ce{CH3-COOH}.
\item \textbf{Chlorure d'acyle :} dérivé d'acide carboxylique dans lequel le groupe \ce{-OH} du carboxyle est remplacé par un atome de chlore ; formule générale $R\text{-}CO\text{-}Cl$. Exemple de formation :
\[ \ce{CH3-COOH + SOCl2 -> CH3-CO-Cl + SO2 + HCl} \]
\item \textbf{Anhydride d'acide :} dérivé obtenu par élimination d'une molécule d'eau entre deux groupes carboxyle ; formule générale $R\text{-}CO\text{-}O\text{-}CO\text{-}R'$. Exemple :
\[ \ce{CH3-COOH + CH3-COOH -> CH3-CO-O-CO-CH3 + H2O} \]
\item \textbf{Polycondensation :} réaction de polymérisation au cours de laquelle des monomères bifonctionnels s'associent par étapes avec élimination d'une petite molécule (eau, \ce{HCl}\ldots). Exemple (nylon 6,6) :
\[ n\,\ce{HOOC-(CH2)4-COOH} + n\,\ce{H2N-(CH2)6-NH2} \longrightarrow \text{[-CO-(CH2)4-CO-NH-(CH2)6-NH-]}_n + 2n\,\ce{H2O} \]
\item \textbf{Saponification :} réaction totale d'un corps gras (triester du glycérol) avec la soude (ou la potasse) à chaud, donnant un savon (carboxylate de sodium) et du glycérol. Exemple :
\[ \ce{(C15H31COO)3C3H5 + 3 NaOH -> 3 C15H31COONa + C3H5(OH)3} \]
\end{enumerate}
\end{corrige}

\begin{corrige}[Exercice 5 — Nommer et reconnaître les familles]
\begin{center}
\begin{tabularx}{\linewidth}{|l|X|X|}
\hline
\rowcolor{popblueL} \textbf{Formule} & \textbf{Nom officiel} & \textbf{Famille} \\
\hline
\ce{CH3-CH2-COOH} & acide propanoïque & acide carboxylique \\
\hline
\ce{CH3-CO-Cl} & chlorure d'éthanoyle & chlorure d'acyle \\
\hline
\ce{CH3-CO-O-CO-CH3} & anhydride éthanoïque & anhydride d'acide \\
\hline
\ce{CH3-COO-CH2-CH3} & éthanoate d'éthyle & ester \\
\hline
\ce{HOOC-(CH2)4-COOH} & acide hexanedioïque (acide adipique) & acide carboxylique (diacide) \\
\hline
\ce{CH3-CH2-CH2-COO-CH3} & butanoate de méthyle & ester \\
\hline
\end{tabularx}
\end{center}

\smallskip
\textbf{Méthode de nomenclature :} pour un ester $R-CO-O-R'$, on nomme d'abord la chaîne carbonée liée à l'oxygène (alkyle : méthyle, éthyle\ldots), puis la partie acide dont le suffixe \emph{-oïque} devient \emph{-oate}. Pour un chlorure d'acyle, le suffixe \emph{-oïque} devient \emph{-oyle} précédé de « chlorure de ».
\end{corrige}

\begin{corrige}[Exercice 6 — Isomères de formule brute \ce{C5H10O2}]
\textbf{1. et 2.} Il existe \textbf{quatre} acides carboxyliques de formule brute \ce{C5H10O2} :

\begin{itemize}
\item \chemfig{CH_3-CH_2-CH_2-CH_2-COOH} : \textbf{acide pentanoïque} (chaîne linéaire) ;
\item \chemfig{CH_3-CH_2-CH(-[2]CH_3)-COOH} : \textbf{acide 2-méthylbutanoïque} ;
\item \chemfig{CH_3-CH(-[2]CH_3)-CH_2-COOH} : \textbf{acide 3-méthylbutanoïque} ;
\item \chemfig{CH_3-C(-[2]CH_3)(-[6]CH_3)-COOH} : \textbf{acide 2,2-diméthylpropanoïque}.
\end{itemize}

\textbf{3.} L'acide possédant une chaîne ramifiée avec \textbf{deux ramifications} (deux groupes méthyle sur le même carbone) est l'\textbf{acide 2,2-diméthylpropanoïque}.

\smallskip
\textbf{Point de vigilance :} la chaîne principale doit toujours contenir le carbone du groupe \ce{-COOH}, et ce carbone porte le numéro 1 (il est inutile de le préciser dans le nom).
\end{corrige}

\begin{corrige}[Exercice 7 — Liaisons hydrogène et températures d'ébullition]
\textbf{1.} Formules semi-développées :
\begin{itemize}
\item acide propanoïque : \chemfig{CH_3-CH_2-C(=[1]O)-[7]OH} ;
\item butan-1-ol : \chemfig{CH_3-CH_2-CH_2-CH_2-OH}.
\end{itemize}

\textbf{2.} Les deux molécules possèdent un atome d'hydrogène lié à un oxygène et peuvent donc former des \cle{liaisons hydrogène} intermoléculaires. Mais l'acide propanoïque possède \emph{deux} sites (le \ce{O-H} donneur et le \ce{C=O} accepteur) : deux molécules d'acide s'associent par \textbf{deux liaisons hydrogène} pour former un \emph{dimère} cyclique, une « super-molécule » de masse apparente double ($2 \times 74 = \SI{148}{g.mol^{-1}}$). Le butan-1-ol ne forme qu'une seule liaison hydrogène par paire de molécules. Il faut donc fournir beaucoup plus d'énergie pour séparer les molécules d'acide : d'où $T_{eb} = \SI{141}{\celsius} > \SI{117}{\celsius}$.

\textbf{3.} Oui, l'acide propanoïque est \textbf{soluble dans l'eau} : son groupe carboxyle établit des liaisons hydrogène avec les molécules d'eau, et sa chaîne carbonée (3 atomes de carbone) est encore courte, donc la partie hydrophile l'emporte sur la partie hydrophobe.
\end{corrige}

\begin{corrige}[Exercice 8 — Acide éthanoïque en solution aqueuse]
\textbf{1.} Réaction avec l'eau :
\[ \ce{CH3COOH + H2O <=> CH3COO- + H3O+} \]
Cette réaction \textbf{n'est pas totale} : c'est un équilibre (double flèche), car l'acide éthanoïque est un \cle{acide faible} ($pK_a = 4,8$, valeur finie).

\textbf{2.} Réaction avec la soude :
\[ \ce{CH3COOH + (Na+ + HO-) -> (CH3COO- + Na+) + H2O} \]
soit, en écriture simplifiée : \ce{CH3COOH + HO- -> CH3COO- + H2O}. Cette réaction est quasi totale.

\textbf{3.} L'acide étant peu dissocié, on pose $[\ce{H3O+}] = [\ce{CH3COO-}] = h$ et $[\ce{CH3COOH}] \approx C$. La constante d'acidité s'écrit :
\[ K_a = \dfrac{[\ce{CH3COO-}][\ce{H3O+}]}{[\ce{CH3COOH}]} \approx \dfrac{h^2}{C} \quad\Longrightarrow\quad h^2 = K_a \cdot C \]
En passant aux logarithmes : $-2\log h = -\log K_a - \log C$, c'est-à-dire :
\[ pH \approx \dfrac{1}{2}\left(pK_a - \log C\right) \]
Application numérique : $\log C = \log(0,10) = -1$, donc :
\[ pH \approx \dfrac{1}{2}\left(4,8 - (-1)\right) = \dfrac{5,8}{2} = \boxed{2,9} \]
Vérification de l'hypothèse : $h = 10^{-2,9} \approx \SI{1,3e-3}{mol.L^{-1}} \ll C = \SI{0,10}{mol.L^{-1}}$ : l'acide est bien peu dissocié, l'approximation est valide.
\end{corrige}

\begin{corrige}[Exercice 9 — Synthèse de l'éthanoate d'éthyle]
\textbf{1.} Chlorure d'éthanoyle : \chemfig{CH_3-C(=[1]O)-[7]Cl}, de formule brute \ce{C2H3ClO}.

\textbf{2.} Équation-bilan :
\[ \ce{CH3-CO-Cl + CH3-CH2-OH -> CH3-COO-CH2-CH3 + HCl} \]
Produits : l'\textbf{éthanoate d'éthyle} (ester) et le \textbf{chlorure d'hydrogène}.

\textbf{3.} Avantages par rapport à l'estérification directe : la réaction est \textbf{rapide} et \textbf{totale} (pas d'équilibre limitant), alors que l'estérification acide + alcool est lente et limitée.

\textbf{4.a.} Masse molaire du chlorure d'éthanoyle \ce{C2H3ClO} :
\[ M = 2 \times 12 + 3 \times 1 + 35,5 + 16 = \SI{78,5}{g.mol^{-1}} \]
\[ n = \dfrac{m}{M} = \dfrac{7,85}{78,5} = \SI{0,100}{mol} \]

\textbf{4.b.} Tableau d'avancement (éthanol en excès) :

\begin{center}
\begin{tabularx}{\linewidth}{|l|X|X|X|X|}
\hline
\rowcolor{popblueL} État & \ce{CH3COCl} & \ce{C2H5OH} & \ce{CH3COOC2H5} & \ce{HCl} \\
\hline
Initial ($x=0$) & 0,100 & excès & 0 & 0 \\
\hline
En cours ($x$) & $0,100 - x$ & excès & $x$ & $x$ \\
\hline
Final ($x_{max}$) & 0 & excès & 0,100 & 0,100 \\
\hline
\end{tabularx}
\end{center}

Le chlorure d'acyle est le réactif limitant : $x_{max} = \SI{0,100}{mol}$.

\textbf{4.c.} Masse d'ester théorique :
\[ m_{th} = n \times M = 0,100 \times 88 = \SI{8,8}{g} \]

\textbf{5.} Rendement :
\[ \eta = \dfrac{m_{exp}}{m_{th}} \times 100 = \dfrac{7,5}{8,8} \times 100 \approx \boxed{85,2\ \%} \]

\smallskip
\textbf{Barème indicatif (sur 10) :} formule et nom 1 pt ; équation 2 pts ; avantages 1 pt ; quantité de matière 1,5 pt ; tableau d'avancement 1,5 pt ; masse théorique 1,5 pt ; rendement 1,5 pt.

\textbf{Points de vigilance :} ne pas oublier le chlore dans le calcul de $M(\ce{CH3COCl})$ ; le rendement se calcule toujours par rapport à la masse \emph{théorique} issue du tableau d'avancement ; \ce{HCl} étant gazeux et irritant, la manipulation se fait sous hotte.
\end{corrige}

\begin{corrige}[Exercice 10 — Le nylon 6,6, une fibre de polyamide]
\textbf{1.} L'acide hexanedioïque porte deux groupes \textbf{carboxyle} \ce{-COOH} ; l'hexane-1,6-diamine porte deux groupes \textbf{amine} \ce{-NH2}. Chaque monomère possède \emph{deux} fonctions réactives (une à chaque extrémité) : ils sont dits \textbf{bifonctionnels}, ce qui permet la formation de longues chaînes.

\textbf{2.} Équation-bilan de polycondensation :
\[ \ce{n HOOC-(CH2)4-COOH + n H2N-(CH2)6-NH2 -> [-CO-(CH2)4-CO-NH-(CH2)6-NH-]_n + 2n H2O} \]
Le petit composé éliminé est l'\textbf{eau} (d'où le nom de \emph{polycondensation}).

\textbf{3.} Le motif est $\ce{[-CO-(CH2)4-CO-\underset{\text{liaison amide}}{\underbrace{NH}}-(CH2)6-NH-]}$ ; la liaison caractéristique formée entre un groupe \ce{-CO-} et un groupe \ce{-NH-} est la liaison \textbf{amide} \ce{-CO-NH-} (encadrer le motif entre crochets avec l'indice $n$, entourer le groupe \ce{-CO-NH-}).

\textbf{4.} Le motif a pour formule brute \ce{C12H22N2O2} :
\[ M_{motif} = 12 \times 12 + 22 \times 1 + 2 \times 14 + 2 \times 16 = 144 + 22 + 28 + 32 = \SI{226}{g.mol^{-1}} \]
La valeur donnée est bien vérifiée. \emph{(Astuce de vérification : $M_{motif} = M_{diacide} + M_{diamine} - 2\,M_{eau} = 146 + 116 - 36 = 226$.)}

\textbf{5.} Masse molaire moyenne de la macromolécule :
\[ M = n \times M_{motif} = 500 \times 226 = \boxed{\SI{113000}{g.mol^{-1}}} = \SI{113}{kg.mol^{-1}} \]

\textbf{6.a.} Formation du tergal :
\[ \ce{n HOOC-C6H4-COOH + n HO-CH2-CH2-OH -> [-CO-C6H4-CO-O-CH2-CH2-O-]_n + 2n H2O} \]

\textbf{6.b.} Dans le tergal, les motifs sont reliés par des liaisons \textbf{ester} \ce{-CO-O-} : c'est un \textbf{polyester}. Dans le nylon 6,6, les motifs sont reliés par des liaisons \textbf{amide} \ce{-CO-NH-} : c'est un \textbf{polyamide}.

\smallskip
\textbf{Barème indicatif (sur 10) :} groupes et bifonctionnalité 1,5 pt ; équation nylon 2 pts ; motif et liaison amide 1,5 pt ; vérification de $M_{motif}$ 1,5 pt ; masse molaire moyenne 1 pt ; équation tergal 1,5 pt ; comparaison des familles 1 pt.

\textbf{Points de vigilance :} bien écrire $2n$ molécules d'eau (deux fonctions par monomère) ; ne pas confondre le nombre d'atomes de carbone des monomères (6 et 6, d'où « nylon 6,6 ») avec celui du motif (12) ; les coefficients $n$ doivent apparaître devant chaque monomère.
\end{corrige}

\begin{corrige}[Exercice 11 — Fabrication artisanale du savon à partir de l'huile de palme]
\textbf{1.} Glycérol (propane-1,2,3-triol) : \chemfig{CH_2(-[2]OH)-CH(-[2]OH)-CH_2(-[6]OH)}, soit \ce{HO-CH2-CH(OH)-CH2-OH}.

\textbf{2.} Équation-bilan de la réaction :
\[ \ce{(C15H31COO)3C3H5 + 3 NaOH -> 3 C15H31COONa + C3H5(OH)3} \]
Produits : le \textbf{palmitate de sodium} \ce{C15H31COONa}, qui est le \textbf{savon}, et le \textbf{glycérol} (propane-1,2,3-triol).

\textbf{3.} Cette réaction s'appelle la \textbf{saponification}. Deux caractéristiques : elle est \textbf{lente} (d'où le chauffage) et \textbf{totale}.

\textbf{4.a.} Quantité de matière de palmitine :
\[ n_P = \dfrac{m}{M_P} = \dfrac{100}{806} \approx \SI{0,124}{mol} \]

\textbf{4.b.} Tableau d'avancement (soude en excès) :

\begin{center}
\begin{tabularx}{\linewidth}{|l|X|X|X|X|}
\hline
\rowcolor{popblueL} État & palmitine & \ce{NaOH} & savon \ce{C15H31COONa} & glycérol \\
\hline
Initial & 0,124 & excès & 0 & 0 \\
\hline
En cours & $0,124 - x$ & excès & $3x$ & $x$ \\
\hline
Final & 0 & excès & $3 \times 0,124 = 0,372$ & 0,124 \\
\hline
\end{tabularx}
\end{center}

La palmitine est le réactif limitant : $x_{max} = \SI{0,124}{mol}$.

Masse de savon obtenue :
\[ m_S = 3\,x_{max} \times M_S = 0,372 \times 278 \approx \boxed{\SI{103,4}{g}} \]

Masse de glycérol obtenue :
\[ m_G = x_{max} \times M_G = 0,124 \times 92 \approx \boxed{\SI{11,4}{g}} \]

\textbf{5.} La soude est utilisée \textbf{en excès} pour déplacer la réaction dans le sens direct et garantir que tout le corps gras réagit (réaction totale, aucun triglycéride résiduel) ; le \textbf{chauffage} augmente la vitesse de cette réaction lente (facteur cinétique : la température).

\textbf{6.} L'opération qui consiste à ajouter une solution concentrée de chlorure de sodium pour faire précipiter le savon s'appelle le \textbf{relargage}. Le glycérol récupéré peut être utilisé dans la fabrication de \textbf{produits pharmaceutiques ou cosmétiques} (crèmes, lotions), de savons liquides, ou encore comme matière première dans l'industrie alimentaire.

\smallskip
\textbf{Barème indicatif (sur 10) :} formule du glycérol 1 pt ; équation-bilan 2 pts ; nom et caractéristiques 1,5 pt ; quantité de matière 1 pt ; tableau d'avancement 1,5 pt ; masses de savon et de glycérol 1,5 pt ; rôle de l'excès et du chauffage 1 pt ; relargage et valorisation du glycérol 0,5 pt.

\textbf{Points de vigilance :} attention au coefficient stœchiométrique \textbf{3} devant le savon (une mole de triglycéride donne trois moles de savon) : c'est l'erreur la plus fréquente au Bac ; ne pas confondre saponification (totale) et estérification/hydrolyse (équilibrées) ; vérifier que la somme des masses des produits ($103,4 + 11,4 + \text{excès de soude}$) reste cohérente avec la conservation de la matière.
\end{corrige}

\newpage

\coursec{Fiche bilan}

\coursec{Bilan synthétique : Les acides carboxyliques et leurs dérivés}

\begin{popfigure}
\begin{center}\tcbox[colback=popink,colframe=popink,coltext=white,arc=9pt,boxsep=4pt,left=6pt,right=6pt]{\bfseries\small Acides carboxyliques et dérivés}\end{center}
\vspace{-2pt}
\begin{tcbraster}[raster columns=3,raster equal height=rows,raster column skip=6pt,raster row skip=6pt,enhanced,blankest]
\begin{tcolorbox}[enhanced,colback=white,colframe=poporange,boxrule=0.9pt,arc=3pt,title={Structure $R-\mathrm{COOH}$},fonttitle=\bfseries\scriptsize,coltitle=white,colbacktitle=poporange,left=4pt,right=4pt,top=2pt,bottom=2pt,boxsep=1pt,toptitle=1.5pt,bottomtitle=1.5pt,halign=left]
\begin{itemize}[nosep,leftmargin=*,itemsep=1pt,font=\scriptsize]
\item Groupe carboxyle $-\mathrm{C(=O)OH}$
\item Nomenclature acide \dots oïque
\item Mono/polyacides saturés, insaturés
\end{itemize}
\end{tcolorbox}
\begin{tcolorbox}[enhanced,colback=white,colframe=popgreen,boxrule=0.9pt,arc=3pt,title={Propriétés physiques},fonttitle=\bfseries\scriptsize,coltitle=white,colbacktitle=popgreen,left=4pt,right=4pt,top=2pt,bottom=2pt,boxsep=1pt,toptitle=1.5pt,bottomtitle=1.5pt,halign=left]
\begin{itemize}[nosep,leftmargin=*,itemsep=1pt,font=\scriptsize]
\item Dimères par liaisons H
\item $T_{\text{éb}}$ élevées
\item Solubilité $C \le 4$
\end{itemize}
\end{tcolorbox}
\begin{tcolorbox}[enhanced,colback=white,colframe=poppurple,boxrule=0.9pt,arc=3pt,title={Propriété acide},fonttitle=\bfseries\scriptsize,coltitle=white,colbacktitle=poppurple,left=4pt,right=4pt,top=2pt,bottom=2pt,boxsep=1pt,toptitle=1.5pt,bottomtitle=1.5pt,halign=left]
\begin{itemize}[nosep,leftmargin=*,itemsep=1pt,font=\scriptsize]
\item $\ce{RCOOH + H2O <=> RCOO- + H3O+}$
\item Acide faible $pK_A \approx 4,8$
\item Dosage par soude
\end{itemize}
\end{tcolorbox}
\begin{tcolorbox}[enhanced,colback=white,colframe=poppink,boxrule=0.9pt,arc=3pt,title={Dérivés activés},fonttitle=\bfseries\scriptsize,coltitle=white,colbacktitle=poppink,left=4pt,right=4pt,top=2pt,bottom=2pt,boxsep=1pt,toptitle=1.5pt,bottomtitle=1.5pt,halign=left]
\begin{itemize}[nosep,leftmargin=*,itemsep=1pt,font=\scriptsize]
\item Chlorure d'acyle $R\mathrm{COCl}$
\item Anhydride $(R\mathrm{CO})_2\mathrm{O}$
\item Réactifs acylation
\end{itemize}
\end{tcolorbox}
\begin{tcolorbox}[enhanced,colback=white,colframe=popgold,boxrule=0.9pt,arc=3pt,title={Estérification},fonttitle=\bfseries\scriptsize,coltitle=white,colbacktitle=popgold,left=4pt,right=4pt,top=2pt,bottom=2pt,boxsep=1pt,toptitle=1.5pt,bottomtitle=1.5pt,halign=left]
\begin{itemize}[nosep,leftmargin=*,itemsep=1pt,font=\scriptsize]
\item Directe : lente, limitée, athermique
\item Via $\mathrm{COCl}$ : rapide, totale
\item Via anhydride : totale
\end{itemize}
\end{tcolorbox}
\begin{tcolorbox}[enhanced,colback=white,colframe=popblue!70,boxrule=0.9pt,arc=3pt,title={Polycondensation},fonttitle=\bfseries\scriptsize,coltitle=white,colbacktitle=popblue!70,left=4pt,right=4pt,top=2pt,bottom=2pt,boxsep=1pt,toptitle=1.5pt,bottomtitle=1.5pt,halign=left]
\begin{itemize}[nosep,leftmargin=*,itemsep=1pt,font=\scriptsize]
\item Nylon 6,6 : diamine + diacide
\item Tergal : diacide aromatique + diol
\item Élimination $\ce{H2O}$
\end{itemize}
\end{tcolorbox}
\begin{tcolorbox}[enhanced,colback=white,colframe=popgreen!70!popdark,boxrule=0.9pt,arc=3pt,title={Saponification},fonttitle=\bfseries\scriptsize,coltitle=white,colbacktitle=popgreen!70!popdark,left=4pt,right=4pt,top=2pt,bottom=2pt,boxsep=1pt,toptitle=1.5pt,bottomtitle=1.5pt,halign=left]
\begin{itemize}[nosep,leftmargin=*,itemsep=1pt,font=\scriptsize]
\item Triglycéride $+$ $\ce{NaOH/KOH}$
\item Savon $+$ Glycérol
\item Totale, à chaud
\end{itemize}
\end{tcolorbox}
\end{tcbraster}

\legende{Carte mentale de la leçon : les sept grandes idées à maîtriser pour le Bac TI.}
\end{popfigure}

\begin{aretenir}[L'essentiel de la leçon]

\textbf{1. Définitions et structures}
\begin{itemize}
  \item \cle{Acide carboxylique} : composé organique de formule générale \textbf{$R-\mathrm{COOH}$} (ou $R-\mathrm{C(=O)OH}$). Le \cle{groupe carboxyle} $-\mathrm{COOH}$ résulte de la réunion sur un même carbone d'un groupe \textbf{carbonyle} $\mathrm{C{=}O}$ et d'un groupe \textbf{hydroxyle} $-\mathrm{OH}$.
  \item \cle{Monoacides saturés} : $C_nH_{2n+1}\mathrm{COOH}$ (ex. acide éthanoïque $\ce{CH3COOH}$). \cle{Insaturés} : présence de doubles liaisons dans $R$ (ex. acide oléique $\ce{C17H33COOH}$).
  \item \cle{Polyacides} : plusieurs groupes $\mathrm{COOH}$ (ex. acide hexanedioïque $\ce{HOOC-(CH2)4-COOH}$, acide benzène-1,4-dicarboxylique).
  \item \cle{Chlorure d'acyle} $R-\mathrm{COCl}$ et \cle{anhydride d'acide} $(R\mathrm{CO})_2\mathrm{O}$ : dérivés où le $-\mathrm{OH}$ du carboxyle est remplacé par $-\mathrm{Cl}$ ou $-\mathrm{OCOR}$.
  \item \cle{Ester} $R-\mathrm{COO}-R'$ : fonction issue de la condensation acide/alcool (ou dérivé/alcool).
  \item \cle{Savon} : sel de sodium ou potassium d'un acide gras (carboxylate $R\mathrm{COO}^- \ce{Na^+}$), obtenu par \cle{saponification} (hydrolyse basique d'un triglycéride).
  \item \cle{Polycondensation} : polymérisation par étapes entre monomères bifonctionnels (ou plus) avec élimination d'une petite molécule ($\ce{H2O}$, $\ce{HCl}$). Donne polyamides (nylon 6,6) ou polyesters (tergal).
\end{itemize}

\textbf{2. Nomenclature (IUPAC français) — Règles clés}
\begin{itemize}
  \item \textbf{Acide} : chaîne la plus longue portant le $\mathrm{COOH}$ $\rightarrow$ suffixe \textbf{-oïque}, préfixe \textbf{acide}. Le carbone du $\mathrm{COOH}$ est \textbf{toujours C1}. Ex. $\ce{CH3CH2COOH}$ : acide propanoïque.
  \item \textbf{Chlorure d'acyle} : suffixe \textbf{-oyle} + \textbf{chlorure}. Ex. $\ce{CH3COCl}$ : chlorure d'éthanoyle.
  \item \textbf{Anhydride} : nom de l'acide + \textbf{anhydride}. Ex. $(\ce{CH3CO})_2\mathrm{O}$ : anhydride éthanoïque.
  \item \textbf{Ester} : alkyl (alcool) + acide (suffixe \textbf{-oate}). Ex. $\ce{CH3COOCH2CH3}$ : éthanoate d'éthyle.
  \item \textbf{Savon} : cation (sodium/potassium) + anion (nom de l'acide en -oate). Ex. $\ce{C17H35COONa}$ : stéarate de sodium.
\end{itemize}

\textbf{3. Propriétés physiques — Rôle des liaisons hydrogène}
\begin{propriete}[Liaisons hydrogène et dimères]
Les acides carboxyliques forment des \textbf{dimères cycliques} stables par deux liaisons hydrogène intermoléculaires $\mathrm{O-H\cdots O=C}$ (énergie $\approx \SI{60}{kJ.mol^{-1}}$ par liaison).
\begin{itemize}
  \item \textbf{Températures d'ébullition} : anormalement élevées (proches de celles des alcools de masse molaire double) car il faut rompre les dimères pour vaporiser.
  \item \textbf{Solubilité dans l'eau} : excellente pour les acides de chaîne courte ($n_C \le 4$) car le groupe $\mathrm{COOH}$ s'hydrate fortement (donneur et accepteur de liaisons H). Chute brutale au-delà ($n_C > 4$) du fait de la chaîne hydrocarbonée hydrophobe.
\end{itemize}
\end{propriete}

\textbf{4. Propriété acide — Équilibre dans l'eau}
\begin{definition}[Acidité des acides carboxyliques]
Dans l'eau, l'acide carboxylique est un \textbf{acide faible} : il cède un proton à l'eau de façon limitée.
\[ \ce{RCOOH + H2O <=> RCOO- + H3O+} \qquad K_A = \frac{[\ce{RCOO-}][\ce{H3O+}]}{[\ce{RCOOH}]} \]
$pK_A \approx 4,8$ pour l'acide éthanoïque. Le couple acide/base est \ce{RCOOH}/\ce{RCOO-}.
\end{definition}
\begin{attention}[Piège : force acide]
Ne pas confondre \emph{acide faible} (équilibre déplacé vers les réactifs) et \emph{acide peu concentré}. Un acide carboxylique concentré reste un acide faible ($pK_A$ constant). La basicité de l'anion carboxylate explique l'hydrolyse basique des esters (saponification).
\end{attention}

\textbf{5. Synthèse des dérivés activés (chlorures, anhydrides)}
\begin{align*}
\textbf{Chlorure d'acyle (réactif de choix : chlorure de thionyle $\ce{SOCl2}$)} &: \\
\ce{RCOOH + SOCl2 &-> RCOCl + SO2 ^ + HCl ^} \quad \text{(produits gazeux : réaction totale, purification aisée)} \\[0.5em]
\textbf{Anhydride d'acide (déshydratation par $\ce{P2O5}$ ou chauffage fort)} &: \\
\ce{2 RCOOH &->[P2O5,\ \Delta] (RCO)2O + H2O} \quad \text{(symétrique ; les anhydrides mixtes ne s'obtiennent pas ainsi)}
\end{align*}

\textbf{6. Estérification : directe vs via dérivés activés}
\begin{center}
\begin{tabularx}{\linewidth}{|l|X|X|X|}
\hline
\rowcolor{popblueL} \textbf{Voie} & \textbf{Équation-bilan} & \textbf{Caractères} & \textbf{Conditions / Remarques} \\
\hline
Directe (Fischer) & $RCOOH + R'OH \rightleftharpoons RCOOR' + H_2O$ & Lente, limitée, athermique ($K \approx 1\text{--}10$) & Chauffage à reflux, catalyse acide forte ($\ce{H2SO4}$), excès d'un réactif ou élimination $\ce{H2O}$ pour déplacer l'équilibre. \\
\hline
Chlorure d'acyle & $RCOCl + R'OH \rightarrow RCOOR' + HCl$ & \textbf{Rapide, totale} (à T° ambiante) & \textbf{Présence de pyridine} (base) pour piéger $\ce{HCl}$ et éviter l'estérification du chlorure. \\
\hline
Anhydride d'acide & $(RCO)_2O + R'OH \rightarrow RCOOR' + RCOOH$ & \textbf{Totale} & Moins réactif que le chlorure, mais plus sélectif. Donne 1 ester + 1 acide. \\
\hline
\end{tabularx}
\end{center}

\textbf{7. Polycondensation : Nylon 6,6 et Tergal}
\begin{exemplebox}[Nylon 6,6 (Polyamide 6,6)]
Monomères : \textbf{hexane-1,6-diamine} $\ce{H2N-(CH2)6-NH2}$ + \textbf{acide hexanedioïque (adipique)} $\ce{HOOC-(CH2)4-COOH}$.
\[ n\,\ce{H2N-(CH2)6-NH2} + n\,\ce{HOOC-(CH2)4-COOH} \xrightarrow{\Delta, \text{pression réduite}} \underbrace{\ce{[-NH-(CH2)6-NH-CO-(CH2)4-CO-]}}_n + 2n\,\ce{H2O} \]
Liaison amide $\ce{-CO-NH-}$. Élimination de $2n$ molécules d'eau.
\end{exemplebox}

\begin{exemplebox}[Tergal / Dacron (Polyester : polyéthylène téréphtalate PET)]
Monomères : \textbf{acide benzène-1,4-dicarboxylique (téréphtalique)} $\ce{HOOC-C6H4-COOH}$ + \textbf{éthane-1,2-diol (éthylène glycol)} $\ce{HO-CH2-CH2-OH}$.
\[ n\,\ce{HOOC-C6H4-COOH} + n\,\ce{HO-CH2-CH2-OH} \xrightarrow{\Delta, \text{catalyseur}} \underbrace{\ce{[-OOC-C6H4-COO-CH2-CH2-O-]}}_n + 2n\,\ce{H2O} \]
Liaison ester $\ce{-COO-}$. Élimination de $2n$ molécules d'eau.
\end{exemplebox}

\textbf{8. Saponification — Fabrication du savon (contexte camerounais : huile de palme)}
\begin{definition}[Saponification]
Hydrolyse \textbf{basique} (soude $\ce{NaOH}$ ou potasse $\ce{KOH}$) d'un \textbf{triglycéride} (triester du glycérol et d'acides gras $R\mathrm{COOH}$). Réaction \textbf{lente à froid, rapide à chaud, totale} (irréversible car l'acide formé est neutralisé en carboxylate).
\[
\mathrm{CH_2(OOC}R)-\mathrm{CH(OOC}R')-\mathrm{CH_2(OOC}R'') + 3\,\mathrm{NaOH} \rightarrow
R\mathrm{COONa} + R'\mathrm{COONa} + R''\mathrm{COONa} + \mathrm{CH_2OH-CHOH-CH_2OH}
\]
\textbf{Produits} : \textbf{Savon} (mélange de carboxylates de sodium/potassium) + \textbf{Glycérol} (propane-1,2,3-triol).
\end{definition}

\begin{methode}[Fabrication artisanale du savon (procédé « chaud » type savonneries locales)]
\begin{enumerate}
  \item \textbf{Saponification} : Chauffer le corps gras (huile de palme rouge, huile de palmiste, graisse de bœuf) avec une solution concentrée de soude ($\ce{NaOH}$) ou potasse ($\ce{KOH}$) sous agitation jusqu'à « trace » (épaississement).
  \item \textbf{Relargage (salage)} : Ajouter du chlorure de sodium solide (sel de table) $\rightarrow$ le savon, peu soluble dans l'eau salée, se sépare en phase supérieure (pâte savonneuse) ; la lessive (eau + glycérol + excès de soude + sel) reste en bas.
  \item \textbf{Lavage} : Retirer la pâte, la rincer à l'eau chaude pour éliminer le sel, la soude résiduelle et les impuretés.
  \item \textbf{Moulage / Séchage} : Couler dans des moules, laisser durcir et sécher plusieurs semaines (cure) pour évacuer l'eau résiduelle et améliorer la dureté.
\end{enumerate}
\end{methode}

\textbf{9. Tableau récapitulatif des formules et noms}
\begin{center}
\begin{tabularx}{\linewidth}{|l|X|X|}
\hline
\rowcolor{popblueL} \textbf{Famille} & \textbf{Formule générale / Structure} & \textbf{Exemple (nom IUPAC)} \\
\hline
Acide carboxylique & $\mathrm{R}-\mathrm{COOH}$ & \ce{CH3COOH} : acide éthanoïque \\
\hline
Chlorure d'acyle & $\mathrm{R}-\mathrm{COCl}$ & \ce{CH3COCl} : chlorure d'éthanoyle \\
\hline
Anhydride d'acide (sym.) & $(\mathrm{RCO})_2\mathrm{O}$ & $(\ce{CH3CO})_2\mathrm{O}$ : anhydride éthanoïque \\
\hline
Ester & $\mathrm{R}-\mathrm{COO}-\mathrm{R}'$ & \ce{CH3COOCH2CH3} : éthanoate d'éthyle \\
\hline
Savon (carboxylate) & $\mathrm{R}-\mathrm{COO}^- \,\ce{Na^+}$ & $\ce{C17H35COONa}$ : stéarate de sodium \\
\hline
Triglycéride (corps gras) & $\mathrm{CH_2(OOCR)-CH(OOCR')-CH_2(OOCR'')}$ & Tristéarine (sat.) / Trioléine (insat.) \\
\hline
\end{tabularx}
\end{center}

\end{aretenir}

\begin{savaistu}[Le savon au Cameroun : de l'huile de palme au « savon noir »]
L'huile de palme (issue du fruit du palmier à huile \emph{Elaeis guineensis}, culture majeure au Cameroun) est riche en acides palmitique ($C_{15}H_{31}\mathrm{COOH}$, saturé) et oléique ($C_{17}H_{33}\mathrm{COOH}$, insaturé). La saponification traditionnelle à la potasse ($\ce{KOH}$, lessive de cendres) donne un \textbf{savon mou/noir} (potasse $\rightarrow$ savon mou ; soude $\rightarrow$ savon dur). Le glycérol récupéré est un sous-produit valorisé en cosmétique ou en nitroglycérine (explosifs). Les savonneries industrielles (ex. \textbf{Savonnerie du Cameroun}, unités artisanales de \textbf{Bafoussam}, \textbf{Foumban}) utilisent ce procédé. Attention : le « savon de Marseille » authentique contient 72\% d'acides gras (huile d'olive) ; nos savons locaux sont souvent à base d'huile de palme/palmiste (plus durs, moussent moins).
\end{savaistu}

\begin{attention}[Erreurs fréquentes au Bac — Vigilance !]
\begin{itemize}
  \item \textbf{Nomenclature} : Oublier que le carbone du $\mathrm{COOH}$ est C1 (ex. $\ce{CH3CH2COOH}$ = acide \textbf{propanoïque}, pas butanoïque). Confondre « -oïque » (acide) et « -oate » (ester/savon).
  \item \textbf{Équations} : Écrire $R\mathrm{COOH} + R'\mathrm{OH} \rightleftharpoons R\mathrm{COO}R' + \ce{H2O}$ avec une flèche simple $\rightarrow$ pour l'estérification directe (\textbf{interdit} : c'est un équilibre $\ce{<=>}$). Oublier l'eau dans la polycondensation (nylon 6,6 : $2n\,\ce{H2O}$, pas $n$).
  \item \textbf{Saponification} : Croire qu'on part du glycérol + acide gras (c'est l'estérification inverse !). Le substrat est le \textbf{triglycéride} (huile/gras). Écrire $\ce{NaOH}$ au lieu de $3\,\ce{NaOH}$ pour un triglycéride.
  \item \textbf{Chlorure d'acyle} : Oublier la pyridine (ou base) pour neutraliser $\ce{HCl}$ ; sans elle, l'ester peut être hydrolysé par l'acide chlorhydrique formé.
  \item \textbf{Propriétés physiques} : Dire « liaisons hydrogène \emph{intramoléculaires} » au lieu d'\emph{intermoléculaires} (dimères). Dire « solubilité totale dans l'eau » sans préciser la limite $C \le 4$.
  \item \textbf{Polycondensation} : Confondre avec polymérisation par addition (pas d'élimination de petite molécule). Oublier la bifonctionnalité des monomères.
\end{itemize}
\end{attention}

\begin{exemplebox}[Nomenclature : cas particuliers (insaturé, diacide, ramifié)]
\begin{itemize}
  \item $\ce{CH2=CH-COOH}$ : acide \textbf{prop-2-énoïque} (acide acrylique). Numérotation : C1=COOH, C2=CH, C3=CH2. Double liaison en 2.
  \item $\ce{HOOC-CH2-CH2-COOH}$ : acide \textbf{butanedioïque} (acide succinique).
  \item $\ce{CH3-CH(CH3)-CH2-COOH}$ : acide \textbf{3-méthylbutanoïque} (chaîne 4C + méthyle en 3).
  \item $\ce{C6H5-COOH}$ : acide \textbf{benzoïque} (nom retenu IUPAC).
\end{itemize}
\end{exemplebox}

\begin{methode}[Écrire l'équation-bilan d'une polycondensation (Nylon 6,6 ou Tergal)]
\begin{enumerate}
  \item Identifier les \textbf{deux monomères} et leurs fonctions (amine/acide ou alcool/acide).
  \item Écrire leur formule semi-développée en faisant apparaître les groupes réactifs ($-\mathrm{NH2}$, $-\mathrm{COOH}$, $-\mathrm{OH}$).
  \item Assembler : lier chaque fonction acide à une fonction amine/alcool en formant la liaison amide/ester et en libérant $\ce{H2O}$.
  \item Vérifier la \textbf{conservation des atomes} : atomes C, H, O, N identiques des deux côtés.
  \item Écrire le coefficient stœchiométrique $n$ devant les monomères et la petite molécule ($2n\,\ce{H2O}$ pour diacide+diamine/diol).
  \item Représenter l'unité répétitive entre crochets $\ce{[-...-]}_n$.
\end{enumerate}
\end{methode}

\begin{aretenir}[Checklist finale — Suis-je prêt pour l'épreuve ?]
\begin{itemize}
  \item $\square$ J'écris la formule générale $R-\mathrm{COOH}$ et je représente le groupe carboxyle (carbonyle + hydroxyle sur même C).
  \item $\square$ Je nomme correctement : acide \dots oïque, chlorure d'\dots oyle, anhydride \dots oïque, \dots oate d'alkyle, \dots oate de sodium.
  \item $\square$ J'écris les formules semi-développées de monoacides saturés ($C_nH_{2n+1}\mathrm{COOH}$), insaturés (double liaison dans $R$) et diacides ($HOOC-R-COOH$).
  \item $\square$ J'explique $T_{\text{éb}}$ élevées et solubilité ($C \le 4$) par la formation de \textbf{dimères par liaisons H intermoléculaires}.
  \item $\square$ J'écris l'équation d'acidité $RCOOH + H_2O \rightleftharpoons RCOO^- + H_3O^+$ et je cite $pK_A \approx 4,8$ (acide faible).
  \item $\square$ Je connais les 3 caractères de l'estérification directe : \textbf{lente, limitée, athermique} ($\rightleftharpoons$).
  \item $\square$ J'écris les synthèses d'esters \textbf{totales} : via chlorure d'acyle ($RCOCl + R'OH \rightarrow$, + pyridine) et via anhydride ($(RCO)_2O + R'OH \rightarrow$).
  \item $\square$ J'écris les équations de polycondensation : \textbf{Nylon 6,6} (hexane-1,6-diamine + acide hexanedioïque $\rightarrow$ polyamide + $2n\,\ce{H2O}$) et \textbf{Tergal} (acide téréphtalique + éthane-1,2-diol $\rightarrow$ polyester + $2n\,\ce{H2O}$).
  \item $\square$ Je définis la \textbf{saponification} : hydrolyse basique d'un \textbf{triglycéride} par $3\,\ce{NaOH}$ (ou $\ce{KOH}$) $\rightarrow$ 3 savons + glycérol. Réaction \textbf{totale}.
  \item $\square$ Je décris les 4 étapes de la fabrication du savon : saponification à chaud $\rightarrow$ relargage au sel $\rightarrow$ lavage $\rightarrow$ moulage/séchage.
  \item $\square$ Je distingue clairement \textbf{estérification} (équilibre, acide+alcool) et \textbf{saponification} (totale, triglycéride+base).
  \item $\square$ Je vérifie \textbf{systématiquement} la conservation des atomes (C, H, O, N, Cl) et des charges dans toutes mes équations-bilans.
\end{itemize}
\end{aretenir}

\findecours

\end{document}
